\documentclass[11pt,oneside]{book}
\usepackage[utf8]{inputenc}
\usepackage[T1]{fontenc}
\usepackage[french]{babel}
\usepackage{lmodern}
\usepackage[a4paper,margin=2.45cm,headheight=28pt]{geometry}
\usepackage{amsmath,amssymb,amsthm,mathtools,bm,stmaryrd}
\usepackage{booktabs,longtable,tabularx,array,enumitem}
\usepackage{xcolor,microtype,url,hyperref,listings,fancyhdr}
\hypersetup{colorlinks=true,linkcolor=blue!50!black,citecolor=blue!50!black,urlcolor=blue!50!black,pdftitle={NGeGenBrain : formalisme unifié de l'observation et de la reconstruction},pdfsubject={Empreintes informationnelles, invariants, variations, observation multi-échelle et reconstruction mécaniste}}
\setlist{itemsep=2pt,topsep=4pt}
\setlength{\emergencystretch}{3em}
\setlength{\headheight}{28pt}
\setlength{\tabcolsep}{4pt}
\setcounter{tocdepth}{2}
\setcounter{secnumdepth}{3}
\pagestyle{fancy}\fancyhf{}
\fancyhead[L]{\small NGeGenBrain -- théorie unifiée}
\fancyhead[R]{\small Observation, empreintes et mécanismes}
\fancyfoot[C]{\thepage}
\newtheorem{theoreme}{Théorème}[chapter]
\newtheorem{proposition}[theoreme]{Proposition}
\newtheorem{lemme}[theoreme]{Lemme}
\newtheorem{corollaire}[theoreme]{Corollaire}
\newtheorem{hypothese}[theoreme]{Hypothèse explicite}
\theoremstyle{definition}
\newtheorem{definition}[theoreme]{Définition}
\newtheorem{convention}[theoreme]{Convention}
\newtheorem{critere}[theoreme]{Critère}
\newtheorem{algorithme}[theoreme]{Algorithme}
\theoremstyle{remark}
\newtheorem{remarque}[theoreme]{Remarque}
\newtheorem{contreexemple}[theoreme]{Contre-exemple}
\newcommand{\X}{\mathcal X}
\newcommand{\Y}{\mathcal Y}
\newcommand{\Z}{\mathcal Z}
\newcommand{\Pcal}{\mathcal P}
\newcommand{\Fcal}{\mathcal F}
\newcommand{\Gcal}{\mathcal G}
\newcommand{\Lcal}{\mathcal L}
\newcommand{\Zcal}{\mathcal Z}
\newcommand{\Ycal}{\mathcal Y}
\newcommand{\Tcal}{\mathcal T}
\newcommand{\Ecal}{\mathcal E}
\newcommand{\Hcal}{\mathcal H}
\newcommand{\Rcal}{\mathcal R}
\newcommand{\Ical}{\mathcal I}
\newcommand{\Bcal}{\mathcal B}
\newcommand{\Mcal}{\mathcal M}
\newcommand{\Acal}{\mathcal A}
\newcommand{\Ucal}{\mathcal U}
\newcommand{\Vcal}{\mathcal V}
\newcommand{\Scal}{\mathcal S}
\newcommand{\Dcal}{\mathcal D}
\newcommand{\U}{\mathcal U}
\newcommand{\W}{\mathcal W}
\newcommand{\N}{\mathbb N}
\newcommand{\R}{\mathbb R}
\newcommand{\Prob}{\mathbb P}
\newcommand{\Exp}{\mathbb E}
\newcommand{\1}{\mathbf 1}
\newcommand{\dd}{\mathrm d}
\newcommand{\cl}{\operatorname{cl}}
\newcommand{\Id}{\operatorname{Id}}
\newcommand{\TV}{\operatorname{TV}}
\newcommand{\KL}{\operatorname{KL}}
\newcommand{\argmin}{\operatorname*{arg\,min}}
\newcommand{\argmax}{\operatorname*{arg\,max}}
\newcommand{\supp}{\operatorname{supp}}
\newcommand{\suppo}{\operatorname{supp}}
\newcommand{\diam}{\operatorname{diam}}
\newcommand{\rank}{\operatorname{rank}}
\newcommand{\Var}{\operatorname{Var}}
\newcommand{\Cov}{\operatorname{Cov}}
\newcommand{\norm}[1]{\left\lVert #1\right\rVert}
\newcommand{\abs}[1]{\left\lvert #1\right\rvert}
\newcommand{\set}[1]{\left\{#1\right\}}
\newcommand{\paren}[1]{\left(#1\right)}
\newcommand{\bracks}[1]{\left[#1\right]}
\newcommand{\eps}{\varepsilon}
\newcommand{\calO}{\mathcal O}
\newcommand{\given}{\,\middle|\,}
\newcommand{\doop}{\operatorname{do}}
\newcommand{\sem}[1]{\llbracket #1 \rrbracket}
\newcommand{\Inv}{\operatorname{Inv}}
\newcommand{\Pred}{\operatorname{Pred}}
\newcommand{\MechEq}{\equiv_{\mathrm m}}
\newcommand{\ObsEq}{\equiv_{\mathrm{obs}}}
\newcommand{\Trace}{\operatorname{Trace}}
\title{\Huge\bfseries NGeGenBrain : formalisme unifié de l'observation et de la reconstruction\\[0.6em]
\Large Empreintes informationnelles, invariants, variations, multi-échelle, temps et mécanismes exécutables}
\author{Théorie consolidée et audit mathématique}
\date{Version full\_0 -- 28 octobre 2026}
\begin{document}
\hypersetup{pageanchor=false}
\frontmatter
\maketitle
\thispagestyle{empty}
\begin{center}\begin{minipage}{0.92\textwidth}
\textbf{Objet du document.} Ce traité refond en une seule chaîne logique les fondations de l'observation par sondes, les empreintes informationnelles, la théorie inverse des mécanismes et la décomposition invariant--variant--dynamique. Il ne prétend pas qu'une observation finie détermine universellement l'intérieur d'un modèle : il donne les conditions exactes sous lesquelles une reconstruction est possible, unique, approchée ou non identifiable.

\medskip
\textbf{Règle de lecture scientifique.} Une définition fixe un vocabulaire ; un axiome ou une hypothèse ajoute une condition ; une proposition ou un théorème n'est déclaré démontré que sous ses hypothèses ; un résultat expérimental reste relatif au protocole exécuté. L'absence d'identifiabilité est un résultat et doit être conservée comme telle.

\medskip
\textbf{Principe central.} Une empreinte n'est pas nécessairement un invariant, n'est pas automatiquement suffisante pour reconstruire une observation, et n'identifie pas automatiquement un mécanisme. L'invariance, la suffisance et l'identification sont des propriétés à définir et à démontrer séparément.
\end{minipage}\end{center}
\clearpage
\tableofcontents
\hypersetup{pageanchor=true}
\mainmatter


\part{I. Fondations, types d’objets et sémantique des empreintes}

\chapter{Objet, portée et statuts épistémiques}
\section{Objet de la théorie}
La théorie ne postule pas qu'un modèle possède une essence immédiatement lisible. Elle étudie les différences que des protocoles d'observation déterminés peuvent révéler. Un objet est représenté par les réponses mesurées à une famille explicite de sondes. La représentation obtenue dépend donc de l'objet, du contexte, du protocole, de l'échelle, de la résolution des mesures et des ressources de l'observateur.

La chaîne de construction est
\[\begin{gathered}
\text{objet / trajectoire}\to\text{contexte}\to\text{sondes}\to\text{mesures}\to\text{empreinte},\\
\text{empreinte}\to\text{fibres / quotient}\to\text{seuils et dynamique},\\
\text{hypothèses}\to\text{interventions}\to\text{mécanismes candidats}\to\text{validation / révision}.
\end{gathered}\]
Chaque flèche est une opération distincte. En particulier, mesurer une activation ne démontre pas sa fonction; observer une différence après intervention ne démontre pas à elle seule un mécanisme complet; une reconstruction comportementale n'établit pas une équivalence mécaniste.

\section{Statuts épistémiques obligatoires}
Tout élément du registre doit porter au moins un statut : définition; convention; hypothèse explicite; mesure; observation; résultat expérimental; inférence; conjecture; mécanisme candidat; mécanisme soutenu; mécanisme réfuté ou révisé; théorème; démonstration; contre-exemple; limite; inconnue. Un même dossier peut contenir des éléments de statuts différents, mais un statut ne doit pas être élevé automatiquement sans satisfaire son critère de transition.

\begin{definition}[Observation et interprétation]
Une \emph{observation} est un enregistrement traçable d'une valeur ou d'un événement acquis sous un protocole déclaré. Une \emph{mesure} est une observation accompagnée d'une règle d'attribution de valeur, d'unités ou d'échelle et, lorsque pertinent, d'une incertitude. Une \emph{interprétation} est une proposition reliant des observations à une explication; elle n'est pas un fait brut.
\end{definition}

\begin{remarque}
Une sortie sérialisée sans contexte, identifiant de sonde, provenance, version du modèle et règle d'acquisition peut être un texte stocké, mais ce n'est pas une observation scientifique suffisamment définie pour soutenir une comparaison reproductible.
\end{remarque}

\section{Hypothèses de fond qui ne peuvent rester implicites}
Le formalisme exige de déclarer, pour chaque étude : (i) l'objet visé; (ii) le domaine des entrées et contextes; (iii) la famille de sondes; (iv) l'accès réel du système à ces sondes; (v) le type d'horloge et le régime temporel; (vi) l'échelle et la résolution; (vii) l'incertitude et les valeurs manquantes; (viii) la définition des interventions; (ix) la population sur laquelle une généralisation est revendiquée; (x) les critères de validation et de réfutation. Les théorèmes ci-dessous indiquent plus précisément lesquelles de ces hypothèses ils utilisent.

\chapter{Contrat mathématique commun et décisions d'harmonisation}
\section{Les objets qu'il est interdit de confondre}
La théorie unifiée distingue au minimum les objets suivants. Une exécution source est une trajectoire dans un espace d'exécutions. Une observation est le résultat d'un protocole appliqué à cette exécution. Un enregistrement est la représentation conservée du résultat, avec ses horodatages, statuts de mesure, provenance et incertitudes. Une empreinte est une application déterminée sur un espace d'observations ou d'objets. Un invariant est une propriété d'une application relativement à une famille de transformations déclarée. Un dossier d'évidence conserve les enregistrements et leur provenance. Un mécanisme est un objet exécutable avec une sémantique et des interfaces.

\begin{definition}[Observation induite par un protocole]
Soient $\Omega$ un espace d'exécutions et $\Ycal$ un espace d'enregistrements. Un protocole déterministe $\pi$ induit une application $O_\pi:\Omega\to\Ycal$. Si l'acquisition ou l'exécution est stochastique, elle est décrite au niveau approprié par un noyau probabiliste $K_\pi(\dd y\mid\omega)$ ; une valeur observée est alors une réalisation, et non le noyau lui-même.
\end{definition}
\begin{definition}[Empreinte relative à un protocole]
Pour un espace d'observations $\Ycal$, une empreinte est une application $I_\pi:\Ycal\to\Zcal$. L'empreinte composée d'une exécution est $F_\pi=I_\pi\circ O_\pi$ dans le cas déterministe. Le qualificatif « invariant » n'est ajouté que si une invariance sous une famille de transformations est explicitement définie et vérifiée.
\end{definition}
\begin{definition}[Dossier d'évidence]
Un dossier d'évidence $\mathfrak D$ est une collection versionnée d'enregistrements, de protocoles, de transformations, d'incertitudes, de résultats négatifs et de liens de provenance. Un dossier peut contenir plusieurs empreintes et les données qui les ont produites ; il ne se réduit pas, par définition, à une empreinte.
\end{definition}
\begin{definition}[Invariance déclarée]
Soit $G$ une famille de transformations agissant sur $\Ycal$. Une application $I:\Ycal\to\Zcal$ est invariante sous $G$ si, chaque fois que $g\cdot y$ est défini, $I(g\cdot y)=I(y)$ pour tout $g\in G$ et $y$ de son domaine déclaré. Toute affirmation d'invariance doit nommer la famille $G$, le domaine et la relation d'égalité considérée.
\end{definition}
\section{Propriétés distinctes}
Les affirmations suivantes sont indépendantes en général : (i) une signature est bien définie ; (ii) elle est invariante ; (iii) elle est injective ; (iv) elle est suffisante pour une propriété cible ; (v) elle identifie une classe de mécanismes ; (vi) elle est estimée avec précision à partir d'un échantillon ; (vii) elle est calculable dans un budget fini. Aucune de ces propriétés ne découle des autres sans hypothèse supplémentaire.

\begin{proposition}[Invariance et reconstruction sont des propriétés distinctes]
Une application invariante non constante peut être non injective ; une application injective n'est pas nécessairement invariante sous une famille de transformations non triviale.
\end{proposition}
\begin{proof}
Pour la première assertion, prenons $\Ycal=\mathbb R^2$, l'action $g_a(x,z)=(ax,z)$ pour $a>0$ et $I(x,z)=z$. Alors $I$ est invariant sous cette action mais n'est pas injectif. Pour la seconde, sur $\Ycal=\mathbb R$, considérons une transformation $g(y)=y+1$ et $I(y)=y$ : $I$ est injectif, mais $I(g(y))=y+1\ne y=I(y)$.\end{proof}
\section{Règles d'harmonisation appliquées dans le traité}
\begin{enumerate}[label=\arabic*.]
\item Le mot \emph{empreinte} désigne une application ou une famille indexée d'applications. Le mot \emph{invariant} désigne une propriété prouvée relativement à une famille de transformations, et non une essence supposée.
\item Une trace temporelle, un enregistrement, une empreinte et un dossier d'évidence sont des types d'objets distincts. Une réduction d'un dossier à une empreinte est une opération dont la perte d'information doit être évaluée.
\item Un seuil, une distance ou une hiérarchie n'est défini qu'après avoir précisé domaine, codomaine, ordre, convention de bord et cas infinis.
\item Une égalité exacte, une égalité presque sûre, une borne d'erreur et une non-différence statistiquement détectée ne sont jamais interchangeables.
\item Une intervention identifie d'abord un effet sous un protocole précis. Une interprétation causale plus large exige des hypothèses d'assignation, de contrôle, de mesure et de transport déclarées.
\item Une composition syntaxique de modules ne garantit ni la compatibilité de leurs interfaces, ni l'autonomie d'exécution, ni la fidélité au système source.
\item Les résultats établis sur une classe finie de tests ne sont pas des garanties universelles sur un domaine infini.
\item Un hash vérifie une intégrité relative à une référence fiable ; il ne prouve ni la vérité d'une donnée ni la validité d'une interprétation.
\item Les fonctions de perte, les scores agrégés et les résumés n'effacent pas les résultats par facette, les contre-exemples, ni les candidats compatibles alternatifs.
\item Une affirmation de réussite de NGeGenBrain doit mentionner le domaine, le protocole, les conditions d'exécution et les niveaux non établis.
\end{enumerate}
\section{Décision sur les sources}
Les trois documents sources ne sont pas traités comme des autorités infaillibles. Leurs définitions et démonstrations sont réintégrées dans l'argument commun ; les répétitions documentaires sont déplacées en annexes ; les divergences terminologiques sont signalées ; les corrections locales identifiées lors de la consolidation sont consignées dans l'audit. La traceabilité en fin de volume indique, pour chaque chapitre source, son emplacement ou sa conservation documentaire.

\chapter{Objets, contextes, sondes et mesures}
\section{Objets et contextes}
\begin{definition}[Objet étudié]
Un objet étudié est un élément $x\in\X$. Selon l'expérience, $x$ peut être un modèle complet, un sous-module, une représentation, un état d'exécution, une version paramétrique, une règle candidate ou une trajectoire. Le type d'objet est déclaré dans le protocole; on ne mélange pas silencieusement des objets de types différents.
\end{definition}

\begin{definition}[Contexte d'observation]
Un contexte $p\in\Pcal$ est le tuple des facteurs contrôlés ou enregistrés qui peuvent influencer la réponse : entrée, tokenizer, masque, longueur de contexte, paramètres de génération, graine, précision numérique, appareil, état du modèle, version du logiciel, environnement, ordre d'exécution, intervention et configuration de la sonde. Le contexte comprend les variables pertinentes au test revendiqué; le fait d'en oublier une est une limite de l'étude, pas une justification pour la traiter comme constante.
\end{definition}

\begin{definition}[Sonde]
Une sonde $f\in\Fcal$ est un opérateur d'observation spécifié par un identifiant et une version, un domaine d'entrée, un codomaine, une procédure d'acquisition, des préconditions, un coût ou budget estimé, un niveau d'intrusion, une règle de sérialisation et une procédure de contrôle. Une sonde peut être passive ou active; une intervention active est aussi décrite par un opérateur d'action distinct et le capteur qui mesure ses conséquences.
\end{definition}

\begin{definition}[Application d'observation]
Pour chaque triplet $(p,x,f)$ admissible, une application d'observation assigne une valeur
\[
\Phi(p,x,f)\in Y_f,
\]
où $Y_f$ est le codomaine propre de la sonde. Si une sonde renvoie une erreur, une valeur absente ou un résultat censuré, cet état doit être représenté explicitement et ne peut être remplacé silencieusement par zéro. Dans le cas scalaire, $Y_f=\R$ et l'on retrouve la forme $\Phi:\Pcal\times\X\times\Fcal\to\R$.
\end{definition}

\begin{definition}[Protocole admissible]
Un protocole $\pi$ est admissible pour $(p,x,f)$ si les préconditions de $f$ sont satisfaites, si l'accès requis est déclaré disponible, si l'exécution est enregistrée et si les défaillances d'acquisition sont signalées. L'application $\Phi$ n'est appelée que sur des triplets admissibles; hors domaine, le résultat est « non défini » ou « non disponible », et non une valeur inventée.
\end{definition}

\section{Mesure exacte, mesure bruitée et modèles stochastiques}
\begin{definition}[Modèle de mesure]
Une observation déterministe exacte renvoie $y=\Phi(p,x,f)$. Une mesure bruitée est représentée par une loi de mesure $Q_f(\cdot\mid p,x)$ sur $Y_f$, ou par un modèle explicite $Y=\Phi(p,x,f)+\epsilon$ lorsque le codomaine vectoriel le permet. La loi d'erreur, les bornes ou la distribution de $\epsilon$ doivent être précisées avant toute garantie de précision. Un modèle neuronal stochastique peut lui-même produire une loi de sortie : cette variabilité est distincte du bruit du capteur.
\end{definition}

\begin{hypothese}[Admissibilité des égalités]
Les égalités exactes de signatures ne sont utilisées que lorsque les sorties sont exactement comparables selon une représentation canonique déclarée. En présence de quantification, bruit, non-déterminisme ou arrondi, l'étude doit annoncer soit une comparaison des lois, soit une relation d'équivalence tolérante munie d'une métrique et d'un seuil choisis à l'avance.
\end{hypothese}

\begin{definition}[Équivalence tolérante]
Si chaque coordonnée observée appartient à un espace métrique $(Y_f,d_f)$, une signature à tolérance $\eps$ est dite compatible lorsque
\[
I_{\Gcal}^{p}(x)\approx_{\eps} I_{\Gcal}^{p}(y)
\iff \forall f\in\Gcal,\ d_f(\Phi(p,x,f),\Phi(p,y,f))\leq\eps_f.
\]
Cette relation est réflexive et symétrique; elle n'est pas nécessairement transitive. Elle ne doit donc pas être appelée relation d'équivalence sans une construction supplémentaire, par exemple une fermeture transitive explicitement documentée.
\end{definition}

\chapter{Empreinte informationnelle et identité observationnelle}
\section{Construction de l'empreinte}
\begin{definition}[Empreinte finie]
Pour une famille finie de sondes $\Gcal_K=\{f_1,\ldots,f_K\}$ et un contexte $p$, l'empreinte informationnelle est
\[
I_{\Phi,\Gcal_K}^{p}:\X_{p,\Gcal_K}\to\prod_{j=1}^{K}Y_{f_j},\qquad
x\mapsto\bigl(\Phi(p,x,f_1),\ldots,\Phi(p,x,f_K)\bigr),
\]
où $\X_{p,\Gcal_K}$ est le domaine des objets pour lesquels toutes les sondes de la famille sont admissibles. Lorsque $K$ est implicite, on écrit $I_{\Gcal}^{p}$.
\end{definition}

\begin{definition}[Empreinte d'une famille arbitraire]
Pour $\Gcal\subseteq\Fcal$, on peut représenter la signature par la fonction
\[
I_{\Phi,\Gcal}^{p}:\X_{p,\Gcal}\to\prod_{f\in\Gcal}Y_f,\qquad
I_{\Phi,\Gcal}^{p}(x)(f)=\Phi(p,x,f).
\]
Si la famille est infinie, l'empreinte est une fonction sur $\Gcal$; elle n'est pas réputée calculable en temps fini sans une hypothèse de représentation effective.
\end{definition}

\begin{definition}[Indiscernabilité observationnelle]
À contexte $p$ et famille de sondes $\Gcal$ fixés, on définit
\[
x\sim_{p,\Gcal}y\quad\Longleftrightarrow\quad I_{\Phi,\Gcal}^{p}(x)=I_{\Phi,\Gcal}^{p}(y).
\]
Deux objets sont indiscernables relativement à ce protocole. Cette relation ne signifie pas que les objets sont intrinsèquement identiques; elle signifie uniquement que la famille actuelle de sondes ne les sépare pas.
\end{definition}

\begin{proposition}[L'indiscernabilité est une relation d'équivalence]
Sur le domaine commun $\X_{p,\Gcal}$, la relation $\sim_{p,\Gcal}$ est réflexive, symétrique et transitive.
\end{proposition}
\begin{proof}
Pour tout $x$, $I(x)=I(x)$, donc la relation est réflexive. Si $I(x)=I(y)$, alors $I(y)=I(x)$, donc elle est symétrique. Si $I(x)=I(y)$ et $I(y)=I(z)$, la transitivité de l'égalité donne $I(x)=I(z)$. Aucune autre hypothèse n'est requise que le caractère bien défini des signatures sur un domaine commun.
\end{proof}

\begin{definition}[Fibres et quotient]
La fibre d'une signature en $v$ est $I^{-1}(v)=\{x\in\X:I(x)=v\}$. Le quotient observationnel $\X/{\sim_{p,\Gcal}}$ est l'ensemble des classes d'indiscernabilité. L'application canonique $q_{p,\Gcal}:\X\to\X/{\sim_{p,\Gcal}}$ envoie $x$ sur sa classe $[x]$. L'image observée $\Lcal_{p,\Gcal}=I_{\Phi,\Gcal}^{p}(\X_{p,\Gcal})$ est l'ensemble des signatures effectivement possibles dans le domaine déclaré.
\end{definition}

\begin{theoreme}[Raffinement monotone des empreintes]
Soient $\Gcal_1\subseteq\Gcal_2\subseteq\Fcal$ deux familles de sondes, toutes admissibles sur le même domaine $D$. Alors
\[
x\sim_{p,\Gcal_2}y\Longrightarrow x\sim_{p,\Gcal_1}y.
\]
La partition induite par $\Gcal_2$ raffine donc celle induite par $\Gcal_1$ : ajouter des coordonnées ne peut fusionner deux classes auparavant distinctes.
\end{theoreme}
\begin{proof}
L'égalité des signatures sur $\Gcal_2$ signifie l'égalité de toutes les coordonnées $f\in\Gcal_2$. Comme $\Gcal_1\subseteq\Gcal_2$, elle implique l'égalité sur chaque $f\in\Gcal_1$, donc celle des signatures restreintes. En contraposée, si une sonde de $\Gcal_1$ sépare deux objets, aucune extension de la famille ne peut faire disparaître cette différence tant que les valeurs de cette sonde sont conservées.
\end{proof}

\begin{corollaire}[Critère d'une nouvelle distinction]
L'ajout d'une sonde $g$ raffine strictement la partition si et seulement s'il existe $x,y\in D$ tels que $I_{\Gcal}^{p}(x)=I_{\Gcal}^{p}(y)$ mais $\Phi(p,x,g)\neq\Phi(p,y,g)$.
\end{corollaire}
\begin{proof}
Si une telle paire existe, sa classe actuelle est scindée par la nouvelle coordonnée. Inversement, si la partition est strictement raffinée, une classe s'est divisée; deux éléments de cette classe ont la même ancienne signature et des valeurs différentes pour au moins une coordonnée ajoutée, notamment $g$ dans le cas d'une seule nouvelle sonde.
\end{proof}

\section{Théorème de factorisation}
\begin{theoreme}[Factorisation par l'empreinte]
Soient $I=I_{\Phi,\Gcal}^{p}:D\to\Lcal$ et $A:D\to Y$ une application. Les trois assertions suivantes sont équivalentes :
\begin{enumerate}[label=(\roman*)]
\item $A$ est constante sur chaque fibre de $I$;
\item il existe une application $\bar A:\Lcal\to Y$ telle que $A=\bar A\circ I$;
\item il existe une application $\widehat A:D/{\sim_I}\to Y$ telle que $A=\widehat A\circ q$.
\end{enumerate}
Quand elle existe, $\bar A$ est unique sur $\Lcal=I(D)$, et $\widehat A$ est unique sur le quotient.
\end{theoreme}
\begin{proof}
(i)$\Rightarrow$(ii). Pour $v\in\Lcal$, choisissons $x$ tel que $I(x)=v$ et posons $\bar A(v)=A(x)$. Si $I(x)=I(x')$, alors $x,x'$ appartiennent à la même fibre et (i) donne $A(x)=A(x')$; la définition ne dépend donc pas du représentant. Par construction $A(x)=\bar A(I(x))$.

(ii)$\Rightarrow$(i). Si $I(x)=I(y)$, alors $A(x)=\bar A(I(x))=\bar A(I(y))=A(y)$.

(i)$\Leftrightarrow$(iii) par le même raisonnement, en remplaçant chaque fibre par sa classe $[x]$. Enfin, l'unicité sur l'image suit parce que toute valeur $v\in\Lcal$ possède au moins un représentant $x$ et la factorisation impose $\bar A(v)=A(x)$; la constance sur les fibres garantit que ce choix est unique.
\end{proof}

\begin{remarque}[Portée]
Le théorème n'affirme pas que l'empreinte contient toute l'information imaginable sur l'objet. Il caractérise exactement les fonctions $A$ calculables à partir de l'empreinte choisie. Si une propriété utile varie à l'intérieur d'une fibre, cette famille de sondes ne suffit pas à la déterminer.
\end{remarque}

\chapter{Pourquoi l'empreinte ne doit pas être confondue avec l'observation entière}
\section{Objets à distinguer}
On distingue quatre objets qui ne sont pas interchangeables : (i) l'exécution source et sa trajectoire ; (ii) l'enregistrement de cette trajectoire sous un protocole donné ; (iii) l'empreinte, qui est une application choisie sur cet enregistrement ; (iv) un mécanisme candidat qui peut produire des enregistrements. Une empreinte peut être exhaustive relativement à un protocole et néanmoins insuffisante pour reconstruire la trajectoire brute.

\begin{definition}[Enregistrement observationnel]
Soit \(\Omega\) l'espace des trajectoires possibles et \(\Y\) l'espace des enregistrements admissibles. Un protocole fixé \(\pi\) induit une application d'observation \(O_\pi:\Omega\to\Y\). En présence d'aléas de mesure ou d'exécution, \(O_\pi(\omega)\) désigne une réalisation d'un noyau probabiliste plutôt qu'une application déterministe.
\end{definition}
\begin{definition}[Empreinte relative]
Une empreinte relative à \(\pi\) est une application \(I_\pi:\Y\to\Z\), où \(\Z\) est son espace de valeurs. L'empreinte composée sur une trajectoire est \(F_\pi=I_\pi\circ O_\pi\). Le terme « empreinte » ne signifie ni « invariant » ni « suffisant » sans une propriété supplémentaire prouvée.
\end{definition}
\begin{definition}[Invariance déclarée]
Soit \(\mathcal G\) un ensemble de transformations de \(\Y\). Une application \(I:\Y\to\Z\) est \emph{invariante sous} \(\mathcal G\) si, pour tous \(y\in\Y\) et \(g\in\mathcal G\), \(I(g\cdot y)=I(y)\), lorsque l'action \(g\cdot y\) est définie. L'invariance est toujours relative à \(\mathcal G\) et au domaine où l'égalité est évaluée.
\end{definition}
\begin{proposition}[Invariance sous transformations et constance sur les fibres sont distinctes]
Soient $P$ un espace de contextes, $\Y$ un domaine d'enregistrements, $\Fcal_0$ une famille non vide de sondes et $W_f$ le codomaine de chaque sonde $f$. Pour un contexte fixé $p\in P$, soit
\[
 I_{p,\Fcal_0}(y)=\bigl(\Phi(p,y,f)\bigr)_{f\in\Fcal_0}\in\prod_{f\in\Fcal_0}W_f.
\]
Soit $\mathcal A$ une famille déclarée d'applications $a:\Y\to\Y$. Alors $I_{p,\Fcal_0}$ est invariant sous $\mathcal A$ si et seulement si
\[
 \forall a\in\mathcal A\ \forall y\in\Y\ \forall f\in\Fcal_0,\qquad
 \Phi(p,a(y),f)=\Phi(p,y,f).
\]
En revanche, $I_{p,\Fcal_0}$ est toujours constante sur chacune de ses propres fibres, par définition de la relation $y\sim y'\iff I_{p,\Fcal_0}(y)=I_{p,\Fcal_0}(y')$. Cette dernière propriété tautologique ne démontre aucune invariance sous une famille de transformations non spécifiée. Fixer le contexte et la famille de sondes rend la signature bien définie; cela ne garantit pas qu'elle soit préservée si l'objet ou le contexte change.
\end{proposition}
\begin{proof}
L'égalité de deux éléments du produit $\prod_{f\in\Fcal_0}W_f$ équivaut à l'égalité de toutes leurs coordonnées. En appliquant cette caractérisation à $I_{p,\Fcal_0}(a(y))=I_{p,\Fcal_0}(y)$, on obtient exactement la condition affichée, dans les deux sens. La constance sur les fibres résulte de leur définition : deux éléments d'une même fibre ont par définition la même valeur de $I_{p,\Fcal_0}$. Elle n'implique pas que $a(y)$ appartienne à la même fibre que $y$ pour une transformation arbitraire $a$.
\end{proof}
\begin{contreexemple}[Une sonde fixée ne garantit pas l'invariance par transformation]
Prenons $\Y=\mathbb R$, une seule sonde $f$ définie par $\Phi(p,y,f)=y$, et la transformation $a(y)=y+1$. L'empreinte vaut $I_{p,\{f\}}(y)=y$, donc $I_{p,\{f\}}(a(y))=y+1\ne y=I_{p,\{f\}}(y)$. La famille de sondes est pourtant fixée. Ainsi, l'invariance ne suit pas du seul fait que le protocole utilise un nombre fini ou déterminé de sondes.
\end{contreexemple}
\begin{remarque}[Sens rigoureux de l'invariance contextuelle]
Si par « invariance contextuelle » on entend que la signature reste égale lorsque le contexte varie dans un ensemble $P_0$, il faut vérifier séparément $I_{p,\Fcal_0}(y)=I_{p',\Fcal_0}(y)$ pour tous $p,p'\in P_0$ et les $y$ du domaine revendiqué. Si l'on entend plutôt la constance sur les classes induites par une signature à contexte fixé, le terme précis est « constance sur les fibres observationnelles ». Les deux notions peuvent être utiles, mais elles ne sont pas équivalentes.
\end{remarque}
\begin{remarque}
Une moyenne stable entre contextes peut être un estimateur d'invariant, mais ne devient pas un invariant exact par le seul fait qu'elle varie peu sur un échantillon fini. Il faut préciser si l'on revendique une propriété exacte, une borne d'erreur ou seulement une régularité empirique.
\end{remarque}

\section{Correction d'une intuition partiellement juste}
L'intuition « l'empreinte capte les invariants » peut être utile pour concevoir un système, mais elle n'est pas vraie par définition. Une empreinte peut contenir des variations brutes; à l'inverse, une empreinte volontairement invariante peut les éliminer. La théorie distingue donc : (a) empreinte générale ; (b) empreinte invariantielle ; (c) empreinte multi-facettes ; (d) empreinte de trajectoire ; (e) dossier d'évidence, plus riche que toute empreinte compressée.

\begin{proposition}[Une empreinte non injective ne permet pas de reconstruire chaque observation]
Soient \(I:\Y\to\Z\) et \(y_1\ne y_2\) tels que \(I(y_1)=I(y_2)\). Il n'existe aucune application \(R:\Z\to\Y\) telle que \(R(I(y))=y\) simultanément pour tout \(y\in\{y_1,y_2\}\).
\end{proposition}
\begin{proof}
Si un tel \(R\) existait, alors \(y_1=R(I(y_1))=R(I(y_2))=y_2\), contradiction avec \(y_1\ne y_2\).\end{proof}
\begin{corollaire}[Condition exacte de reconstruction]
Il existe une application de reconstruction $R:I(\Y)\to\Y$ satisfaisant $R\circ I=\Id_\Y$ si et seulement si $I$ est injective. Pour reconstruire seulement une propriété cible $A:\Y\to\W$, la condition nécessaire et suffisante est que $A$ soit constante sur chaque fibre de $I$. Le domaine de $R$ est explicitement restreint à $I(\Y)$, qui contient toutes les signatures produites; une extension aux valeurs de $\Z\setminus I(\Y)$ n'est pas requise pour reconstruire les observations effectivement représentées.
\end{corollaire}
\begin{proof}
Si $R\circ I=\Id_\Y$, alors $I(y)=I(y')$ implique $y=R(I(y))=R(I(y'))=y'$, donc $I$ est injective. Réciproquement, si $I$ est injective, chaque $z\in I(\Y)$ possède un unique antécédent; poser $R(z)$ égal à cet antécédent définit bien une application et donne $R\circ I=\Id_\Y$. Pour la propriété cible, si $A=\bar A\circ I$, elle est constante sur les fibres. Réciproquement, si $A$ est constante sur les fibres, définir $\bar A(I(y))=A(y)$ est indépendant du choix de $y$ et fournit la factorisation.\end{proof}

\chapter{Quotient invariant et information résiduelle}
\section{Relation d'indiscernabilité structurale}
\begin{definition}[Fibre invariantielle]
Pour une application invariante \(I:\Y\to\Z\), la fibre de \(y\) est \([y]_I=\{y'\in\Y:I(y')=I(y)\}\). Elle contient exactement les enregistrements que l'empreinte ne distingue pas.
\end{definition}
\begin{proposition}[Une empreinte invariante factorise par le quotient]
La relation \(y\sim_I y'\iff I(y)=I(y')\) est une relation d'équivalence. L'application \(I\) induit une bijection entre le quotient \(\Y/\!\sim_I\) et l'image \(I(\Y)\).
\end{proposition}
\begin{proof}
Réflexivité, symétrie et transitivité découlent de l'égalité dans \(\Z\). L'application induite envoie la classe \([y]\) sur \(I(y)\). Elle est bien définie par définition des classes, surjective sur l'image et injective puisque deux classes ayant la même image coïncident.\end{proof}
\begin{definition}[Coordonnée variante ou résiduelle]
Une coordonnée variante est une application \(V:\Y\to\U\) conçue pour distinguer certaines différences que \(I\) oublie. Elle peut encoder le contexte, la sonde, l'intervention, l'état temporel, la position dans la trajectoire, la réalisation stochastique ou un résidu mesuré. Le nom « variant » n'implique pas qu'elle soit bruit pur ou négligeable.
\end{definition}
\begin{theoreme}[Critère de reconstruction par empreinte et variante]
La paire \((I,V):\Y\to\Z\times\U\) permet de reconstruire exactement chaque \(y\in\Y\) si et seulement si elle est injective. De façon équivalente, pour tous \(y,y'\), \(I(y)=I(y')\) et \(V(y)=V(y')\) impliquent \(y=y'\).
\end{theoreme}
\begin{proof}
Appliquer le corollaire de reconstruction à l'application produit \(J(y)=(I(y),V(y))\). La condition écrite est exactement l'injectivité de \(J\).\end{proof}
\begin{remarque}
Ce théorème ne prétend pas que la séparation invariant/variant existe automatiquement. Il donne le critère exact qu'une telle décomposition doit satisfaire. La construction d'un \(V\) suffisamment discriminant constitue un problème effectif, à résoudre par le protocole et à valider.
\end{remarque}

\chapter{Quand une décomposition additive est-elle définie de façon unique ?}
\section{Non-unicité sans contrainte}
\begin{contreexemple}[Décomposition additive non identifiable]
Pour toute donnée \(y\in\R\), les couples \((s,v)=(y,0)\), \((0,y)\) et \((y-c,c)\) pour tout \(c\in\R\) satisfont \(y=s+v\). L'équation seule ne définit donc ni un invariant ni un variant unique.
\end{contreexemple}
\begin{proposition}[Unicité par somme directe]
Soient \(U,W\) deux sous-espaces vectoriels d'un espace vectoriel \(E\), avec \(U\cap W=\{0\}\). Si \(y=s+v=s'+v'\), où \(s,s'\in U\) et \(v,v'\in W\), alors \(s=s'\) et \(v=v'\). Une décomposition existe pour chaque \(y\in U+W\).
\end{proposition}
\begin{proof}
De \(s+v=s'+v'\), on déduit \(s-s'=v'-v\). Le membre gauche appartient à \(U\) et le droit à \(W\); leur valeur commune appartient donc à \(U\cap W=\{0\}\). Ainsi \(s=s'\) et \(v=v'\). L'existence pour \(y\in U+W\) découle de la définition de cette somme.\end{proof}
\begin{remarque}
L'hypothèse \(U\cap W=\{0\}\) est un contrat de modèle, non un fait découvert automatiquement. Dans une application statistique, elle peut être remplacée par une condition d'identifiabilité explicitement vérifiée, par exemple une matrice de conception de rang plein; sinon l'ambiguïté doit être conservée.
\end{remarque}

\section{Invariance, stabilité temporelle et stationnarité ne sont pas synonymes}
\begin{definition}[Trois prédicats distincts]
Une application $q:X\to Z$ est invariante sous une famille déclarée $\mathcal G$ de transformations de $X$ si $q(gx)=q(x)$ pour tout $g\in\mathcal G$ et tout $x$ où l'action est définie. Une trajectoire observée $s:T\to Z$ est $\varepsilon$-stable sur une fenêtre $J\subseteq T$ relativement à une métrique $d$ si $\sup_{t,u\in J}d(s(t),s(u))\leq\varepsilon$. Un processus stochastique est stationnaire au sens strict si toutes ses distributions finies sont invariantes par translation temporelle admissible; d'autres notions, comme la stationnarité au second ordre, doivent être nommées séparément. Ces prédicats ont des définitions et des quantificateurs différents; aucune implication entre eux n'est présumée par la seule terminologie.
\end{definition}
\begin{proposition}[L'invariance sous une action d'état n'implique pas la constance temporelle]
Il existe un espace d'état $X$, une action non triviale $\mathcal G$, une application $q:X\to\mathbb R$ invariante sous $\mathcal G$ et une trajectoire $x:T\to X$ telle que $t\mapsto q(x(t))$ ne soit pas constante.
\end{proposition}
\begin{proof}
Prenons $X=\mathbb R^2$, $\mathcal G=\{g_a:a>0\}$ avec $g_a(u,v)=(au,v)$ et $q(u,v)=v$. Pour tout $a>0$, $q(g_a(u,v))=v=q(u,v)$, donc $q$ est invariant sous cette action non triviale. Pour $T=\mathbb R$, définissons $x(t)=(0,t)$. Alors $q(x(t))=t$, qui n'est pas constant. L'invariance de $q$ sous $\mathcal G$ n'implique donc pas la constance temporelle de son observation le long d'une trajectoire.
\end{proof}
\begin{remarque}[Implications entre prédicats]
Une implication entre invariance, stabilité sur une fenêtre et stationnarité ne peut être affirmée qu'après avoir spécifié les espaces, l'action, la fenêtre, la métrique, la loi du processus et les hypothèses qui les relient. Cette section ne pose aucune implication générale supplémentaire.
\end{remarque}
\begin{remarque}
Pour un LLM, un mécanisme stable peut produire des états cachés différents à chaque token. L'invariant pertinent peut être la loi de transition, le contrat de calcul ou la relation de dépendance, non une valeur d'activation constante.
\end{remarque}

\chapter{Redondance et fermeture observationnelle}
\section{Équivalence des sondes}
\begin{definition}[Équivalence de sondes]
À contexte et domaine fixés, deux sondes $f,g$ de même codomaine sont équivalentes si
\[
f\equiv_p g\quad\Longleftrightarrow\quad\forall x\in D,\ \Phi(p,x,f)=\Phi(p,x,g).
\]
Si leurs codomaines diffèrent, l'égalité doit être remplacée par une correspondance canonique déclarée. Une égalité observée sur un échantillon fini ne démontre pas l'équivalence universelle; elle établit uniquement l'équivalence empirique sur cet échantillon.
\end{definition}

\begin{proposition}[Retrait d'une sonde redondante]
Si $g\equiv_p f$ sur tout le domaine $D$, alors remplacer $g$ par $f$ dans une famille de sondes ne change ni l'empreinte à permutation et codage canonique près, ni les fibres induites.
\end{proposition}
\begin{proof}
Pour tout $x\in D$, la coordonnée $g$ est égale à la coordonnée $f$. Ajouter $g$ à une signature contenant $f$ ne crée donc aucune nouvelle distinction; les égalités des signatures avant et après ajout sont équivalentes.
\end{proof}

\begin{definition}[Fermeture observationnelle]
Pour $\Gcal\subseteq\Fcal$, définissons
\[
\cl_p(\Gcal)=\{f\in\Fcal:\exists h_f:\Lcal_{p,\Gcal}\to Y_f\text{ telle que }\forall x\in D,\ \Phi(p,x,f)=h_f(I_{p,\Gcal}(x))\}.
\]
C'est la famille des sondes dont la réponse est entièrement déterminée par l'empreinte déjà disponible sur $D$. Cette définition est globale sur $D$; une fermeture empirique sur un jeu de données doit être notée séparément et n'autorise aucune généralisation hors échantillon sans hypothèse supplémentaire.
\end{definition}

\begin{theoreme}[La fermeture observationnelle est un opérateur de fermeture]
Sous réserve que chaque sonde de $\Fcal$ soit définie sur le même domaine $D$, l'application $\cl_p$ est extensive, monotone et idempotente :
\[
\Gcal\subseteq\cl_p(\Gcal),\qquad
\Gcal\subseteq\Hcal\Rightarrow\cl_p(\Gcal)\subseteq\cl_p(\Hcal),\qquad
\cl_p(\cl_p(\Gcal))=\cl_p(\Gcal).
\]
Elle induit les mêmes fibres que $\Gcal$ et est maximale parmi les familles qui conservent ces fibres.
\end{theoreme}
\begin{proof}
\emph{Extensivité.} Pour $g\in\Gcal$, la réponse $\Phi(p,x,g)$ est la projection de coordonnée $g$ de $I_{p,\Gcal}(x)$; elle se factorise donc à travers cette empreinte.

\emph{Monotonie.} Si $\Gcal\subseteq\Hcal$, il existe une projection de restriction $r:\Lcal_{p,\Hcal}\to\Lcal_{p,\Gcal}$ telle que $I_{p,\Gcal}=r\circ I_{p,\Hcal}$. Si $f$ se factorise par $I_{p,\Gcal}$ via $h_f$, alors $\Phi_f=h_f\circ r\circ I_{p,\Hcal}$; ainsi $f\in\cl_p(\Hcal)$.

\emph{Fibres.} L'extensivité implique que l'empreinte de fermeture contient toutes les coordonnées initiales, donc elle ne peut fusionner des fibres. Chaque coordonnée ajoutée par fermeture est, par définition, une fonction des coordonnées initiales, donc elle ne peut scinder une fibre initiale. Les fibres sont donc identiques.

\emph{Idempotence.} Puisque $\cl_p(\Gcal)$ et $\Gcal$ induisent la même fibre, leurs empreintes déterminent mutuellement leurs coordonnées : toute sonde déterminée par l'empreinte fermée est déterminée par l'empreinte initiale, et inversement par extensivité. Les deux fermetures sont donc égales.

\emph{Maximalité.} Si une famille $\Hcal$ conserve exactement les fibres de $\Gcal$, toute sonde $f\in\Hcal$ est constante sur chaque fibre de $I_{\Gcal}$. Le théorème de factorisation fournit alors $h_f$; donc $f\in\cl_p(\Gcal)$. Toute famille conservant ces fibres est incluse dans la fermeture.
\end{proof}

\begin{remarque}[Limite de calculabilité]
La définition est sémantique : elle quantifie sur tout $D$ et sur l'existence d'une fonction $h_f$. Dans un domaine infini, cette propriété n'est généralement pas décidée par un nombre fini d'observations sans hypothèse de classe de fonctions, borne de complexité ou oracle. Le logiciel doit séparer la fermeture théorique, une approximation apprise et une vérification empirique.
\end{remarque}

\chapter{De l'empreinte à la reconstruction : objets et statuts}
\section{Données disponibles et objet recherché}
\begin{definition}[Système source]
Un système source $S$ est l'objet étudié. Dans l'application visée, il peut être un LLM complet, un sous-réseau, un état d'exécution ou un ensemble de composants. On ne suppose pas que sa description interne soit intégralement accessible. Les accès réellement disponibles sont consignés dans le protocole.
\end{definition}

\begin{definition}[Protocole expérimental]
Un protocole $\pi$ est le tuple
\[
\pi=(\Pcal,\Ucal,\Acal,\Fcal,\mathcal Q,\mathcal B,\mathcal V),
\]
où $\Pcal$ est l'ensemble des contextes, $\Ucal$ les entrées, $\Acal$ les interventions admissibles, $\Fcal$ les sondes, $\mathcal Q$ les règles d'ordonnancement et de mesure, $\mathcal B$ les budgets et limites d'acquisition, et $\mathcal V$ les versions des modèles, sondes, données et logiciels. Une composante non déclarée est inconnue, non réputée constante.
\end{definition}

\begin{definition}[Dossier d'évidence]
Un dossier d'évidence $D$ est un ensemble fini ou dénombrable d'enregistrements
\[
e=(\mathrm{id},p,u,a,f,y,t,\rho,\eta),
\]
où $p$ est le contexte, $u$ l'entrée, $a$ l'intervention ou $\varnothing$, $f$ la sonde, $y$ la valeur ou le statut d'échec, $t$ les informations temporelles, $\rho$ la provenance et $\eta$ les métadonnées d'incertitude. Un enregistrement peut contenir une trajectoire et non une seule valeur. Un résultat absent, censuré ou invalide conserve un statut explicite.
\end{definition}

\begin{definition}[Dossier empirique structuré]
Le dossier empirique structuré dérivé de $D$ est la représentation
\[
\mathfrak D_D=\bigl(\pi,D,\text{règles de comparaison},\text{incertitudes déclarées}\bigr).
\]
Il conserve, selon le contrat de stockage déclaré, les conditions d’acquisition, les essais négatifs, les interventions et la provenance. Il ne doit pas être confondu avec une empreinte au sens strict, qui est une application (ou famille d’applications) calculée à partir d’un objet d’observation. Toute réduction des données brutes en résumés est déclarée comme potentiellement irréversible.
\end{definition}

\begin{definition}[Classe de mécanismes admissibles]
La classe $\Mcal$ est un ensemble explicitement représenté de modèles candidats. Chaque $M\in\Mcal$ possède une syntaxe ou sérialisation, une sémantique exécutable, des interfaces d'entrée/sortie, des règles d'intervention et une version. La classe peut être restreinte par des contraintes architecturales; ces contraintes doivent être annoncées avant l'évaluation et leur effet sur l'identifiabilité doit être documenté.
\end{definition}

\section{Pourquoi l'inverse est une relation avant d'être une fonction}
\begin{definition}[Compatibilité d'un mécanisme avec un dossier]
Pour un protocole $\pi$, une classe $\Mcal$ et un critère de compatibilité $C$, on définit
\[
\Inv_{\pi,C}(D)=\{M\in\Mcal:C(M,D,\pi)=\mathrm{vrai}\}.
\]
Cet ensemble est appelé \emph{ensemble inverse} ou \emph{ensemble des mécanismes compatibles}. Pour des mesures exactes déterministes, $C$ peut imposer l'égalité aux observations. Pour des mesures bruitées ou stochastiques, $C$ doit préciser un modèle de mesure, un score, une tolérance ou un niveau de confiance. Ces deux cas ne sont pas interchangeables.
\end{definition}

\begin{proposition}[L'inverse n'est pas une fonction sans condition d'unicité]
Soit $D$ un dossier et $\Inv_{\pi,C}(D)$ l'ensemble inverse associé. Une fonction de reconstruction $R(D)\in\Mcal$ qui prétend donner \emph{le mécanisme identifié par $D$} n'est justifiée par la seule compatibilité que si l'ensemble inverse détermine une classe unique selon une relation d'équivalence mécaniste déclarée. Si plusieurs classes non équivalentes appartiennent à l'ensemble inverse, tout choix d'un représentant nécessite une règle supplémentaire et ne découle pas des données seules.
\end{proposition}
\begin{proof}
Par définition, chaque $M\in\Inv_{\pi,C}(D)$ satisfait le critère de compatibilité. Si deux éléments $M_1,M_2$ compatibles ne sont pas équivalents selon la relation mécaniste retenue, le dossier n'élimine ni l'un ni l'autre. Une règle qui choisit $M_1$ plutôt que $M_2$ doit donc utiliser une information supplémentaire : une contrainte de classe, une préférence, un coût, un prior ou une nouvelle observation. Elle ne peut être déduite de la seule assertion que les deux sont compatibles.\end{proof}

\begin{definition}[Sélection et identification]
Une règle de sélection $s$ choisit $s(D)\in\Inv_{\pi,C}(D)$ lorsqu'un mécanisme compatible existe. L'\emph{identification} est plus forte : elle établit que tous les mécanismes compatibles appartiennent à la même classe d'équivalence pertinente. Sélectionner un candidat n'implique donc pas l'avoir identifié.
\end{definition}

\section{Statuts épistémiques de la reconstruction}
Toute sortie du système doit être étiquetée comme l'un des statuts suivants : (i) donnée directement observée; (ii) mesure dérivée d'une donnée; (iii) hypothèse de modèle; (iv) mécanisme candidat; (v) propriété démontrée dans la classe annoncée; (vi) propriété corroborée sur un domaine d'essai; (vii) propriété réfutée; (viii) question non identifiée. Le statut « mécanisme validé » signifie validation relativement à un protocole et à un domaine déclarés, jamais vérité absolue pour toute entrée imaginable.

\part{II. Échelles, facettes, temps et observation en direct}

\chapter{Échelles de sondes et capacité observationnelle}
\section{Échelle monotone}
\begin{definition}[Échelle d'observation]
Une échelle est un ensemble totalement ordonné $(\Lambda,\preceq)$ et une famille croissante $(\Fcal_\lambda)_{\lambda\in\Lambda}$ telle que $\lambda\preceq\mu\Rightarrow\Fcal_\lambda\subseteq\Fcal_\mu$. Une intensité $\iota:\Lambda\to\R_{\geq0}$ est monotone : $\lambda\preceq\mu\Rightarrow\iota(\lambda)\leq\iota(\mu)$. Si $\Fcal_\lambda$ est finie, la capacité brute vaut $C(\lambda)=|\Fcal_\lambda|$; si les sondes ont des coûts $c(f)\geq0$, le coût de la famille est $c(\Fcal_\lambda)=\sum_{f\in\Fcal_\lambda}c(f)$.
\end{definition}

\begin{remarque}
Le nombre de sondes, leur coût et leur capacité à distinguer des objets sont trois quantités différentes. Une famille plus grande peut contenir des sondes redondantes; une sonde coûteuse peut n'ajouter aucune distinction; une sonde peu coûteuse peut séparer une paire importante. Il n'existe donc pas de score universel « profondeur = nombre de sondes ».
\end{remarque}

\begin{theoreme}[Raffinement avec l'échelle]
Si $\lambda\preceq\mu$, alors la partition de $D$ induite par $I_{p,\Fcal_\mu}$ raffine celle induite par $I_{p,\Fcal_\lambda}$.
\end{theoreme}
\begin{proof}
La propriété suit directement du théorème de raffinement monotone avec $\Fcal_\lambda\subseteq\Fcal_\mu$.
\end{proof}

\begin{proposition}[Borne de stabilisation sur un domaine fini]
Supposons $|D|=n<\infty$. Dans une suite d'échelles croissantes, chaque raffinement strict de la partition augmente son nombre de classes d'au moins une unité. Il peut donc y avoir au plus $n-1$ raffinement(s) strict(s) avant d'atteindre la partition discrète; après cela, aucune sonde supplémentaire ne peut rendre la partition plus fine.
\end{proposition}
\begin{proof}
Une partition de $D$ possède entre $1$ et $n$ classes. Un raffinement strict remplace au moins une classe par au moins deux classes, accroissant le nombre total d'au moins un. Partant d'une seule classe, il ne peut y avoir plus de $n-1$ accroissements stricts avant d'obtenir $n$ classes. La partition discrète est déjà la partition la plus fine d'un ensemble de $n$ objets.
\end{proof}

\section{Capacité réellement discriminante}
\begin{definition}[Taux de séparation empirique]
Pour un jeu fini de paires $\Pi\subseteq\binom{D}{2}$ et une famille $\Gcal$, le taux de séparation est
\[
S_{\Pi}(\Gcal)=\frac{1}{|\Pi|}\sum_{\{x,y\}\in\Pi}\1\{I_{p,\Gcal}(x)\neq I_{p,\Gcal}(y)\},
\]
si $|\Pi|>0$. Cette valeur mesure uniquement le jeu de paires annoncé; elle n'est pas une garantie sur les paires non échantillonnées.
\end{definition}

\begin{proposition}[Monotonie de séparation]
Si $\Gcal\subseteq\Hcal$, alors $S_{\Pi}(\Gcal)\leq S_{\Pi}(\Hcal)$ pour le même jeu de paires et des mesures comparables.
\end{proposition}
\begin{proof}
Pour chaque paire, si la signature de $\Gcal$ diffère, alors la signature étendue de $\Hcal$ diffère aussi. L'indicatrice de séparation ne peut donc que rester égale ou passer de $0$ à $1$. La somme puis la division par $|\Pi|$ préservent cette inégalité.
\end{proof}

\chapter{Fonction de seuil et géométrie des fibres}
\section{Seuil de séparation d'une paire}
\begin{definition}[Seuil informationnel]
Sous une échelle totalement ordonnée, on pose
\[
\Theta_p(\{x,y\})=\sup\{\iota(\lambda): I_{p,\Fcal_\lambda}(x)=I_{p,\Fcal_\lambda}(y)\}.
\]
Le supremum peut être $+\infty$ si aucun niveau ne sépare la paire, et le seuil exact peut dépendre d'une convention de bord si le supremum n'est pas atteint. Le domaine de sondes et les règles d'égalité utilisées à chaque échelle doivent être constants dans cette définition.
\end{definition}

\begin{hypothese}[Chaîne d'échelles et seuils]
La famille des échelles est une chaîne totalement ordonnée; les familles de sondes sont emboîtées; la fonction $\iota$ est monotone. Pour la loi des minima, on suppose aussi que les ensembles de niveaux d'indiscernabilité sont des segments initiaux (ce qui découle de l'emboîtement si les intensités respectent l'ordre). Les conventions pour les cas vides et non bornés sont déclarées.
\end{hypothese}

\begin{theoreme}[Loi du minimum double]
Pour tous $x,y,z\in D$,
\[
\Theta_p(\{x,z\})\geq\min\bigl(\Theta_p(\{x,y\}),\Theta_p(\{y,z\})\bigr).
\]
\end{theoreme}
\begin{proof}
Posons $a=\Theta_p(\{x,y\})$ et $b=\Theta_p(\{y,z\})$. Pour tout $r<\min(a,b)$ non négatif, la définition du supremum et le caractère segment initial de l'indiscernabilité impliquent que l'on peut choisir des échelles où $x,y$ sont indiscernables et où $y,z$ le sont à une intensité au moins $r$; par emboîtement, on descend au niveau commun le moins élevé de ces échelles. À ce niveau, $I(x)=I(y)$ et $I(y)=I(z)$. La transitivité de l'égalité impose $I(x)=I(z)$. Le seuil de $x,z$ est ainsi au moins $r$. En faisant tendre $r$ vers $\min(a,b)$, on obtient l'inégalité. Les cas où le minimum vaut $0$ ou $+\infty$ suivent par les conventions de borne et par le même raisonnement d'ordre.
\end{proof}

\begin{definition}[Ultramétrique dérivée]
Pour $x=y$, posons $d_\Theta(x,x)=0$. Pour $x\neq y$, posons $d_\Theta(x,y)=\exp(-\Theta_p(\{x,y\}))$, avec $e^{-\infty}=0$. Si un seuil négatif était permis par une autre convention, il faudrait décaler ou renormaliser l'échelle avant cette transformation; ici les intensités sont non négatives.
\end{definition}

\begin{theoreme}[La distance de seuil est une pseudoultramétrique]
Sous les hypothèses du théorème de minimum double, $d_\Theta$ est symétrique, non négative, nulle sur la diagonale et satisfait
\[
d_\Theta(x,z)\leq\max\{d_\Theta(x,y),d_\Theta(y,z)\}.
\]
Elle est une métrique si, pour chaque $x\neq y$, le seuil $\Theta_p(\{x,y\})$ est fini.
\end{theoreme}
\begin{proof}
La symétrie découle de l'usage des paires non ordonnées et la diagonalité de la définition. De la loi du minimum double et de la décroissance de $u\mapsto e^{-u}$, on obtient
\[
e^{-\Theta(x,z)}\leq e^{-\min(\Theta(x,y),\Theta(y,z))}=\max(e^{-\Theta(x,y)},e^{-\Theta(y,z)}).
\]
C'est l'inégalité ultramétrique. Dans le cas $x\neq y$, $d_\Theta(x,y)>0$ exactement lorsque $\Theta(x,y)<+\infty$; cette condition donne la séparation nécessaire à une métrique. Si deux objets restent indiscernables à toutes les échelles, la distance vaut zéro entre objets distincts et l'on n'a qu'une pseudoultramétrique.
\end{proof}

\begin{remarque}[Bord critique]
La valeur d'un supremum ne détermine pas si l'indiscernabilité tient exactement au niveau critique. Le squelette de seuil décrit les coupures strictement inférieures et supérieures; une convention fermée aux seuils ou une mesure explicite du bord est nécessaire pour reconstruire l'appartenance à chaque seuil critique. Cette donnée ne peut être déduite du seul nombre $\Theta$.
\end{remarque}

\chapter{Un ordre produit des échelles observationnelles}
\section{Indices indépendants}
\begin{definition}[Indexation multi-échelle]
Soient \((\Lambda,\preceq_\Lambda)\) un poset d'échelles de résolution, \((\mathfrak F,\preceq_F)\) un poset de familles de sondes et \((\Delta,\preceq_\Delta)\) un poset de résolutions temporelles. Une échelle composite est un élément \(s=(\lambda,F,\delta)\) de l'ordre produit \(\Sigma=\Lambda\times\mathfrak F\times\Delta\), muni de l'ordre coordonnée par coordonnée. L'ordre n'affirme une amélioration de résolution que si chaque coordonnée a une sémantique de raffinement déclarée.
\end{definition}
\begin{definition}[Empreinte en tour multi-échelle]
Pour chaque \(s\in\Sigma\), soit \(I_s:\Y\to\Z_s\) l'empreinte à l'échelle \(s\). La famille complète est \(\mathbf I_\Sigma(y)=(I_s(y))_{s\in\Sigma}\). Une empreinte à une échelle instrumentale fixée peut donc demeurer multi-facettes ou multi-temporelle, si les composantes correspondantes sont conservées.
\end{definition}
\begin{proposition}[Raffinement des empreintes par ajout de coordonnées]
Soit \(A\subseteq B\) des indices de sondes et de temps, et supposons que l'empreinte sur \(B\) conserve exactement toutes les composantes de celle sur \(A\) plus éventuellement des composantes supplémentaires. Alors l'indiscernabilité induite par l'empreinte de \(B\) est plus fine ou égale à celle induite par \(A\).
\end{proposition}
\begin{proof}
Si les empreintes de \(B\) sont égales, toutes leurs coordonnées le sont, notamment celles qui constituent l'empreinte de \(A\). Ainsi l'égalité sous \(B\) implique l'égalité sous \(A\), ce qui est la définition du raffinement.\end{proof}
\begin{remarque}
La monotonie ne vaut pas si l'on remplace les anciennes composantes par un résumé différent ou si le protocole de mesure change sans relation de restriction. Il faut alors établir un théorème de comparaison adapté, pas présumer une hiérarchie.
\end{remarque}

\chapter{Trajectoires et signatures temporelles}
\section{Deux objets distincts : état et trajectoire}
\begin{definition}[État et trajectoire d'exécution]
Un état d'exécution est $z(t)\in\Zcal$ à un temps $t$ d'un domaine temporel $\Tcal$, qui peut être $\N$ en temps discret ou un intervalle de $\R_{\geq0}$ en temps continu. Une trajectoire est une fonction $z:\Tcal\to\Zcal$. Elle est déterminée par le système, les entrées exogènes, les états internes, les décisions d'ordonnancement et les éventuelles variables aléatoires. Si ces variables ne sont pas accessibles, la trajectoire complète n'est pas présumée connue.
\end{definition}

\begin{definition}[Réponse temporelle d'une sonde]
Pour un contexte éventuellement variable $p(t)$ et une sonde $f$, la réponse idéale au temps $t$ est
\[
y_f(t)=\Phi(p(t),z(t),f),
\]
quand cette expression est définie. La fonction $t\mapsto y_f(t)$ est le signal idéal visé. Une acquisition effective renvoie généralement un sous-ensemble discret d'échantillons et d'événements; il ne faut pas confondre le signal idéal avec les échantillons effectivement acquis.
\end{definition}

\begin{definition}[Empreinte temporelle idéale]
Pour une fenêtre $J\subseteq\Tcal$ et une famille $\Gcal$, l'empreinte temporelle idéale est
\[
\Ical_{p,J,\Gcal}(z)=\bigl(y_f(t)\bigr)_{(t,f)\in J\times\Gcal}.
\]
Deux trajectoires sont indiscernables sur $J$ relativement à $\Gcal$ si leurs empreintes temporelles idéales sont égales sur tout $J\times\Gcal$. Cette notion est plus forte que l'égalité d'un état à un seul instant.
\end{definition}

\begin{definition}[Empreinte temporelle échantillonnée]
Un plan d'acquisition $S_f\subseteq J$ induit l'empreinte échantillonnée $\Ical^S(z)=(y_f(t))_{f\in\Gcal,t\in S_f}$. L'indiscernabilité échantillonnée se définit par l'égalité de ces valeurs uniquement aux instants effectivement acquis. Elle ne permet aucune conclusion directe sur les intervalles non échantillonnés.
\end{definition}

\section{Événements d'observation et horloges}
\begin{definition}[Enregistrement d'événement]
Un événement d'observation est un tuple
\[
e_i=(\mathrm{id}_i,\mathrm{model}_i,\mathrm{hash}_i,\mathrm{probe}_i,\mathrm{target}_i,
 t_i^{\mathrm{event}},t_i^{\mathrm{capture}},t_i^{\mathrm{arrival}},n_i,v_i,u_i,c_i,s_i),
\]
où $t^{\mathrm{event}}$ désigne le temps annoncé de l'événement source, $t^{\mathrm{capture}}$ le temps d'acquisition, $t^{\mathrm{arrival}}$ l'arrivée au collecteur, $n_i$ un numéro de séquence dans son espace de numérotation, $v_i$ la valeur ou le statut d'échec, $u_i$ les informations d'incertitude, $c_i$ le contexte et $s_i$ la provenance et le statut d'intégrité. Une donnée inconnue reste inconnue; une absence de timestamp ne doit pas être remplacée par un temps inventé.
\end{definition}

\begin{definition}[Flux live]
Un flux live est une suite d'événements reçus au fil de l'exécution, accompagnée d'un protocole de démarrage, pause, reprise et arrêt; d'une politique d'horodatage; d'une politique de tampon; d'une détection des numéros de séquence manquants; d'un contrôle de débit; d'une politique d'erreurs et de backpressure; d'une provenance; et d'une règle de finalisation du flux. Le flux est \emph{continu au sens opérationnel} s'il reste ouvert pendant l'exécution et reçoit des observations successives selon un plan annoncé. Cela ne veut pas dire qu'il mesure mathématiquement chaque instant du temps réel.
\end{definition}

\begin{definition}[Filtration de l'information disponible]
Soit $E_i$ un processus d'événements ordonnés par temps d'arrivée. La filtration $\Fcal_t^{\mathrm{arr}}=\sigma(E_i:t_i^{\mathrm{arrival}}\leq t)$ représente l'information disponible au collecteur à la date d'arrivée $t$. Les décisions prises au temps $t$ sont non anticipatives si elles sont mesurables par rapport à $\Fcal_t^{\mathrm{arr}}$; elles ne peuvent utiliser des événements qui n'étaient pas encore arrivés.
\end{definition}

\begin{proposition}[Non-anticipativité d'un observateur en ligne]
Si une décision $a_t$ est une fonction mesurable uniquement du préfixe d'événements reçus jusqu'à $t$, alors deux exécutions ayant le même préfixe reçu jusqu'à $t$ induisent la même décision à $t$.
\end{proposition}
\begin{proof}
Les deux exécutions fournissent exactement le même argument à la fonction de décision sur le préfixe. Une fonction déterministe appliquée à des arguments identiques produit la même sortie. Pour un décideur randomisé, la conclusion vaut conditionnellement à la même graine ou à la même variable aléatoire interne; sinon la variation est une randomisation du décideur, pas une anticipation de l'avenir.
\end{proof}

\section{Ce qu'on peut et ne peut pas déduire d'un échantillonnage}
\begin{theoreme}[Non-identifiabilité sans hypothèse de régularité]
Pour tout plan d'échantillonnage fini $S=\{t_1,\ldots,t_n\}$ contenant un intervalle non échantillonné, il existe deux fonctions continues distinctes $y_0,y_1$ qui ont les mêmes valeurs sur $S$. Par conséquent, les échantillons seuls ne permettent pas de reconstruire exactement un signal continu arbitraire.
\end{theoreme}
\begin{proof}
Prenons $y_0(t)=0$. Choisissons un intervalle ouvert $(a,b)$ ne contenant aucun $t_i$ et définissons une fonction triangulaire continue $g$ nulle hors de $(a,b)$, nulle aux extrémités et strictement positive à l'intérieur. Posons $y_1=g$. Alors $y_1(t_i)=0=y_0(t_i)$ pour tout échantillon, mais $y_1\neq y_0$ à l'intérieur de l'intervalle. Le même argument s'applique à deux trajectoires dont la sonde mesure ces fonctions. Donc l'exactitude continue ne découle pas d'un nombre fini d'échantillons.
\end{proof}

\begin{theoreme}[Erreur d'interpolation sous une hypothèse de Lipschitz]
Supposons que le signal idéal $y:[0,T]\to\R^d$ est $L$-Lipschitz, c'est-à-dire $\norm{y(t)-y(s)}\leq L|t-s|$. Si chaque instant $t$ est à distance au plus $\Delta/2$ d'un échantillon $t_i$, l'estimateur du plus proche échantillon $\widehat y(t)=y(t_i)$ vérifie
\[
\norm{y(t)-\widehat y(t)}\leq L\Delta/2.
\]
\end{theoreme}
\begin{proof}
Par hypothèse, il existe $t_i$ tel que $|t-t_i|\leq\Delta/2$. L'inégalité de Lipschitz donne directement $\norm{y(t)-y(t_i)}\leq L|t-t_i|\leq L\Delta/2$. Sans borne connue sur $L$ ou sans borne sur le pas d'échantillonnage, cette conclusion ne s'applique pas.
\end{proof}

\begin{proposition}[La détection de pertes ne restaure pas les valeurs perdues]
Si un flux numéroté révèle une lacune entre deux numéros de séquence, le collecteur peut conclure qu'au moins un événement attendu est absent de la séquence reçue, sous l'hypothèse que le producteur attribue les numéros sans réutilisation. Il ne peut pas déduire la valeur manquante à partir du seul numéro.
\end{proposition}
\begin{proof}
Un numéro identifie la place d'un événement dans l'ordre du producteur, mais n'encode pas nécessairement son contenu. Deux flux source peuvent être identiques sur tous les événements reçus et différer uniquement par la valeur de l'événement perdu. Le reçu est identique dans les deux cas; la valeur perdue n'est donc pas identifiable à partir du reçu seul.
\end{proof}

\section{Critères minimaux d'une observation live sérieuse}
Le caractère live exige au minimum : un identifiant de session; l'empreinte de l'artefact et du runtime au démarrage; des identifiants versionnés de sondes; un temps d'événement et un temps de capture distincts lorsque ces notions diffèrent; un numéro de séquence ou une autre preuve d'ordre; la valeur mesurée et son statut d'erreur; des bornes de retard et de fréquence; un détecteur de lacunes; une politique de tampon et de saturation; une politique de pause/reprise/arrêt; des métadonnées de contexte; un hash par événement ou lot; la possibilité de rejouer la séquence; et une règle de fin qui marque les événements attendus mais non reçus.

Une observation à fréquence élevée qui surcharge le modèle peut changer le phénomène observé. Le protocole doit donc mesurer le coût de l'instrumentation et comparer, sur un contrôle, le comportement avec et sans instrumentation. Une capture dite « passive » ne signifie pas « sans effet » tant que cet effet n'a pas été borné ou évalué.

\chapter{Traces live, histoire et reconstruction temporelle}
\section{Trace, événements et états}
\begin{definition}[Trace d'observation]
Une trace discrète est une suite \(\tau=((t_i,e_i))_{i\in J}\) où \(J\) est un ensemble d'indices ordonné, \(t_i\) un temps enregistré et \(e_i\) un événement de mesure, de commande ou de perte. Une trace continue idéale est une fonction de temps, mais une trace réelle est généralement un échantillonnage sous une horloge et un ordonnanceur déclarés.
\end{definition}
\begin{definition}[Dossier temporel complet relatif au protocole]
Un dossier temporel est le tuple constitué des événements bruts disponibles, horodatages, identifiants de sondes, contextes, interventions, versions du modèle, erreurs d'acquisition, règles de synchronisation, données manquantes et métadonnées de provenance. « Complet » signifie complet relativement au protocole et au domaine déclarés, non complet par rapport à toute propriété imaginable du système.
\end{definition}
\begin{proposition}[Un échantillonnage fini ne détermine pas un signal arbitraire]
Pour tout ensemble fini de temps d'échantillonnage distincts \(t_1,\ldots,t_n\), il existe deux fonctions continues distinctes ayant les mêmes valeurs en chacun de ces temps.
\end{proposition}
\begin{proof}
Soit \(g(t)=\prod_{i=1}^n(t-t_i)\). Cette fonction polynomiale est non nulle hors de ses racines et vérifie \(g(t_i)=0\). Pour toute fonction continue \(f\), les fonctions \(f\) et \(f+g\) sont distinctes mais ont les mêmes valeurs aux temps \(t_i\).\end{proof}
\begin{proposition}[Borne d'interpolation sous régularité Lipschitz]
Soit \(f:[a,b]\to\R\) une fonction \(L\)-Lipschitz, c'est-à-dire \(\abs{f(t)-f(s)}\le L\abs{t-s}\). Si les temps d'échantillonnage couvrent l'intervalle avec un maillage maximal \(h\), alors l'estimation par valeur au plus proche échantillon satisfait \(\abs{\hat f(t)-f(t)}\le Lh/2\) lorsque chaque point a un échantillon à distance au plus \(h/2\); avec la seule définition d'un écart maximal entre échantillons \(h\), la borne générale sûre est \(Lh\).
\end{proposition}
\begin{proof}
Soit \(t_i\) un échantillon le plus proche. Par Lipschitz, \(\abs{f(t)-f(t_i)}\le L\abs{t-t_i}\). La borne découle de la distance garantie par le maillage. Le facteur \(1/2\) nécessite la couverture au plus proche dans un maillage borné aux deux extrémités; sans cette garantie, la distance maximale ne dépasse que \(h\).\end{proof}
\begin{remarque}
Le résultat ne reconstruit pas exactement le signal. Il transforme une hypothèse de régularité en une borne d'erreur. Si la régularité n'est pas justifiée, la garantie n'est pas disponible.
\end{remarque}

\chapter{Stabilité, non-stationnarité et estimation en ligne}
\section{État courant et fenêtres d'analyse}
\begin{definition}[Empreinte fenêtre]
Pour un flux $y_f(t)$, une fenêtre $W_k=[t_k,t_k+w)$ et une statistique $\psi$ définie sur des observations disponibles, l'empreinte de fenêtre est $J_k=\psi(\{(t_i,y_i):t_i\in W_k\})$. La valeur $J_k$ décrit cette fenêtre et ne remplace pas le flux brut; la granularité et la définition de $\psi$ doivent être conservées.
\end{definition}

\begin{definition}[Dérive observationnelle]
Pour une métrique $d$ fixée sur les empreintes de fenêtres, une dérive au seuil $\delta$ entre fenêtres $k$ et $k+1$ est l'événement $d(J_k,J_{k+1})>\delta$. Ce fait détecte un changement relativement à la métrique et au seuil spécifiés; il n'identifie pas automatiquement sa cause.
\end{definition}

\begin{proposition}[La dérive dépend de la sonde]
Il existe deux trajectoires distinctes dont une sonde $f$ donne des empreintes stables, tandis qu'une autre sonde $g$ révèle une dérive.
\end{proposition}
\begin{proof}
Considérons des états $z_0,z_1$ tels que $\Phi(p,z_0,f)=\Phi(p,z_1,f)$ mais $\Phi(p,z_0,g)\neq\Phi(p,z_1,g)$. Une trajectoire passant de $z_0$ à $z_1$ apparaît constante sous $f$ et change sous $g$. L'existence de tels états est exactement la possibilité de non-injectivité de la première sonde; le théorème de factorisation montre qu'une propriété qui varie dans une fibre ne peut être déterminée par cette sonde.
\end{proof}

\begin{remarque}
La stabilité d'un résumé n'établit pas la stabilité du modèle entier. Elle établit la stabilité de ce résumé, sur les fenêtres examinées, pour les sondes utilisées. Une absence de dérive observée peut signifier que le phénomène est stable, que le pas d'échantillonnage est trop large, que la sonde est insensible, ou que les événements sont manquants; ces possibilités doivent être distinguées expérimentalement.
\end{remarque}

\chapter{Empreinte de trajectoire et prédiction continue}
\section{De la trace observée à la loi de transition}
\begin{definition}[Dossier de trajectoires]
Un dossier temporel $D_{\mathrm{traj}}$ contient des séquences d'événements ordonnées avec temps d'événement et temps d'arrivée, identifiants de session, état de séquence, entrées, interventions, valeurs mesurées, provenance et marqueurs d'événements manquants. Les deux horloges sont distinguées lorsqu'elles peuvent diverger. Si seul l'ordre est fiable, l'horodatage est traité comme intervalle ou ordre partiel, pas comme temps exact.
\end{definition}

\begin{definition}[Compatibilité temporelle]
Un mécanisme $M$ est temporellement compatible avec $D_{\mathrm{traj}}$ si ses prédictions d'événements, de transitions et de sorties satisfont le critère de compatibilité préenregistré pour chaque trajectoire valide, en tenant compte des événements manquants selon le modèle d'observation déclaré. Un algorithme ne peut considérer un événement absent comme une sortie nulle sans une règle qui établit explicitement cette sémantique.
\end{definition}

\begin{proposition}[Un échantillonnage irrégulier non modélisé peut confondre deux dynamiques]
Si le dossier n'enregistre que les valeurs échantillonnées et n'enregistre ni les instants réels d'échantillonnage ni les événements manquants, alors deux dynamiques distinctes peuvent induire le même dossier observé.
\end{proposition}
\begin{proof}
Considérons deux fonctions temporelles qui coïncident à tous les instants consignés mais diffèrent strictement entre deux instants consignés. Si l'observateur ne stocke que les valeurs consignées, les deux fonctions donnent le même dossier. Leur comportement entre observations est néanmoins différent. Le dossier ne permet donc pas d'identifier la dynamique continue sans hypothèse de régularité ou mesure supplémentaire.\end{proof}

\section{Prédiction des transitions et tests de rejeu}
\begin{definition}[Rejeu déterministe]
Un rejeu est déterministe relativement à une exécution si le paquet enregistre toutes les entrées, interventions, versions, paramètres et sources d'aléa nécessaires pour reproduire la même trace, et si le mécanisme d'exécution fournit le déterminisme annoncé. Si une dépendance ou un état externe manque, le rejeu n'est que partiel et son statut doit le signaler.
\end{definition}

\begin{proposition}[Identité des seules entrées non suffisante pour un rejeu bit-à-bit]
Deux exécutions ayant la même séquence visible d'entrées peuvent produire des sorties différentes si leur état initial, leurs paramètres stochastiques, leur environnement numérique ou leurs interventions diffèrent.
\end{proposition}
\begin{proof}
Par la définition de mécanisme, les sorties dépendent de l'état et des actions autant que de l'entrée. Des états initiaux distincts peuvent produire des sorties distinctes à entrée identique. De même, un mécanisme stochastique peut tirer des réalisations distinctes. Les seules entrées ne déterminent donc pas nécessairement la sortie; il faut enregistrer les autres variables pertinentes au niveau de reproductibilité revendiqué.\end{proof}

\chapter{Mise à jour en ligne de l'ensemble inverse}
\section{Assimilation d'un nouvel événement}
\begin{definition}[Mise à jour incrémentale]
Soit $\mathcal V_t$ l'ensemble de candidats compatibles avec le dossier à l'instant logique $t$. Pour un nouvel événement $e_{t+1}$, la mise à jour exacte est
\[
\mathcal V_{t+1}=\{M\in\mathcal V_t:C(M,e_{t+1},\pi)=\mathrm{vrai}\}.
\]
La définition suppose que la règle de compatibilité est stable et que l'événement est validé selon les contrôles d'intégrité. Si l'événement révise une donnée antérieure, le système doit recalculer les contraintes affectées plutôt que d'appliquer aveuglément une intersection monotone.
\end{definition}

\begin{proposition}[Monotonie de la mise à jour par ajout d'évidence]
Si un nouvel événement ajoute seulement une contrainte et n'invalide ni ne modifie les observations précédentes, alors $\mathcal V_{t+1}\subseteq\mathcal V_t$.
\end{proposition}
\begin{proof}
La définition de $\mathcal V_{t+1}$ exige à la fois l'appartenance à $\mathcal V_t$ et la satisfaction de la nouvelle contrainte. Il s'agit donc d'un sous-ensemble de $\mathcal V_t$.\end{proof}

\begin{proposition}[Une mise à jour vide est une alerte, pas une raison de forcer un candidat]
Si $\mathcal V_{t+1}=\varnothing$, aucun candidat de la classe courante ne satisfait simultanément le dossier et le nouveau critère. Il faut rapporter cette incompatibilité et examiner données, classe, protocole et hypothèses; retourner silencieusement un modèle non compatible ne respecte pas la définition d'un candidat reconstruit.
\end{proposition}
\begin{proof}
L'égalité à l'ensemble vide signifie qu'il n'existe aucun $M$ de la classe courante qui satisfasse toutes les contraintes requises. Retourner un tel $M$ tout en le déclarant compatible contredirait la définition de l'ensemble inverse. Un modèle peut être proposé comme approximation ou comme nouvelle hypothèse, mais son statut doit changer et l'incompatibilité doit rester attachée au registre.\end{proof}

\part{III. Sondes, vues multiples et interventions}

\chapter{Facettes observationnelles complémentaires}
\section{Les facettes sont des modalités de mesure}
On indexe les familles de sondes par un ensemble $J$ : $\Fcal=\bigcup_{j\in J}\Fcal^{(j)}$. L'index $j$ nomme un domaine de mesure, et non une interprétation. Les facettes ci-dessous sont des familles typiques; elles ne sont ni universelles, ni nécessairement indépendantes.

\begin{longtable}{p{0.20\textwidth}p{0.38\textwidth}p{0.32\textwidth}}
\toprule
\textbf{Facette} & \textbf{Exemples de sondes mesurables} & \textbf{Question observable}\\
\midrule
\endhead
Comportement externe & sorties, logits, probabilités, distributions de génération, latence, erreurs & Quelles réponses le modèle produit-il sous des entrées déclarées ?\\
Structure statique & modules, paramètres, formes, tokenizer, normes, connexions, config & Quels composants existent et quels chemins d'exécution peuvent être instrumentés ?\\
Représentations & activations, états cachés, tenseurs intermédiaires, distances entre états & Quelles variations internes coïncident avec les variations d'entrée et de sortie ?\\
Attention / routage & matrices d'attention, scores de routage ou portes, si l'architecture les expose & Comment les poids d'allocation calculés varient-ils ? Ces valeurs ne sont pas automatiquement des explications causales.\\
Sensibilité locale & gradients, jacobiens partiels, différences finies, stabilité numérique & À quelles perturbations locales la sortie est-elle sensible dans le régime mesuré ?\\
Intervention & ablation, remplacement, patching, permutation, mise à l'échelle avec contrôle & Que change une modification définie sous un contexte apparié ?\\
Dynamique & états et sorties répétés, trajectoires autoregressives, caches, états récurrents, événements live & Comment évoluent les états et réponses dans le temps d'exécution ?\\
Stochasticité & répétitions par graines, distributions de sorties, entropie, variance & Quelle variabilité vient du modèle, de l'échantillonnage ou du runtime ?\\
Robustesse / OOD & perturbations, permutations, domaines exclus, changements de distribution & Quelles régularités survivent hors du domaine de découverte ?\\
Opérationnelle & temps, mémoire, appels, exceptions, saturation, coût des hooks & À quel coût et avec quel effet de perturbation l'observation est-elle acquise ?\\
Provenance & hash de modèle, poids, prompt, sonde, configuration et environnement & Peut-on relier un résultat à l'artefact et au protocole réellement exécutés ?\\
\bottomrule
\end{longtable}

\section{Complémentarité des vues}
\begin{definition}[Vue et empreinte multi-facettes]
Une vue $j$ est déterminée par une famille $\Fcal^{(j)}$, un protocole de mesure et une fenêtre d'exécution. Pour une famille $J_0\subseteq J$, l'empreinte multi-facettes est la concaténation canonique des signatures admissibles des vues :
\[
I_{J_0}^{p}(x)=\bigl(I_{p,\Fcal^{(j)}}(x)\bigr)_{j\in J_0}.
\]
Les valeurs indisponibles sont marquées par leur statut; elles ne sont pas remplacées par des zéros qui les rendraient comparables à une mesure réelle.
\end{definition}

\begin{theoreme}[Complémentarité stricte d'une nouvelle vue]
Ajouter une vue $j$ ne fusionne jamais les classes d'indiscernabilité. L'ajout crée une distinction nouvelle exactement lorsqu'il existe deux objets indiscernables pour toutes les vues précédentes mais que la nouvelle vue sépare.
\end{theoreme}
\begin{proof}
Appliquer le théorème de raffinement monotone aux unions de familles de sondes. Le critère de stricteté suit du corollaire de séparation d'une nouvelle sonde en considérant la signature complète de la nouvelle vue comme un tuple de coordonnées.
\end{proof}

\begin{definition}[Redondance empirique d'une vue]
Sur un ensemble fini $D_{\mathrm{obs}}$, une vue $j$ est redondante relativement à $J_0$ si sa signature est constante sur chaque fibre empirique de $I_{J_0}$. Cette propriété ne concerne que $D_{\mathrm{obs}}$; la redondance universelle exige la même constance sur tout le domaine $D$.
\end{definition}

\begin{remarque}[Pourquoi le nombre de vues ne suffit pas]
Deux vues peuvent être fortement redondantes. Une vue supplémentaire n'est utile que par les distinctions ou tests qu'elle apporte, compte tenu de son coût, de ses incertitudes et des hypothèses scientifiques en compétition. Le but n'est pas de maximiser le nombre de tableaux, mais de maximiser l'information discriminante pertinente.
\end{remarque}

\chapter{Sélection active de sondes et recherche discriminante}
\section{Hypothèses concurrentes}
Soit $\Hcal=\{H_1,\ldots,H_m\}$ une famille d'hypothèses explicites. Chaque hypothèse fournit une prédiction sur la distribution de réponse d'une sonde $a$ dans un contexte donné. L'action $a$ peut être une observation passive ou une intervention active, à condition que ses préconditions et son coût soient spécifiés.

\begin{definition}[Gain d'information bayésien]
Si une distribution a priori $\pi(H)$ est explicitement justifiée et que les vraisemblances $p(y\mid H,a,D)$ sont définies, le gain d'information attendu de l'action $a$ est
\[
\operatorname{IG}(a\mid D)=H(\pi(H\mid D))-\Exp_{y\mid a,D}\bigl[H(\pi(H\mid D,y,a))\bigr],
\]
où $H$ est l'entropie de Shannon. Si aucun a priori ni modèle de vraisemblance n'est défendable, ce score n'est pas calculé comme un fait; on peut utiliser une mesure de séparation minimax ou publier les prédictions sans agréger.
\end{definition}

\begin{definition}[Action discriminante sous contrainte]
Pour un budget $B$ et un coût explicite $c(a)$, une action admissible maximise une fonction déclarée, par exemple $\operatorname{IG}(a\mid D)/c(a)$ sous $c(a)\leq B$, ou maximise la séparation minimale entre prédictions des hypothèses. Le choix est enregistré avant l'obtention du résultat pour éviter de réécrire a posteriori la motivation du test.
\end{definition}

\begin{proposition}[L'observation confirmatoire seule est insuffisante pour éliminer des hypothèses équivalentes]
Si deux hypothèses prédisent la même loi de réponse pour toutes les sondes exécutées, les données de ces sondes ne les distinguent pas, quelle que soit leur quantité, si elles ne couvrent que cette même famille de sondes.
\end{proposition}
\begin{proof}
Si $p(y\mid H_1,a)=p(y\mid H_2,a)$ pour chaque action effectivement exécutée, la vraisemblance de toute suite de résultats générée par ces actions est identique sous les deux hypothèses, sous le modèle probabiliste spécifié. Le rapport de vraisemblance reste donc égal à $1$ et les données de ces actions n'apportent aucun pouvoir discriminant entre elles. Une autre sonde ou intervention dont les prédictions divergent est nécessaire.
\end{proof}

\chapter{Interventions, contrôles et effets mesurés}
\section{Définition d'une intervention}
\begin{definition}[Intervention opérationnelle]
Une intervention $a$ est un opérateur versionné qui modifie une variable interne ou une partie d'un calcul sous une sémantique explicite. Sa fiche comprend au minimum : identifiant; cible exacte; emplacement dans le graphe de calcul; tenseur ou état modifié; opération mathématique; portée (exemple, token, séquence, batch ou couche); paramètres; contrôles; procédure de restauration; hash du modèle avant et après; contexte; conditions de validité; et méthode de mesure du résultat.
\end{definition}

\begin{remarque}[L'étiquette seule ne définit pas l'intervention]
Les mots « ablation », « inhibition », « patching » ou « scale » ne suffisent pas. Par exemple, mettre à zéro la sortie d'un bloc Transformer peut aussi annuler une connexion résiduelle selon le point où le hook est posé. Une expérience qui prétend retirer uniquement la transformation doit définir une opération qui préserve les chemins qui ne sont pas la cible, ou reconnaître qu'elle teste une intervention plus large.
\end{remarque}

\section{Effet déterministe apparié}
\begin{definition}[Effet sous contexte fixé]
Pour une entrée $u$, un contexte fixe $p$ et une intervention $a$, soit $Y_0(u,p)$ la sortie de référence et $Y_a(u,p)$ la sortie après l'intervention. Pour une métrique de sortie $d_Y$, l'effet local est $\Delta_a(u,p)=d_Y(Y_0(u,p),Y_a(u,p))$. Le vecteur différence $Y_0-Y_a$ peut aussi être publié si le codomaine le permet. Le résultat est local à l'entrée, au contexte, à la cible et à l'opération exécutée.
\end{definition}

\begin{theoreme}[Ce que démontre un effet ponctuel]
Si $\Delta_a(u,p)>0$, il est établi que, pour l'exécution mesurée et la métrique déclarée, la sortie de référence diffère de la sortie après intervention. Ce seul fait ne permet pas de conclure que la cible est nécessaire à toute sortie, suffisante pour produire une fonction, ni qu'un mécanisme sémantique complet a été identifié.
\end{theoreme}
\begin{proof}
La première phrase découle de la définition de $\Delta_a$ et de l'inégalité observée. Pour démontrer la non-implication des conclusions plus fortes, considérons deux systèmes ayant la même différence pour cette entrée : dans le premier, la cible intervient dans le mécanisme supposé; dans le second, l'intervention modifie un calcul auxiliaire ou une distribution numérique sans jouer le rôle sémantique allégué. Comme une seule paire de sorties est compatible avec les deux systèmes, le résultat ponctuel ne permet pas de les distinguer. Il faut des contrôles et des contextes supplémentaires pour éliminer le second modèle explicatif.
\end{proof}

\section{Effet moyen pour un modèle stochastique}
\begin{hypothese}[Expérience répétée]
Pour estimer un effet moyen, les contextes sont un échantillon indépendant $U_i\sim\mu$ ou proviennent d'un plan d'échantillonnage connu dont les poids sont utilisés. Les répétitions, distributions de graines et conditions de remise à l'état initial sont déclarées. L'espérance de $g(Y(a,U))-g(Y(0,U))$ doit exister; les observations doivent être assez indépendantes, ou leur dépendance doit être modélisée, pour utiliser l'intervalle de confiance annoncé.
\end{hypothese}

\begin{definition}[Effet moyen sous deux régimes]
Pour une population de contextes $U\sim\mu$ et des sorties potentielles $Y(0,U)$ et $Y(a,U)$, l'effet moyen d'une statistique vectorielle $g$ est
\[
\tau_a=\Exp_{U\sim\mu}[g(Y(a,U))-g(Y(0,U))],
\]
si l'espérance existe. Pour une sortie vectorielle, $g$ peut être l'identité ou une métrique choisie à l'avance. Un estimateur fondé sur un échantillon fini possède une incertitude qui doit être rapportée.
\end{definition}

\begin{proposition}[Appariement des contextes]
Si, pour chaque contexte observé $u_i$, les sorties baseline et intervention sont mesurées sous le même contexte contrôlé, la moyenne des différences $g(Y(a,u_i))-g(Y(0,u_i))$ estime l'effet moyen sur la distribution d'échantillonnage des $u_i$ sous le plan d'échantillonnage déclaré. Elle n'autorise pas une généralisation à une autre distribution $\mu'$ sans hypothèse de transport.
\end{proposition}
\begin{proof}
La différence appariée est calculée pour chaque élément du même échantillon. Sa moyenne empirique est l'estimateur naturel de l'espérance des différences pour la loi qui a généré les $u_i$; sous échantillonnage aléatoire et intégrabilité, sa valeur attendue égale $\tau_a$ pour cette loi. Si la loi cible change, l'espérance peut changer; aucun terme de l'échantillon n'établit l'égalité avec $\mu'$ sans condition supplémentaire.
\end{proof}

\section{Conditions de validité causale}
Une conclusion causale plus forte exige que la cible d'intervention soit effectivement modifiée; que le reste du contexte soit contrôlé ou que les variations soient modélisées; que le système soit restauré entre interventions si elles peuvent persister; que les sorties soient comparables; que les effets de l'instrumentation soient bornés; que les contrôles négatifs et positifs soient pertinents; et que la conclusion soit répliquée dans la population visée. Ces conditions sont logiques et expérimentales, pas une formule qui garantit la réussite.

\begin{theoreme}[Non-identifiabilité mécaniste par une famille finie de sondes]
Soit $D_{\mathrm{test}}\subsetneq D$ un domaine de contextes et une famille de sondes finie observée uniquement sur $D_{\mathrm{test}}$. Sans hypothèse de régularité ou de structure du modèle, il existe deux systèmes distincts qui ont les mêmes résultats sur $D_{\mathrm{test}}$ et diffèrent sur au moins un contexte de $D\setminus D_{\mathrm{test}}$. Ils sont indiscernables sur les données collectées, mais pas équivalents sur tout $D$.
\end{theoreme}
\begin{proof}
Choisissons un contexte $u^*\in D\setminus D_{\mathrm{test}}$. Construisons deux fonctions de réponse $F_0,F_1$ telles que $F_0(u)=F_1(u)$ pour tout $u\in D_{\mathrm{test}}$, mais $F_0(u^*)\neq F_1(u^*)$. Elles induisent les mêmes observations sur le domaine testé et des réponses différentes en $u^*$. Puisque l'étude ne pose aucune hypothèse qui lie les réponses en $u^*$ aux valeurs testées, elle ne peut déterminer laquelle décrit le système sur tout $D$. Le même argument vaut pour des mécanismes internes distincts ayant le même comportement sur les sondes exécutées.
 Une différence mesurée ne prouve pas, à elle seule, un mécanisme identifié.
\end{proof}



\chapter{Conditions d'identifiabilité des composants invariant et variant}
\section{Plan d'expériences apparié}
\begin{definition}[Plan apparié]
Un plan est apparié lorsque les contextes ou interventions comparés partagent les variables contrôlées nécessaires et que les différences restantes sont enregistrées. Le mot « apparié » décrit le protocole; la validité causale exige en plus des hypothèses sur l'assignation, les confondeurs, les effets de report et la mesure.
\end{definition}
\begin{proposition}[Une différence observée n'identifie pas sans condition l'effet de la sonde]
Si deux enregistrements diffèrent à la fois par la sonde et par un contexte non contrôlé, leur différence ne peut être attribuée de façon unique à la sonde à partir de ces deux enregistrements seuls.
\end{proposition}
\begin{proof}
Considérons le modèle additif \(Y=\beta_F F+\beta_C C\). Deux observations où \(F\) et \(C\) changent simultanément ne fournissent qu'une équation pour les deux contributions; plusieurs couples \((\beta_F,\beta_C)\) peuvent produire la même différence. L'attribution unique est donc impossible sans variation séparée ou contrainte additionnelle.\end{proof}
\begin{proposition}[Identifiabilité dans un modèle linéaire]
Pour le modèle \(y=X\theta\) avec matrice connue \(X\), le paramètre \(\theta\) est unique pour toute donnée compatible si et seulement si \(X\) est de rang colonne plein.
\end{proposition}
\begin{proof}
Si \(X\) est de rang colonne plein, son noyau est réduit à \(\{0\}\); si \(X\theta=X\theta'\), alors \(X(\theta-\theta')=0\), donc \(\theta=\theta'\). Si le rang n'est pas plein, il existe \(u\ne0\) tel que \(Xu=0\), et \(\theta\) comme \(\theta+u\) produisent la même donnée; l'unicité échoue.\end{proof}
\begin{remarque}
Ce résultat ne démontre pas que la décomposition d'un LLM est linéaire. Il précise seulement le test d'identifiabilité dans la classe de modèles linéaires choisie. Les modèles non linéaires nécessitent leur propre analyse.
\end{remarque}

\chapter{Interventions et reconstruction causale}
\section{Sémantique d'une intervention}
\begin{definition}[Intervention structurelle]
Pour un mécanisme factorisé en équations $z_j:=f_j(z_{\mathrm{pa}(j)},u)$, une intervention $\doop(z_k:=v)$ remplace l'équation de la variable $z_k$ par $z_k:=v$ dans le modèle intervenu, tout en conservant les autres équations sauf si la définition de l'intervention précise explicitement une modification plus large. La cible, la portée, la durée, le point d'application et les effets de bord sont enregistrés.
\end{definition}

\begin{definition}[Effet interventionnel apparié]
Pour une métrique de sortie $d$, un contexte $u$ et une intervention $a$, l'effet apparié est $\Delta_a(u)=d(Y_M(u,a),Y_M(u,a_0))$, où $a_0$ est le contrôle de référence. Pour les systèmes stochastiques, on doit comparer une quantité spécifiée de leurs distributions (par exemple distance totale variation, différence d'espérance sous intégrabilité, ou différence de probabilité d'un événement). La valeur dépend du contexte et du protocole, et ne doit pas être présentée comme une propriété universelle du composant.
\end{definition}

\begin{theoreme}[Effet mesuré et identification d'une arête causale sont distincts]
Un delta de sortie non nul après une intervention valide sur $z_k$ établit une différence de sortie sous cette intervention dans le contexte testé. Il n'établit pas, à lui seul, que $z_k$ est nécessaire ou suffisant à un mécanisme sur tout le domaine, ni que l'arête structurelle proposée est unique.
\end{theoreme}
\begin{proof}
Le delta est défini pour une intervention et un contexte donnés; il n'énonce donc qu'une différence pour cette paire de conditions. La nécessité universelle requerrait que la sortie perde la propriété cible sous toutes les conditions pertinentes quand le composant est retiré. La suffisance requerrait des tests séparés où l'activation de la cible produit l'effet prévu sous des contrôles appropriés. L'unicité de l'arête nécessiterait en outre d'exclure des structures concurrentes produisant les mêmes réponses aux interventions observées. Rien de cela ne suit du seul delta non nul.\end{proof}

\begin{proposition}[Contrôles négatifs et contrôles positifs répondent à des questions différentes]
Un contrôle négatif teste si la procédure d'intervention produit un effet sans l'action ciblée; un contrôle positif teste si le dispositif et la métrique peuvent détecter un effet connu ou attendu dans les conditions d'essai. La réussite d'un contrôle ne remplace pas l'autre.
\end{proposition}
\begin{proof}
Les deux contrôles portent sur des hypothèses différentes. Le contrôle négatif compare l'effet au cas où l'action pertinente est absente, alors que le contrôle positif vérifie la sensibilité de la procédure à un effet pertinent. Il existe un dispositif qui déclenche toujours un signal (faux positifs sur le contrôle négatif) mais détecte un signal, et un dispositif qui n'émet jamais de signal (pas de faux positifs) mais échoue au contrôle positif. Les deux résultats ne sont donc pas logiquement équivalents.\end{proof}

\section{Interventions réalisables dans un LLM}
Les interventions possibles dépendent de l'interface du modèle. Sur un LLM, elles peuvent porter sur des activations, des têtes d'attention, des résidus, des MLP, des embeddings, des poids ou des couches, uniquement si le framework et l'architecture exposent ces éléments. Un adapter qui ne peut modifier qu'une sortie textuelle ne doit pas simuler une intervention interne. Toute capacité non supportée est déclarée indisponible pour ce modèle et ce protocole; l'étude peut proposer un nouvel adapter, mais ne peut inférer les effets d'une intervention qui n'a pas été exécutée.

\begin{proposition}[Transportabilité limitée des effets]
Un effet interventionnel mesuré dans un contexte $p$ ne se transporte vers un contexte $p'$ que si une hypothèse de stabilité ou une série d'expériences de transport justifie l'égalité ou la borne revendiquée entre les réponses sous intervention dans ces deux contextes.
\end{proposition}
\begin{proof}
Par définition, un effet est une fonction du contexte et de l'intervention. Sans hypothèse reliant ces deux contextes, les valeurs de cette fonction peuvent être arbitrairement différentes hors du point observé. Un seul point ne détermine pas une fonction générale.\end{proof}

\part{IV. Du problème inverse aux mécanismes exécutables}

\chapter{Définition d'un mécanisme exécutable}
\section{Mécanisme comme transducteur}
\begin{definition}[Mécanisme déterministe]
Un mécanisme déterministe est un tuple $M=(U,S,Y,s_0,T,O)$ où $U$ est un espace d'entrées, $S$ un espace d'états internes, $Y$ un espace de sorties, $s_0$ un état initial déclaré, $T:S\times U\to S$ une transition et $O:S\times U\to Y$ une fonction de sortie. Pour une entrée temporelle $u_1,\ldots,u_n$, la trajectoire est définie par $s_{t+1}=T(s_t,u_t)$ et $y_t=O(s_t,u_t)$. Un mécanisme stochastique remplace $T,O$ par des noyaux probabilistes spécifiés.
\end{definition}

\begin{definition}[Contrat d'interface]
Le contrat d'un mécanisme spécifie son domaine d'entrées, les invariants d'état, son domaine de sortie, les préconditions, les postconditions, les ressources, les exceptions et l'incertitude. Une composition n'est admissible que si les sorties du parent immédiat satisfont les préconditions du mécanisme suivant.
\end{definition}

\begin{proposition}[Composition fonctionnelle]
Si $M_1:U\to V$ et $M_2:V\to Y$ sont des fonctions totales sur leurs domaines et si $M_1(U)\subseteq\operatorname{dom}(M_2)$, alors $M_2\circ M_1:U\to Y$ est une fonction bien définie. Pour une chaîne finie de mécanismes, la composition est associative.
\end{proposition}
\begin{proof}
La condition d'inclusion garantit que $M_2(M_1(u))$ est défini pour tout $u\in U$. Pour trois fonctions compatibles, $(M_3\circ M_2)\circ M_1(u)=M_3(M_2(M_1(u)))=M_3\circ(M_2\circ M_1)(u)$; l'égalité point par point prouve l'associativité. Cette propriété est fonctionnelle; elle ne prouve pas la validité scientifique des contrats ni des mécanismes choisis.
\end{proof}

\section{Abstraction hiérarchique}
\begin{definition}[Niveau d'abstraction]
Un niveau est un ordre partiel sur des représentations ou mécanismes, défini par des relations d'implémentation et d'interface déclarées. Dire qu'un parent est à un niveau supérieur signifie qu'il expose une interface plus abstraite ou regroupe des composants, pas qu'il est plus vrai, plus intelligent ou mieux validé.
\end{definition}

\begin{theoreme}[Un objet composé peut devenir un objet d'observation]
Si un mécanisme composé $M^*$ dispose d'une interface déterministe ou probabiliste définie et d'un protocole d'observation applicable, il peut être traité comme nouvel objet $x^*$ dans une seconde phase d'observation. La construction de $M^*$ ne garantit pas que son empreinte soit égale à celle du système source; cette égalité constitue une proposition expérimentale distincte.
\end{theoreme}
\begin{proof}
Par définition, le domaine $\X$ accepte tout objet pour lequel les sondes appliquées sont définies. Si $M^*$ possède une telle interface, il est admissible à ce domaine et l'application $\Phi$ peut lui être appliquée. Aucune égalité avec le modèle source n'intervient dans cette construction; celle-ci doit être testée séparément.
\end{proof}

\begin{proposition}[Critère d'équivalence comportementale relative]
Deux mécanismes $M_1,M_2$ sont équivalents relativement à un protocole $\Pi$ si leurs empreintes sont égales pour chaque contexte et sonde admissibles dans $\Pi$. Cette équivalence ne permet pas de conclure à leur identité structurelle si $\Pi$ n'est pas séparant.
\end{proposition}
\begin{proof}
La première partie est la définition d'équivalence relative. La seconde découle du fait que deux objets distincts peuvent appartenir à la même fibre de l'empreinte; la famille de sondes non injective ne peut distinguer les éléments d'une fibre.
\end{proof}

\chapter{Sémantique exécutable d'un mécanisme}
\section{Transducteur mécaniste}
\begin{definition}[Mécanisme discret déterministe]
Un mécanisme discret déterministe est un tuple
\[
M=(Z,z_0,U,A,Y,\delta,\omega,\Gamma,\mathcal C),
\]
où $Z$ est l'espace d'état, $z_0$ l'état initial, $U$ l'espace des entrées, $A$ l'espace des actions ou interventions, $Y$ l'espace des sorties, $\delta:Z\times U\times A\to Z$ la fonction de transition, $\omega:Z\times U\times A\to Y$ la fonction d'émission, $\Gamma$ le graphe de composants et de dépendances, et $\mathcal C$ les contrats des interfaces. Les fonctions doivent être totales sur le domaine annoncé ou déclarer explicitement où elles sont partielles.
\end{definition}

Pour une séquence d'entrées et d'actions $h_t=((u_0,a_0),\ldots,(u_{t-1},a_{t-1}))$, l'état est défini récursivement par $z_{t+1}=\delta(z_t,u_t,a_t)$ et la sortie par $y_t=\omega(z_t,u_t,a_t)$. Cette définition rend l'exécution indépendante d'un modèle source appelé en cachette : l'exécuteur doit calculer $\delta$ et $\omega$ à partir du mécanisme livré.

\begin{definition}[Mécanisme stochastique]
Un mécanisme stochastique spécifie un noyau conjoint $L(\cdot\mid z,u,a)$ sur $Y\times Z$ pour le couple (sortie, prochain état). Ses marges définissent les lois d'émission et de transition, mais le noyau conjoint est conservé afin de ne pas supposer silencieusement l'indépendance conditionnelle entre la sortie et le prochain état. La loi de trajectoire est obtenue par itération de $L$ à partir de $z_0$, lorsque les espaces sont mesurables, $L$ est un noyau mesurable et les entrées-actions sont admissibles. Toute garantie probabiliste suppose ce cadre; hors de celui-ci, aucune loi de trajectoire ne doit être invoquée implicitement.
\end{definition}

\begin{definition}[Exécution autonome]
Une implémentation $\widehat M$ est autonome relativement au modèle source si, après son chargement, ses sorties sur une entrée ne dépendent que de son propre état, de l'entrée, des interventions autorisées et de ses dépendances déclarées. Un appel caché au modèle source, à un service distant ou à un oracle invalide le statut autonome, sauf si cette dépendance fait explicitement partie du modèle évalué.
\end{definition}

\begin{proposition}[Test nécessaire d'autonomie par fermeture des dépendances]
Si l'exécuteur de $\widehat M$ dépend d'une fonction externe $g$ non incluse dans le paquet déclaré et si la sortie peut changer lorsque $g$ change à entrée et état de $\widehat M$ identiques, alors $\widehat M$ n'est pas autonome relativement à ce paquet.
\end{proposition}
\begin{proof}
L'autonomie exige que les variables déclarées de l'exécuteur déterminent la sortie. Par hypothèse, deux environnements qui ne diffèrent que par $g$ produisent des sorties différentes. La sortie n'est donc pas fonction des seuls états et dépendances inclus dans le paquet. La condition d'autonomie échoue.\end{proof}

\section{Mécanisme, statistique et capture : trois objets distincts}
\begin{definition}[Résumé statistique]
Un résumé $T(D)$ est une application du dossier de données vers un espace de statistiques. Il peut conserver des moments, histogrammes, corrélations ou fréquences tout en éliminant l'ordre temporel, l'identité des unités ou la condition d'intervention. Il n'est pas un mécanisme par sa seule forme mathématique.
\end{definition}

\begin{proposition}[Une même statistique peut correspondre à des dynamiques différentes]
Il existe deux systèmes dont la moyenne empirique de sortie sur un jeu d'essai fini est identique, alors que leurs sorties sur une séquence ordonnée différente ne sont pas identiques.
\end{proposition}
\begin{proof}
Considérons deux entrées binaires $u_0,u_1$. Le système $S_1$ émet $y_t=u_t$. Le système $S_2$ mémorise la première entrée et émet $y_0=u_0$, $y_1=u_0$. Sur l'essai $(u_0,u_1)=(0,0)$, les sorties sont $(0,0)$ pour les deux systèmes, donc leur moyenne vaut zéro. Sur l'essai $(0,1)$, $S_1$ émet $(0,1)$ tandis que $S_2$ émet $(0,0)$. La statistique observée sur le premier essai ne suffit pas à déterminer la dynamique.\end{proof}

\begin{remarque}
Le fait qu'une statistique soit utile pour sélectionner des candidats ne la rend pas suffisante pour reconstruire une transition, un graphe causal ou une règle d'intervention. Le dossier inverse doit préserver autant que possible l'ordre, les conditions de contrôle et les métadonnées nécessaires à ces questions.
\end{remarque}

\chapter{Équivalence observationnelle, équivalence mécaniste et non-identifiabilité}
\section{Relations d'équivalence pertinentes}
\begin{definition}[Équivalence observationnelle relative à un protocole]
Deux mécanismes $M,N\in\Mcal$ sont observationnellement équivalents relativement à un protocole $\pi$ et à un domaine d'expériences $\mathcal E$ si les lois ou valeurs de sorties observables qu'ils prédisent sont identiques pour chaque expérience de $\mathcal E$. On note $M\equiv_{\mathrm{obs},\pi,\mathcal E}N$. Le domaine $\mathcal E$ inclut les entrées, contextes, sondes et interventions réellement autorisés; il n'est jamais remplacé silencieusement par tous les scénarios possibles.
\end{definition}

\begin{definition}[Isomorphisme mécaniste]
Deux mécanismes discrets $M$ et $N$ sont isomorphes si une bijection $h:Z_M\to Z_N$ préserve les états initiaux, les transitions, les sorties, les interfaces, le graphe de composants et les interventions, à renommage déclaré des identifiants près. La notion peut être adaptée aux états infinis ou aux modèles stochastiques en exigeant la préservation des noyaux de transition et d'émission.
\end{definition}

\begin{definition}[Identifiabilité relative]
Une classe $\Mcal$ est identifiable modulo $\MechEq$ sous le protocole $\pi$ sur $\mathcal E$ si
\[
M\equiv_{\mathrm{obs},\pi,\mathcal E}N\ \Longrightarrow\ M\MechEq N
\qquad\text{pour tous }M,N\in\Mcal.
\]
Cette propriété porte sur la classe, le protocole et le domaine annoncés. Elle n'est pas une propriété absolue d'une empreinte isolée.
\end{definition}

\begin{theoreme}[Non-identifiabilité par observations passives limitées]
Si deux mécanismes non isomorphes prédisent les mêmes observations pour toutes les expériences effectuées dans $\mathcal E_D$, alors aucun algorithme ne peut, à partir du seul dossier issu de $\mathcal E_D$, garantir lequel des deux est le mécanisme source sans une hypothèse supplémentaire qui exclut l'un d'eux ou une expérience qui les distingue.
\end{theoreme}
\begin{proof}
Soit $D$ le dossier produit par ces expériences. Par hypothèse, chacun des deux mécanismes est compatible avec tous les enregistrements de $D$. Un algorithme ne recevant que $D$ reçoit donc exactement la même entrée dans le cas où $M$ est la source et dans celui où $N$ est la source. Il doit produire la même sortie dans ces deux cas s'il est déterministe; dans le cas randomisé, sa distribution de sortie est également identique conditionnellement à la même entrée. Il ne peut garantir de distinguer les deux sources puisque la sortie correcte diffère entre ces cas. Une information nouvelle ou une restriction de classe est nécessaire.\end{proof}

\begin{contreexemple}[Deux mécanismes ayant la même entrée-sortie]
Sur le domaine d'essai déclaré $U_D=\{0,1\}$, considérons $M_1$ qui calcule $y=u$ directement et $M_2$ qui calcule d'abord $z=u$, puis $y=z$. Les sorties externes sont identiques sur tout $U_D$, mais les graphes internes comportent respectivement un et deux composants. Sans sonde ou intervention capable de distinguer $z$, la décomposition interne n'est pas identifiable.
\end{contreexemple}

\section{Renommage des états et limite de l'identité interne}
\begin{theoreme}[Invariance par renommage d'état]
Soit $M$ un mécanisme déterministe et $h:Z\to Z'$ une bijection. Construisons $M^h$ en posant $z'_0=h(z_0)$,
\[
\delta^h(z',u,a)=h\bigl(\delta(h^{-1}(z'),u,a)\bigr),\qquad
\omega^h(z',u,a)=\omega(h^{-1}(z'),u,a).
\]
Alors $M$ et $M^h$ produisent exactement les mêmes sorties pour toute séquence d'entrées et d'actions. Ils sont isomorphes comme systèmes d'états, mais leurs étiquettes internes ne sont pas identifiables par les sorties externes seules.
\end{theoreme}
\begin{proof}
Par induction sur $t$. À $t=0$, les états sont liés par $z'_0=h(z_0)$. Si $z'_t=h(z_t)$, alors
\[
z'_{t+1}=\delta^h(z'_t,u_t,a_t)=h(\delta(z_t,u_t,a_t))=h(z_{t+1}).
\]
Les sorties coïncident car $\omega^h(z'_t,u_t,a_t)=\omega(z_t,u_t,a_t)$. L'induction établit l'égalité des sorties pour chaque horizon fini.\end{proof}

\begin{corollaire}[Cible exacte de la reconstruction]
Sauf si le protocole fixe une représentation interne canonique ou inclut des sondes qui l'identifient, l'objectif défendable est une classe d'équivalence mécaniste, non une copie des noms, coordonnées ou bases internes arbitraires du système source.
\end{corollaire}
\begin{proof}
Le théorème précédent donne des mécanismes isomorphes et observationnellement identiques après transformation bijective des états. Les sorties externes ne peuvent donc sélectionner les noms d'état d'origine. Une assertion plus forte nécessite une observation interne additionnelle ou une convention canonique qui lève cette ambiguïté.\end{proof}

\section{Les données finies n'impliquent pas une généralisation universelle}
\begin{theoreme}[Limite d'identification sur un domaine infini]
Supposons l'espace d'entrées $U$ infini et le dossier $D$ ne contenant qu'un ensemble fini $U_D\subset U$. Dans une classe $\Mcal$ qui autorise des mécanismes à différer sur au moins une entrée $u^*\in U\setminus U_D$ sans modifier leurs sorties sur $U_D$, le dossier $D$ ne détermine pas les sorties sur tout $U$.
\end{theoreme}
\begin{proof}
Par hypothèse, il existe $M_1,M_2\in\Mcal$ tels que leurs prédictions concordent sur les entrées du dossier mais diffèrent sur $u^*$. Les deux sont donc compatibles avec $D$ et donnent des prédictions distinctes sur $u^*$. Le dossier ne permet pas de déduire laquelle est correcte.\end{proof}

\begin{remarque}
Cette proposition ne signifie pas qu'une généralisation est impossible. Elle signifie qu'une garantie de généralisation requiert une classe inductive, des hypothèses de régularité, une distribution cible, des données suffisantes ou une autre structure explicitement déclarée. La préférence pour un modèle simple est une règle de sélection, non une preuve que ce modèle est le bon mécanisme.
\end{remarque}

\chapter{Quand l'empreinte identifie-t-elle le mécanisme ?}
\section{Application expérimentale de signature}
Soit $\mathcal E$ le domaine d'expériences possibles. La signature expérimentale complète d'un mécanisme est
\[
\Psi_{\pi,\mathcal E}(M)=\bigl(\sem{M}(e)\bigr)_{e\in\mathcal E},
\]
où $\sem{M}(e)$ est la prédiction observable du mécanisme pour l'expérience $e$. Pour un système stochastique, cette composante est une loi, et non un échantillon unique, sauf si le protocole vise explicitement une réalisation particulière.

\begin{theoreme}[Critère suffisant d'identifiabilité]
Si l'application $\Psi_{\pi,\mathcal E}$ est injective sur les classes de $\Mcal$ modulo $\MechEq$, alors la signature expérimentale complète identifie une unique classe d'équivalence mécaniste. Réciproquement, si deux classes distinctes possèdent la même signature complète, cette signature ne les identifie pas.
\end{theoreme}
\begin{proof}
Si la signature de $M$ est connue et que $\Psi$ est injective sur le quotient, une autre classe ayant la même signature doit être la même classe; c'est l'identifiabilité. Inversement, si deux classes distinctes ont la même signature, elles satisfont les mêmes observations pour tout $e\in\mathcal E$, donc les données définies par cette signature ne peuvent les distinguer.\end{proof}

\begin{definition}[Expérience discriminante]
Pour des candidats déterministes, une expérience $e$ est discriminante pour $M,N$ si leurs valeurs prédites diffèrent. Pour des candidats stochastiques, elle est discriminante si leurs lois prédites sur la sortie observable diffèrent. Pour un ensemble $\mathcal V\subseteq\Mcal$, les résultats ou lois prédits doivent séparer au moins deux classes non vides de candidats.
\end{definition}

\begin{proposition}[Séparation exacte par expérience déterministe]
Si deux candidats déterministes $M,N$ sont compatibles avec le dossier courant, si une expérience admissible $e$ leur fait prédire deux sorties exactes distinctes et si l'exécution est valide avec mesure exacte, alors le résultat source élimine au moins l'un des deux candidats, ou bien contredit les deux et signale une inadéquation du domaine, de la classe ou du protocole.
\end{proposition}
\begin{proof}
Les deux sorties prédites étant distinctes, une valeur observée ne peut être égale aux deux. Si elle égale l'une d'elles, l'autre candidat est incompatible. Si elle n'égale aucune des deux, les deux sont incompatibles avec cet enregistrement sous le modèle déterministe exact; ce résultat doit être conservé comme alerte sur la classe, le protocole ou la mesure, et non corrigé en forçant un candidat.\end{proof}

\begin{samepage}
\begin{proposition}[Séparation de modèles stochastiques]
Si deux candidats stochastiques induisent des lois différentes pour une expérience, il existe un événement mesurable dont les probabilités prédites diffèrent. Une observation finie ne garantit toutefois pas toujours de déterminer avec certitude laquelle des lois est correcte; une conclusion statistique doit préciser le plan d'échantillonnage, le test, son seuil d'erreur et sa puissance visée.
\end{proposition}
\begin{proof}
Deux mesures de probabilité distinctes sont, par définition de leur inégalité, différentes sur au moins un ensemble mesurable $A$; donc $P_M(A)\ne P_N(A)$. Cette différence rend les lois séparables en principe. Elle ne signifie pas qu'une réalisation unique appartienne à la seule loi source : si les lois ont des supports qui se recouvrent, le même résultat peut avoir une probabilité positive sous les deux. Une décision à partir d'un échantillon fini requiert donc un test probabiliste avec risque d'erreur annoncé, sauf cas de supports disjoints et observation suffisante.
\end{proof}
\end{samepage}

\section{Identification locale, globale et dans la distribution cible}
\begin{definition}[Identifiabilité globale et locale]
L'identifiabilité est globale si l'ensemble inverse sur toute la classe $\Mcal$ ne contient qu'une classe modulo $\MechEq$. Elle est locale relativement à un voisinage ou sous-classe $\Mcal_0\subseteq\Mcal$ si cette propriété n'est établie que dans $\Mcal_0$. Une identifiabilité locale ne permet pas d'exclure les candidats hors de $\Mcal_0$.
\end{definition}

\begin{definition}[Domaine de validité empirique]
Le domaine de validité $\mathcal D_{\mathrm{val}}$ est l'ensemble des contextes, entrées, interventions, longueurs et régimes pour lesquels une validation a effectivement été réalisée et son critère satisfait. Toute extrapolation à un domaine plus grand est une revendication distincte nécessitant des hypothèses ou des tests dédiés.
\end{definition}

\begin{proposition}[La validation ne s'étend pas d'elle-même au-delà de son domaine]
Si deux candidats concordent sur $\mathcal D_{\mathrm{val}}$ mais diffèrent hors de ce domaine, alors la validation effectuée ne détermine pas lequel prédit correctement hors de $\mathcal D_{\mathrm{val}}$.
\end{proposition}
\begin{proof}
La conformité du domaine de validation ne compare les candidats que sur les expériences qui lui appartiennent. Par hypothèse, les prédictions sont égales sur ce domaine et différentes sur une expérience extérieure. Aucun résultat du domaine validé ne distingue donc ces prédictions extérieures.\end{proof}

\chapter{Reconstruction comme génération conditionnelle d'une observation}
\begin{definition}[Générateur observationnel]
Un générateur observationnel est un noyau \(Q(\dd y\mid i,v,h,p)\) ou, dans le cas déterministe, une fonction \(G(i,v,h,p)\), qui produit un enregistrement à partir d'une empreinte, de coordonnées variantes, d'une histoire \(h\) et d'un protocole \(p\). Ce générateur ne doit être considéré comme identifié que relativement à une classe déclarée et à des expériences discriminantes.
\end{definition}
\begin{proposition}[Un modèle invariant seul ne reproduit les variantes que si elles sont déterminées par lui]
Soit \(I:\Y\to\Z\). Il existe un générateur déterministe \(G:\Z\to\Y\) satisfaisant \(G(I(y))=y\) pour tout \(y\in\Y\) si et seulement si \(I\) est injective. Si \(I\) n'est pas injective, un générateur \(G(I(y))\) peut au mieux choisir un représentant ou reproduire une distribution conditionnelle, mais pas garantir la reconstruction exacte de l'observation individuelle.
\end{proposition}
\begin{proof}
La première assertion est le critère de reconstruction précédemment démontré. Si l'empreinte est non injective, il existe au moins deux entrées avec la même empreinte, donc un seul résultat \(G(I(y))\) ne peut être simultanément égal aux deux. Pour un générateur stochastique, il peut représenter la distribution conditionnelle des observations compatibles si celle-ci est estimable, mais cela n'identifie pas la réalisation individuelle perdue.\end{proof}
\begin{definition}[Reconstruction fidèle à un protocole]
Un candidat est fidèle à un protocole \(\pi\) sur un domaine \(D\) si la distance déclarée entre la loi de ses enregistrements et celle du source est sous une tolérance préétablie pour toute condition testée dans \(D\). La fidélité sur \(D\) ne vaut pas hors de \(D\) sans hypothèse de généralisation.
\end{definition}

\chapter{Synthèse de mécanismes et représentation des dépendances}
\section{Structure d'un mécanisme candidat}
\begin{definition}[Spécification mécaniste]
Une spécification mécaniste $\mathfrak m$ comprend au minimum : (1) les variables d'état et leur type; (2) les entrées et conditions de contexte; (3) les équations de transition; (4) les fonctions ou lois d'émission; (5) le graphe des composants; (6) les interfaces; (7) les interventions définies; (8) le domaine de validité; (9) les incertitudes; (10) les prédictions réfutables; (11) les expériences justificatives et contradictoires; (12) les données nécessaires à la reproduction. Une spécification qui n'indique pas comment calculer sa sortie n'est pas encore un mécanisme exécutable.
\end{definition}

\begin{definition}[Graphe mécaniste]
Un graphe mécaniste est un graphe dirigé $\Gamma=(V,E)$ dont les sommets représentent des composants ou variables et dont les arêtes expriment une dépendance de calcul ou une dépendance causale candidate. Chaque arête possède un statut : structure déclarée, dépendance mesurée, effet d'intervention, ou relation validée. Une arête de graphe ne reçoit pas automatiquement le statut causal.
\end{definition}

\begin{proposition}[Nécessité d'un contrat sémantique pour chaque composant]
Un graphe de composants ne détermine pas à lui seul un comportement exécutable, sauf si chaque composant, chaque connexion, l'ordre d'exécution et les règles de mise à jour sont spécifiés de façon à définir une transition unique ou une loi de transition bien définie.
\end{proposition}
\begin{proof}
Considérons un même graphe $A\to B$. Il peut représenter $B:=A$, $B:=2A$, ou $B:=\1_{A>0}$; le graphe est inchangé alors que la fonction calculée diffère. Sans spécification des fonctions et de leur sémantique d'exécution, le graphe admet plusieurs comportements et ne définit donc pas un mécanisme unique.\end{proof}

\section{Problème de synthèse contraint}
\begin{definition}[Problème de synthèse]
La synthèse à partir de $D$ consiste à trouver des candidats dans $\Inv_{\pi,C}(D)$ ou à déclarer cet ensemble vide relativement à la classe explorée. Si une fonction de coût $J(M;D)$ est utilisée, la sélection est
\[
\widehat M\in\argmin_{M\in\Inv_{\pi,C}(D)}J(M;D).
\]
Le coût peut combiner erreur de prédiction, complexité et coût d'exécution, mais ses termes, normalisations et priorités doivent être annoncés. Si le minimum n'est pas atteint ou si l'ensemble est vide, l'algorithme doit retourner un statut d'échec/indétermination et non un candidat déclaré compatible.
\end{definition}

\begin{proposition}[Un optimum de synthèse n'est pas une preuve d'unicité]
Même si le problème de sélection possède un unique minimiseur $\widehat M$, l'ensemble inverse peut contenir d'autres candidats compatibles. L'unicité du minimiseur pour $J$ n'implique donc pas l'identifiabilité du mécanisme source.
\end{proposition}
\begin{proof}
L'unicité du minimiseur signifie seulement que $J(\widehat M;D)<J(N;D)$ pour chaque autre candidat compatible $N$, selon les conventions du problème. Elle n'implique pas que $N$ est incompatible avec $D$, ni que $\widehat M$ est isomorphe au mécanisme source. Le coût est un critère supplémentaire de préférence, pas une mesure de vérité mécanique à moins qu'un théorème supplémentaire ne relie le coût à cette vérité.\end{proof}

\section{Conserver les contradictions}
\begin{definition}[Contradiction empirique]
Une contradiction survient lorsqu'un enregistrement valide est incompatible avec une prédiction précédente selon le critère et le modèle de mesure préenregistrés. La contradiction est stockée avec son dossier de provenance, son contexte et la version exacte du mécanisme mis en défaut; elle ne doit être effacée ni rétroactivement réétiquetée comme bruit sans justification indépendante.
\end{definition}

\begin{proposition}[Une contradiction invalide une revendication universelle correspondante]
Si une prédiction universelle $P$ affirme que $M$ reproduit une propriété pour tous les contextes d'un domaine $\mathcal D$, et qu'un essai valide dans $\mathcal D$ est incompatible avec $P$, alors cette revendication universelle est fausse, sauf si l'essai ou son interprétation a été invalidé par une règle indépendante et documentée.
\end{proposition}
\begin{proof}
Une proposition universelle implique la propriété pour chaque élément de son domaine. Un contre-exemple valide dans ce domaine contredit directement cette implication. Pour rejeter le contre-exemple, il faut établir que les prémisses d'application ne sont pas satisfaites; cela constitue une justification supplémentaire, non une suppression arbitraire du résultat.\end{proof}

\chapter{Doubles parcours top-down et bottom-up}
\section{Top-down}
Le parcours top-down part des comportements globaux et des conditions de tâche : exemples préenregistrés, changements d'entrée, distributions de sorties, erreurs, latence et dépendance au contexte. Il descend ensuite vers les composants qui pourraient expliquer les différences. Un comportement global est une contrainte sur l'explication, pas une preuve de structure interne.

\section{Bottom-up}
Le parcours bottom-up commence avec les composants accessibles : paramètres, couches, représentations, portes, neurones ou dimensions, puis examine les dépendances entre ces éléments et les effets de leur perturbation. Une dimension fortement variable n'est pas automatiquement une feature sémantique; une forte corrélation n'est pas un circuit causal.

\begin{theoreme}[Condition de confrontation des deux parcours]
Soit $H$ une hypothèse de mécanisme, $D_{\mathrm{top}}$ des observations comportementales et $D_{\mathrm{bottom}}$ des mesures internes. Si $H$ implique une prédiction testable $q$ sur chacun des deux domaines, et si l'une des observations indépendantes contredit $q$ au-delà de l'incertitude déclarée, $H$ ne peut être consolidée sous la forme testée.
\end{theoreme}
\begin{proof}
La consolidation de $H$ sous la forme testée exige que ses conséquences observables soient compatibles avec les données retenues et leur modèle d'incertitude. Une contradiction validée réfute au moins cette formulation de $H$ dans le domaine considéré. Le résultat peut conduire à réviser $H$, restreindre son domaine ou contester la mesure; il n'autorise pas la suppression silencieuse de la contradiction.
\end{proof}

\section{Composition récursive}
Le pipeline récursif est : observer un objet; former des signatures; construire des hypothèses; définir des interventions; éliminer des alternatives; proposer des mécanismes; composer des interfaces; observer le système composé comme un nouvel objet; comparer son empreinte à celle de la source; réviser les mécanismes. L'arrêt d'une étape ne doit jamais être remplacé par une déclaration narrative de réussite.

\begin{proposition}[Pas de montée en abstraction par simple renommage]
Renommer une représentation ou regrouper des mécanismes sans changer leur domaine, leurs opérations, leurs contrats ou les preuves ne crée pas de nouvelle distinction observationnelle ni de preuve de validité causale.
\end{proposition}
\begin{proof}
Le renommage est une bijection de labels qui préserve toutes les valeurs mesurées. Il ne modifie donc pas les fibres de l'empreinte et n'ajoute aucune observation. La validité causale est une propriété des résultats d'intervention et de leurs contrôles; aucun de ces résultats n'est créé par le changement de nom.
\end{proof}

\chapter{Abstraction mécaniste et quotient des états}
\section{Pourquoi l'abstraction doit préserver une sémantique}
\begin{definition}[Application d'abstraction]
Soit $M$ un mécanisme à espace d'état $Z$. Une abstraction est une application $\alpha:Z\to Q$ vers un espace d'état abstrait $Q$. L'abstraction associe les états concrets à des états abstraits, mais ne définit un mécanisme quotient que si les transitions et sorties abstraites sont indépendantes du représentant concret choisi.
\end{definition}

\begin{definition}[Congruence exacte pour un mécanisme déterministe]
Une application $\alpha$ induit une abstraction exacte si, pour tous $z_1,z_2$ avec $\alpha(z_1)=\alpha(z_2)$ et pour chaque entrée $u$ et action $a$ admissibles, (i) les sorties observables coïncident, $\omega(z_1,u,a)=\omega(z_2,u,a)$, et (ii) les prochains états appartiennent à la même classe, $\alpha(\delta(z_1,u,a))=\alpha(\delta(z_2,u,a))$.\end{definition}

\begin{theoreme}[Existence du mécanisme quotient exact]
Si $\alpha$ satisfait la congruence exacte sur un domaine fermé par les transitions, alors il existe un mécanisme abstrait $M_\alpha$ dont l'état est $q=\alpha(z)$, dont la sortie est $\omega_\alpha(q,u,a)=\omega(z,u,a)$ et dont la transition est $\delta_\alpha(q,u,a)=\alpha(\delta(z,u,a))$, avec $z$ un représentant quelconque de $q$. Ces définitions sont bien posées et l'abstraction reproduit exactement les sorties de $M$ sur ce domaine.
\end{theoreme}
\begin{proof}
Pour que $\omega_\alpha$ soit bien définie, les sorties doivent être identiques pour tous les représentants $z$ de $q$; c'est la condition (i). Pour que $\delta_\alpha$ soit bien définie, l'image par $\alpha$ du prochain état doit être identique pour tous ces représentants; c'est la condition (ii). La fermeture garantit que le prochain état reste dans le domaine. Par induction sur le nombre de pas, le mécanisme abstrait suit l'image $\alpha(z_t)$ de chaque état concret et émet la même sortie à chaque pas.\end{proof}

\begin{proposition}[Contre-exemple à l'abstraction arbitraire]
Si deux états regroupés par $\alpha$ produisent des sorties différentes sous une entrée admissible, alors aucune fonction de sortie déterministe $\omega_\alpha(q,u,a)$ ne peut reproduire exactement les deux comportements pour ce même $q,u,a$.
\end{proposition}
\begin{proof}
Une fonction déterministe assigne une unique valeur à un triplet donné. Si deux représentants du même état abstrait demandent deux sorties différentes, la fonction devrait prendre deux valeurs distinctes au même argument. C'est impossible.\end{proof}

\section{Abstraction stochastique}
\begin{definition}[Congruence stochastique]
Pour un mécanisme à noyau conjoint $L$, l'abstraction mesurable $\alpha$ est exacte si, pour tout couple $z_1,z_2$ tel que $\alpha(z_1)=\alpha(z_2)$, toute entrée-action admissible $(u,a)$, tout ensemble mesurable $C\subseteq Y$ et tout ensemble mesurable $B\subseteq Q$, on a
\[
L(C\times\alpha^{-1}(B)\mid z_1,u,a)=L(C\times\alpha^{-1}(B)\mid z_2,u,a).
\]
Cette condition compare la loi conjointe de la sortie et du prochain état abstrait, et non seulement ses deux marges séparées.\end{definition}

\begin{theoreme}[Mécanisme stochastique quotient]
Sous la congruence stochastique, avec $\alpha$ mesurable et les noyaux bien définis, la loi conjointe de la sortie et du prochain état abstrait est indépendante du représentant concret. Elle définit un mécanisme stochastique abstrait reproduisant les lois des trajectoires observables abstraites du mécanisme concret.
\end{theoreme}
\begin{proof}
La condition d'égalité sur tous les rectangles mesurables $C\times B$ assure que la loi conjointe de $(y_t,\alpha(z_{t+1}))$ ne dépend pas du représentant concret. Les rectangles déterminent la loi produit sur l'espace mesurable produit (sous le cadre mesurable déclaré), donc un noyau quotient est bien défini. En partant de l'état initial abstrait $\alpha(z_0)$ et en itérant ce noyau, une induction sur les distributions conditionnelles à chaque pas établit l'égalité des lois de trajectoires abstraites.\end{proof}

\chapter{Composition exécutable des mécanismes}
\section{Interfaces et composition synchrone}
\begin{definition}[Contrat d'interface]
Un contrat d'interface spécifie les types, unités, domaines, codomaines, conditions de validité, ordre temporel et modalités d'erreur des ports d'un mécanisme. Deux ports ne sont composables que si les contraintes de compatibilité sont satisfaites ou qu'un adaptateur explicite et testé réalise la conversion.
\end{definition}

\begin{definition}[Composition synchrone sans ambiguïté]
Deux mécanismes déterministes $M_1,M_2$ sont composables en parallèle si leurs états forment un produit $Z_1\times Z_2$, leurs mises à jour sont définies au même pas, les ports reliés ont des types compatibles et chaque entrée interne est une fonction explicitement spécifiée des états et sorties disponibles au début de ce pas. Les mises à jour simultanées utilisent ces valeurs anciennes; aucune boucle algébrique implicite n'est autorisée.
\end{definition}

\begin{theoreme}[Existence de la composition synchrone]
Sous les conditions de la définition précédente, la composition de $M_1$ et $M_2$ est un mécanisme déterministe bien défini sur l'espace d'état $Z_1\times Z_2$.
\end{theoreme}
\begin{proof}
Pour chaque état produit $(z_1,z_2)$, entrée externe et action, les contrats calculent sans ambiguïté les entrées internes à partir des valeurs disponibles au début du pas. Chaque mécanisme possède une fonction de transition totale sur son domaine déclaré; le couple des prochains états est donc unique. Les sorties externes sont obtenues par les fonctions de sortie et le routage déclarés. Ces opérations définissent une transition et une sortie sur le produit.\end{proof}

\begin{proposition}[La composition ne valide pas ses composants]
Même si une composition de mécanismes est mathématiquement bien définie, cette existence ne garantit pas que chacun de ses composants représente fidèlement le système source ni que la composition reproduit le comportement du source.
\end{proposition}
\begin{proof}
La définition de composition exige la cohérence des interfaces et l'existence des fonctions de calcul, non l'égalité de leurs prédictions avec celles du système source. Des fonctions erronées restent composables. La fidélité est une propriété empirique distincte qui doit être évaluée sur des expériences de comparaison.\end{proof}

\section{Provenance de composition}
\begin{definition}[Provenance structurelle]
Chaque composant reconstruit possède un identifiant stable, une version, ses dépendances, les données qui ont motivé sa création, les expériences utilisées, les contre-exemples rencontrés et la transformation d'abstraction qui le relie aux composants enfants. Une composition hérite d'un graphe de provenance qui relie chaque sortie aux opérations et preuves d'origine.
\end{definition}

\begin{proposition}[Propagation de provenance]
Si chaque composant d'une composition possède une provenance complète et si chaque connexion enregistre les composants amont et la transformation appliquée, alors la provenance de chaque sortie calculée peut être représentée comme un graphe fini de dépendances pour toute exécution finie d'un nombre fini de composants.
\end{proposition}
\begin{proof}
À l'exécution, une sortie de composant dépend d'un nombre fini de ses entrées, de son état et de sa version. Les entrées reliées portent elles-mêmes une provenance. La réunion finie de ces graphes et l'ajout des nœuds de calcul et de connexion construisent un graphe fini qui relie la sortie à ses dépendances. Le résultat ne garantit pas que la provenance soit vraie; il garantit seulement sa traçabilité structurelle si tous les liens sont correctement enregistrés.\end{proof}

\chapter{Hiérarchie récursive et raffinement des mécanismes}
\section{Montée et descente hiérarchiques}
Une reconstruction récursive alterne (i) une descente vers les détails observables, (ii) la création de composants exécutables, (iii) une montée par abstraction, puis (iv) la redescente expérimentale afin de tester si l'abstraction préserve les distinctions pertinentes. Une unité de niveau supérieur n'est pas réputée correcte parce que ses enfants le sont; elle est soumise à ses propres tests.

\begin{definition}[Échelle hiérarchique de mécanismes]
Une échelle hiérarchique est une suite de mécanismes $M_0,M_1,\ldots,M_k$ et d'applications de raffinement ou d'abstraction entre niveaux consécutifs. Chaque flèche précise le sens de la transformation, les ports préservés, les propriétés revendiquées et le domaine d'expériences sur lequel la relation a été vérifiée.
\end{definition}

\begin{proposition}[Préservation par composition des abstractions exactes]
Si $\alpha:Z_0\to Z_1$ définit une abstraction exacte de $M_0$ et $\beta:Z_1\to Z_2$ définit une abstraction exacte du mécanisme quotient obtenu, alors $\beta\circ\alpha$ définit une abstraction exacte de $M_0$ sur le domaine commun de validité.
\end{proposition}
\begin{proof}
L'abstraction $\alpha$ préserve les sorties et les transitions abstraites. L'abstraction $\beta$ préserve à son tour les sorties et transitions du mécanisme au niveau $Z_1$. Pour toute paire d'états concrets qui sont identifiés par $\beta\circ\alpha$, leurs images par $\alpha$ sont identifiées par $\beta$. La congruence de $\beta$ garantit alors l'égalité des sorties et des images d'états suivants au niveau $Z_2$. Le domaine commun de validité assure que les transitions considérées restent dans les domaines des deux abstractions.\end{proof}

\begin{theoreme}[Stabilisation finie sous espace de candidats fini]
Supposons que l'ensemble initial de candidats $\mathcal V_0$ soit fini, que chaque expérience ajoute une contrainte sans réécrire les anciennes données, et que chaque étape de mise à jour soit l'intersection exacte de l'ensemble courant avec les candidats satisfaisant cette contrainte. Alors la suite des ensembles compatibles $\mathcal V_0\supseteq\mathcal V_1\supseteq\cdots$ ne peut subir qu'un nombre fini de réductions strictes. Elle peut toutefois se stabiliser sur plusieurs candidats ou sur l'ensemble vide.
\end{theoreme}
\begin{proof}
Chaque réduction stricte enlève au moins un élément de l'ensemble fini $\mathcal V_0$. Il ne peut donc y avoir plus de $|\mathcal V_0|$ réductions strictes. La stabilité ne garantit pas l'unicité; si tous les candidats sont rejetés, cela indique une incompatibilité entre les contraintes et la classe explorée ou les données, et doit être signalé comme tel.\end{proof}

\begin{remarque}
Pour une classe infinie, une suite décroissante d'ensembles compatibles peut ne jamais se stabiliser après un nombre fini d'étapes. Aucune convergence algorithmique ne découle de la seule idée de récursion; elle requiert un résultat de terminaison, une limite de ressources ou une règle d'arrêt explicitement justifiée.
\end{remarque}

\part{V. Incertitude, provenance, validation et niveaux de preuve}

\chapter{Incertitude et données manquantes}
\section{Quatre sources à séparer}
Le protocole distingue : (1) variabilité intrinsèque du modèle ou de l'échantillonnage; (2) incertitude de mesure ou quantification; (3) variabilité de l'environnement et du runtime; (4) incertitude épistémique sur l'explication. Ces quatre sources peuvent interagir, mais ne doivent pas être fusionnées dans un seul score sans modèle explicite.

\begin{definition}[Valeur manquante typée]
Une absence d'observation est enregistrée avec un code, par exemple : sonde indisponible; timeout; erreur d'exécution; événement perdu; valeur censurée; accès interne refusé par le backend; capture interrompue; sortie non comparable. Ces états ne sont pas des valeurs numériques et ne peuvent participer à une moyenne comme si la valeur était zéro.
\end{definition}

\begin{proposition}[Remplacement de données manquantes par zéro invalide généralement l'empreinte]
Si une sonde peut produire une valeur réelle $0$ et qu'un état manquant est encodé par $0$, l'encodage n'est pas injectif sur l'ensemble des statuts de réponse : une mesure nulle et une non-mesure deviennent indiscernables.
\end{proposition}
\begin{proof}
L'application d'encodage envoie deux statuts distincts -- « valeur observée égale à zéro » et « valeur absente » -- vers la même représentation. Elle ne permet donc pas de retrouver le statut original; la signature sérialisée a perdu une distinction. Une représentation typée avec un champ valeur et un champ statut évite cette collision.
\end{proof}

\chapter{Provenance et intégrité de la preuve}
\section{Chaîne de provenance}
Chaque conclusion devrait pointer vers : hypothèse; expérience préenregistrée; modèle et hash; entrée et hash; sondes et versions; contexte; intervention; résultats bruts; métriques; contre-exemples; code et environnement; et rapport de validation. Un registre chaîne-de-hash permet de détecter certaines modifications postérieures à l'enregistrement; il ne démontre ni que les données initiales étaient vraies, ni que la bonne sonde a été utilisée, ni qu'aucune capture n'a été omise avant l'écriture.

\begin{definition}[Registre à chaînage de hachage]
Un registre contient des enveloppes $r_i$ avec payload $p_i$, hash du payload $h_i=H(p_i)$ et liaison au précédent $h_{i-1}$; le hash d'enveloppe est $g_i=H(h_{i-1}\Vert h_i\Vert m_i)$, où $m_i$ comprend les métadonnées canoniques. Le choix de sérialisation canonique et de fonction $H$ est versionné.
\end{definition}

\begin{proposition}[Portée de détection d'une chaîne de hachage]
Si la fonction de hachage résiste aux collisions dans le contexte considéré et si les valeurs de référence des hashes sont conservées hors de la portée d'un réécriveur, une modification d'un payload ou d'un lien antérieur entraîne une vérification incohérente, sauf découverte d'une collision ou réécriture de toute la chaîne et de sa référence externe.
\end{proposition}
\begin{proof}
Une modification du payload change son hash, sous l'hypothèse de résistance aux collisions. Le hash d'enveloppe qui inclut ce hash change ensuite, puis celui de tout successeur qui référence le précédent. La vérification compare les hashes calculés aux hashes enregistrés et signale la divergence. Si l'attaquant peut recalculer toute la chaîne et remplacer aussi la référence externe, le registre seul n'apporte plus de garantie; c'est pourquoi l'hypothèse de référence externe est explicitée.
\end{proof}

\section{Rejouabilité et identité d'une expérience}
La reproductibilité exige de conserver le modèle exacte (ou un hash cryptographique attestant les poids archivés), la version du code, les entrées, les tokens, les masques, la précision, les graines, les dépendances, le hardware pertinent, les versions des sondes, le plan d'acquisition et les paramètres d'intervention. Pour les environnements nondéterministes, le critère de reprise doit être statistique ou tolérant, pas une égalité bit-à-bit présumée.

\chapter{Validation indépendante du modèle reconstruit}
\section{Séparer ajustement, sélection et test}
\begin{definition}[Partitions expérimentales]
Le dossier est réparti, selon une règle déclarée avant le test final, en données de découverte $D_{\mathrm{dec}}$, données de sélection $D_{\mathrm{sel}}$ et données d'évaluation finale $D_{\mathrm{test}}$. Si les données sont temporelles, la partition respecte les dépendances temporelles pertinentes; si les mêmes prompts, états ou segments apparaissent dans plusieurs partitions par fuite, le test ne peut être qualifié d'indépendant sur ces unités.
\end{definition}

\begin{theoreme}[L'ajustement sur le test ne constitue pas une évaluation indépendante]
Si le choix du mécanisme final dépend des résultats de $D_{\mathrm{test}}$, alors ce même résultat ne peut être interprété comme une mesure indépendante non biaisée de la performance du mécanisme sélectionné sans analyse supplémentaire de la procédure adaptative.
\end{theoreme}
\begin{proof}
Le mécanisme sélectionné est une fonction des résultats du test. Les performances rapportées sont donc conditionnées par une sélection qui a exploité ces résultats. Elles ne proviennent plus d'un mécanisme fixé indépendamment du test. Une garantie exige alors un protocole d'inférence tenant compte des choix adaptatifs ou de nouvelles données de test jamais utilisées pour la sélection.\end{proof}

\section{Familles de métriques non substituables}
Un audit sérieux distingue au minimum : fidélité de sortie; fidélité de distribution; prédiction des états mesurables; fidélité interventionnelle; fidélité des trajectoires; généralisation hors distribution; robustesse; coût; autonomie; intégrité; provenance. Une amélioration d'une métrique n'implique pas automatiquement une amélioration des autres.

\begin{proposition}[Une moyenne agrégée peut cacher une régression par dimension]
Pour deux métriques $m_1,m_2$, il est possible qu'un modèle améliore leur moyenne pondérée tout en régressant sur $m_2$. Par conséquent, un score composite unique ne remplace pas le rapport de chacune des composantes.
\end{proposition}
\begin{proof}
Choisissons les scores du modèle A comme $(0.50,0.90)$ et ceux du modèle B comme $(0.95,0.55)$, avec poids égaux. Les moyennes valent respectivement $0.70$ et $0.75$ : B a donc une moyenne supérieure, alors que sa deuxième métrique régresse de $0.90$ à $0.55$.\end{proof}

\section{Contre-exemples actifs}
\begin{definition}[Contre-exemple discriminant]
Un contre-exemple est discriminant s'il réfute une prédiction déclarée ou sépare au moins deux mécanismes encore compatibles. Un cas simplement rare ne constitue pas en soi un contre-exemple; il doit être valide, dans le domaine revendiqué et mesuré avec un protocole interprétable.
\end{definition}

\begin{proposition}[La recherche de contre-exemples ne peut garantir leur découverte]
Aucun algorithme de recherche limité à un ensemble fini de requêtes ne peut garantir de trouver un contre-exemple pour toute paire de mécanismes qui ne diffèrent que sur une entrée non interrogée dans un domaine infini.
\end{proposition}
\begin{proof}
Un algorithme finit par interroger un ensemble fini $U_D$. Comme le domaine est infini, il existe une entrée $u^*$ non interrogée. On peut construire deux mécanismes qui coïncident sur $U_D$ et divergent uniquement sur $u^*$. Les réponses reçues sont identiques, donc l'algorithme ne peut savoir qu'un contre-exemple se trouve à $u^*$.\end{proof}

\chapter{Niveaux de réussite pour un modèle de langage reconstruit}
\begin{definition}[Niveaux de revendication]
On distingue les niveaux suivants :
\begin{enumerate}[label=\textbf{N\arabic*.}]
\item artefact chargeable et exécutable;
\item modèle de langage autonome dont les sorties dépendent de ses propres poids et de son propre état;
\item reproduction comportementale sur un domaine de test déclaré;
\item généralisation évaluée sur des contextes non utilisés pour l'ajustement;
\item reproduction des réponses à une famille d'interventions;
\item reproduction de trajectoires et d'états mesurables;
\item reconstruction de composants et dépendances identifiés modulo une équivalence déclarée;
\item composition hiérarchique validée, avec prédictions et interventions aux niveaux composites;
\item identification mécaniste relativement à une classe de modèles et à un protocole dont le pouvoir discriminant est démontré.
\end{enumerate}
Aucun niveau supérieur ne découle automatiquement d'un niveau inférieur.
\end{definition}

\begin{proposition}[La réussite comportementale n'implique pas l'équivalence mécaniste]
Il existe deux mécanismes qui produisent les mêmes sorties sur toutes les entrées du domaine externe considéré, mais qui possèdent des décompositions internes non isomorphes. La fidélité comportementale sur ce domaine ne suffit donc pas à démontrer l'équivalence mécaniste.
\end{proposition}
\begin{proof}
Le contre-exemple du calcul direct $y=u$ et du calcul composé $z=u$, $y=z$ donne deux graphes internes non isomorphes avec les mêmes sorties sur tout le domaine externe. Ils atteignent une fidélité comportementale parfaite sur ce domaine, mais ne possèdent pas la même décomposition.\end{proof}

\begin{proposition}[Condition minimale pour revendiquer une reconstruction mécaniste]
Une revendication de reconstruction mécaniste exige au minimum : un mécanisme exécutable autonome; une classe d'équivalence mécaniste déclarée; des prédictions pour des interventions; un jeu d'essais indépendant; la provenance des composants; et la description des candidats alternatifs qui restent compatibles. Si l'un de ces éléments manque, la revendication est soit réduite, soit déclarée non établie.
\end{proposition}
\begin{proof}
L'autonomie écarte la simple délégation cachée au source. L'équivalence déclarée précise ce qui est reconstruit. Les prédictions interventionnelles rendent testables les dépendances annoncées. Le test indépendant empêche de confondre sélection et validation. La provenance établit la traçabilité, et les candidats alternatifs rendent explicite l'identifiabilité. Sans chacune de ces pièces, au moins une ambiguïté fondamentale décrite dans les chapitres précédents demeure non résolue.\end{proof}

\chapter{Garanties possibles et hypothèses qu'il est interdit de cacher}
\section{Catalogue des dépendances logiques}
Toute affirmation de reconstruction doit préciser les hypothèses qu'elle utilise. Les principales sont :
\begin{enumerate}
\item \textbf{Adéquation de la classe :} le mécanisme recherché est représentable dans $\Mcal$, si l'on revendique une récupération du source.
\item \textbf{Adéquation de mesure :} les sondes mesurent bien les variables annoncées, avec un modèle d'erreur déclaré.
\item \textbf{Couverture d'expérience :} les entrées, contextes et interventions séparant les candidats sont accessibles et exécutés.
\item \textbf{Stabilité :} le mécanisme source n'a pas changé sans versionnement entre les expériences, ou ses changements sont modélisés.
\item \textbf{Absence de fuite :} les données de test n'influencent ni sélection ni réglages avant l'évaluation finale.
\item \textbf{Sémantique d'intervention :} l'action réellement exécutée correspond à l'intervention décrite, sans effets collatéraux omis dans l'interprétation.
\item \textbf{Calculabilité :} les procédures et espaces ont des représentations effectives suffisantes pour l'algorithme revendiqué.
\item \textbf{Capacité de calcul :} l'optimisation ou la recherche s'est terminée et le statut de terminaison est consigné.
\item \textbf{Fermeture des dépendances :} le modèle produit n'utilise pas un oracle non déclaré.
\item \textbf{Périmètre de généralisation :} la distribution cible et les conditions de validité sont explicitement définies.
\end{enumerate}

\begin{proposition}[Déplacer une hypothèse ne la supprime pas]
Si une preuve requiert une hypothèse $H$ à une étape intermédiaire, réexprimer cette étape au moyen d'une définition ou d'une procédure qui utilise encore $H$ ne rend pas la conclusion indépendante de $H$.
\end{proposition}
\begin{proof}
Une implication logique dont les prémisses incluent $H$ établit seulement la conclusion conditionnellement à $H$. Remplacer le nom de l'étape ou déplacer l'opération dans une fonction ne modifie pas les prémisses de l'implication. Pour supprimer la dépendance, il faut démontrer la propriété sans $H$ ou fournir une autre prémisse indépendante impliquant ce qui était nécessaire de $H$.\end{proof}

\section{Ce qui ne peut pas être garanti universellement}
Sans restriction de classe, les systèmes arbitraires peuvent coder des comportements non calculables ou exiger des ressources non bornées; sans couverture complète, des candidats peuvent différer hors des expériences vues; sans hypothèse sur les mesures, un biais peut être confondu avec le mécanisme; sans accès d'intervention, des dépendances causales peuvent rester non identifiables. Le formalisme ne transforme pas ces limites logiques en limites du projet : il les rend visibles afin de concevoir les expériences ou extensions qui peuvent les réduire, et de ne jamais annoncer une preuve que les données n'autorisent pas.

\chapter{Ce que cette théorie autorise et n'autorise pas à conclure}
\begin{longtable}{p{0.3\textwidth}p{0.29\textwidth}p{0.31\textwidth}}
\toprule Revendication & Conditions minimales & Conclusion non autorisée sans preuve supplémentaire \\\midrule
\endhead
Empreinte invariante & Transformations et domaine déclarés, propriété d'invariance établie & Dire que toute empreinte est intrinsèquement invariante \\
Reconstruction exacte & Injectivité de l'empreinte ou de la paire (invariant, variant) & Reconstituer une observation perdue avec l'invariant seul \\
Séparation invariant/variant unique & Contraintes d'identifiabilité, par exemple somme directe ou rang plein dans la classe choisie & Traiter une décomposition illustrative comme canonique \\
Observation continue fidèle & Horloges, échantillonnage, pertes et modèle d'erreur spécifiés & Déduire un signal continu arbitraire d'échantillons finis \\
Effet de sonde causal & Plan valide, interventions et confondeurs traités & Confondre différence temporelle ou corrélation avec causalité \\
Mécanisme reconstruit & Exécuteur autonome, expériences discriminantes et validation indépendante & Confondre une statistique, un graphe ou une bonne sortie avec une compréhension mécaniste \\
\bottomrule
\end{longtable}

\part{VI. Procédure complète, architecture et critères d’arrêt}

\chapter{Protocole d'observation profonde}
\section{Pré-enregistrement de la question}
Avant la collecte, consigner la question scientifique, la propriété qui pourrait varier, les hypothèses concurrentes, les prédictions différentielles, les conditions de réfutation, le domaine de validité visé et les résultats qui conduiraient à ne pas conclure. Les sondes exploratoires peuvent être ajoutées ensuite, mais elles sont étiquetées comme exploratoires et ne remplacent pas la validation tenue à l'écart.

\section{Phase A -- objet et intégrité}
\begin{enumerate}
\item Identifier l'artefact exact, le modèle/tokenizer, les poids, les fichiers de configuration, le runtime et le hash.
\item Construire une carte de capacités par sonde et ne jamais simuler une capacité absente.
\item Distinguer données d'audit, données de découverte, données de validation et données de contre-test.
\item Définir le contrat de l'observateur et mesurer le coût de l'instrumentation.
\end{enumerate}

\section{Phase B -- observation passive multi-échelle}
\begin{enumerate}
\item Niveau global : tâches, entrées, distributions de sorties, erreurs, latence, répétitions et contextes.
\item Niveau composants : structure du graphe, paramètres, dimensions, chemins de calcul et normalisations.
\item Niveau sous-structures : tenseurs, activations, têtes, états, gradients ou scores lorsque l'accès le permet.
\item Niveau temporel : répétitions, préfixes, trajectoires génératives, évolution de cache/état, incidents et dérives.
\item Niveau opérationnel : mémoire, temps, débit, exceptions, effet de l'instrumentation et événements manquants.
\item Pour chaque niveau, publier le brut utile et les résumés; les résumés n'effacent pas les données sources nécessaires à la réanalyse.
\end{enumerate}

\section{Phase C -- choix et comparaison de sondes}
\begin{enumerate}
\item Construire des familles de sondes définies par propriété mesurée, pas uniquement par nom de catégorie.
\item Évaluer les duplications, les coûts et les distinctions ajoutées.
\item Construire les signatures par contexte, échelle et fenêtre temporelle.
\item Examiner les fibres : quels objets/états/trajectoires restent indiscernables ?
\item Pour chaque conclusion importante, rechercher une sonde susceptible de distinguer l'hypothèse choisie de ses alternatives.
\end{enumerate}

\section{Phase D -- interventions et identification}
\begin{enumerate}
\item Préenregistrer la cible, l'emplacement, l'opérateur exact, l'effet attendu et les contrôles.
\item Mesurer baseline et intervention sous contextes appariés; restaurer l'état entre les essais lorsque requis.
\item Utiliser des contrôles positifs, négatifs et des interventions de comparaison.
\item Répliquer sur plusieurs entrées, séquences, graines ou régimes temporels selon la question.
\item Mesurer les résultats bruts, tailles d'effet, incertitudes et contre-exemples.
\item Ne conclure à un mécanisme que si des hypothèses alternatives pertinentes sont discriminées dans un domaine annoncé.
\end{enumerate}

\section{Phase E -- abstraction et reconstruction}
\begin{enumerate}
\item Formaliser le mécanisme candidat comme transducteur avec interfaces, conditions et opérations.
\item Relier chaque composant à ses preuves et contre-preuves.
\item Composer uniquement les interfaces compatibles et conserver la provenance des enfants.
\item Étudier le modèle composé comme un nouvel objet, avec ses propres sondes, contextes et signatures.
\item Comparer au système source sur des tests indépendants : comportement, généralisation, robustesse et interventions.
\item Conserver les divergences et remonter à l'hypothèse ou au composant qui pourrait les expliquer.
\end{enumerate}

\section{Critères de sortie}
Une étape ne peut être déclarée terminée que si les artefacts existent, que les tests associés ont été exécutés, que les résultats et erreurs sont enregistrés, que la provenance est complète et que la règle de réussite avait été définie avant de consulter les données de test. Un niveau ne réussit pas automatiquement parce que le niveau précédent a réussi.

\chapter{Spécification d'architecture d'un observateur continu}
Cette spécification est une conséquence du formalisme; elle ne prétend pas que les interfaces décrites sont déjà présentes dans le code v3.1.

\section{Contrat abstrait d'une sonde}
Une sonde logicielle devrait exposer au minimum :
\begin{itemize}
\item \texttt{probe\_id}, version, description et facette;
\item domaine de cible et de contexte; préconditions; capacité exigée;
\item fonction d'acquisition ponctuelle et, si supporté, fonction d'acquisition en flux;
\item codomaine, unités, règles de normalisation et statut des valeurs absentes;
\item coût attendu, taux maximal, perturbation possible et mécanisme d'arrêt;
\item méthode de vérification, sérialisation canonique et provenance.
\end{itemize}

\section{Contrat d'une session live}
Une session devrait avoir des opérations explicites \texttt{start}, \texttt{pause}, \texttt{resume}, \texttt{snapshot}, \texttt{subscribe}, \texttt{stop} et \texttt{finalize}. Chaque opération retourne un reçu. La session possède un plan multi-échelle, les cibles actives, les fréquences, les budgets, l'état du tampon, la dernière séquence reçue, les erreurs, les watermarks temporels et les hashes des lots. Si le backend ne peut pas fournir une propriété, il déclare cette incapacité au lieu de simuler la continuité.

\section{Contrat d'un événement et d'un lot}
L'événement individuel porte les champs de la définition précédente. Un lot ajoute son numéro, la première et la dernière séquence, les séquences manquantes connues, le temps de début/fin, les volumes, le hash du lot et le reçu de stockage. Les événements tardifs ne sont pas réordonnés en silence : la politique de traitement des retards est versionnée, et l'ordre causal est distingué de l'ordre d'arrivée.

\section{Ordonnanceur de sondes}
Un ordonnanceur répartit le budget entre sondes coûteuses et légères. Il doit estimer le coût, interrompre ou réduire une fréquence selon une politique déclarée, enregistrer toute adaptation de fréquence et éviter de traiter des fenêtres de fréquences différentes comme si elles étaient identiques. Les choix adaptatifs sont basés exclusivement sur le préfixe d'information déjà reçu; la décision et sa justification sont enregistrées.

\section{Sonde intrusive et effet de l'observation}
Lorsqu'un hook, un wrapper, une capture de gradient ou une intervention modifie le runtime, l'observateur conduit une expérience de perturbation : comparer avec instrumentation, sans instrumentation, et sous plusieurs taux d'échantillonnage. Si l'écart dépasse le budget toléré, les observations sont marquées comme potentiellement intrusives. Le qualificatif « passif » est attaché à la procédure; il ne signifie pas zéro coût ou zéro influence.

\chapter{Procédure minimale invariant--variant--dynamique}
La procédure suivante est une spécification de recherche, pas une garantie de réussite universelle.
\begin{enumerate}[label=\textbf{P\arabic*.}]
\item Déclarer le système, le protocole, le domaine d'entrée, les états accessibles, les familles de sondes, les contextes, les interventions et les contraintes de ressources.
\item Collecter et conserver les événements bruts, y compris les pertes, horodatages et erreurs; construire les empreintes à plusieurs facettes sans remplacer les données brutes par des seules statistiques.
\item Déclarer la relation d'invariance : quelles transformations, changements de contexte ou interventions sont supposés préserver la composante structurelle.
\item Définir séparément les coordonnées variantes : contexte, facette, intervention, histoire, trajectoire, stochastique et résidu de mesure.
\item Tester l'identifiabilité de la décomposition dans la classe choisie. Si elle échoue, préserver plusieurs décompositions compatibles au lieu d'en sélectionner arbitrairement une.
\item Définir un ensemble de mécanismes candidats, leur sémantique exécutable, leurs interfaces et leurs dépendances.
\item Choisir des expériences qui discriminent les candidats, avec contrôles positifs et négatifs adaptés; enregistrer les prédictions avant exécution.
\item Construire et exécuter le candidat sans appel oracle caché au source; rejouer des trajectoires lorsque le contrat le permet.
\item Ré-observer le candidat sous le même ensemble de sondes, de contextes et d'interventions; comparer les signatures multi-échelles, les trajectoires et les incertitudes.
\item Valider sur des données indépendantes; enregistrer divergences, contre-exemples, alternatives encore compatibles et domaine de validité.
\end{enumerate}
\begin{proposition}[Préservation de l'incertitude dans une reconstruction honnête]
Si le dossier et la classe de mécanismes laissent plusieurs candidats non équivalents compatibles, un rapport qui ne retient qu'un candidat doit déclarer une règle de sélection distincte de la preuve d'identification. Sans cette déclaration, la conclusion d'unicité ne suit pas des données.
\end{proposition}
\begin{proof}
Par définition, plusieurs candidats non équivalents satisfont les contraintes disponibles. Le dossier ne permet donc pas de déduire que l'un d'eux est le seul mécanisme possible. Une règle de sélection peut ordonner les candidats, par exemple par coût, mais l'ordre dépend alors du critère et ne réduit pas l'ensemble des compatibles au sens logique.\end{proof}

\chapter{Procédure inverse complète}
\section{Étapes opérationnelles}
\begin{algorithme}[Reconstruction disciplinée à partir d'une empreinte]
\begin{enumerate}
\item \textbf{Geler le protocole.} Identifiants, versions, contexte, sondes, interventions, budgets, métriques, règle de validité et partitions sont préenregistrés.
\item \textbf{Valider le dossier.} Vérifier provenance, ordre temporel, événements manquants, hashes, erreurs d'acquisition, doublons et statut de chaque mesure.
\item \textbf{Déclarer la classe $\Mcal$.} Décrire la grammaire des mécanismes candidats, les interfaces et les restrictions architecturales avant de voir le test final.
\item \textbf{Créer une version compatible.} Chercher un ensemble de mécanismes satisfaisant les données de découverte. Si aucun candidat ne convient, conserver le statut « classe ou hypothèses incompatibles ».
\item \textbf{Identifier les ambiguïtés.} Chercher des candidats non équivalents ayant les mêmes prédictions sur le dossier actuel.
\item \textbf{Planifier des sondes discriminantes.} Sélectionner des entrées ou interventions qui séparent les candidats, sous contraintes de coût, d'éthique, d'accès et de sécurité expérimentale.
\item \textbf{Exécuter et enregistrer.} Acquérir les résultats, y compris négatifs et indéterminés, avec leurs métadonnées et leur provenance.
\item \textbf{Mettre à jour.} Réduire l'ensemble compatible seulement selon le critère préenregistré; si l'ensemble devient vide, arrêter la prétention de compatibilité et diagnostiquer.
\item \textbf{Construire un exécuteur.} Sérialiser états, transitions, sorties, graphe, interfaces, interventions et dépendances; vérifier qu'aucun appel caché au source n'existe.
\item \textbf{Abstraire et composer.} Utiliser uniquement des abstractions dont les conditions de préservation sont prouvées ou testées comme approximations avec bornes déclarées.
\item \textbf{Valider hors sélection.} Exécuter les essais réservés, les contre-exemples, les interventions et les changements de distribution prévus.
\item \textbf{Publier le périmètre.} Déclarer métriques, limites, alternatives encore compatibles, versions, hashes, erreurs, commandes exactes et niveau de réussite établi.
\end{enumerate}
\end{algorithme}

\section{Propriété de correction limitée de la mise à jour}
\begin{theoreme}[Correction logique de l'ensemble compatible]
Si (i) la classe $\Mcal$ est fixée, (ii) le critère $C$ est correctement appliqué, (iii) chaque événement valide est ajouté sans réécriture silencieuse et (iv) l'ensemble compatible initial est exact, alors après $n$ mises à jour l'ensemble calculé est exactement l'ensemble des candidats de $\Mcal$ qui satisfont toutes les contraintes enregistrées.
\end{theoreme}
\begin{proof}
Par induction sur $n$. À $n=0$, c'est l'hypothèse (iv). Supposons l'égalité vraie à $n$. La mise à jour définie par intersection enlève précisément les candidats qui échouent à la contrainte nouvellement ajoutée. Le résultat est donc l'ensemble des candidats qui satisfont les $n+1$ contraintes.\end{proof}

\begin{remarque}
Ce théorème porte sur la correction logique de l'ensemble dans la classe définie. Il ne garantit ni que la classe contient le mécanisme source, ni que le critère est scientifiquement adéquat, ni que les données sont suffisantes, ni que le calcul pratique résout exactement le problème. Ces obligations sont distinctes et doivent être auditées séparément.
\end{remarque}

\chapter{Comment NGeGenBrain doit convertir une empreinte en mécanisme}
\section{Contrat minimal du pipeline}
Le pipeline doit séparer les modules suivants :
\begin{enumerate}
\item \textbf{EvidenceStore :} conserve valeurs brutes, événements temporels, métadonnées, erreurs et provenance.
\item \textbf{ProbeRegistry :} déclare la sémantique, portée, coût, version et possibilité d'intervention de chaque sonde.
\item \textbf{HypothesisSpace :} représente les classes et candidats, leurs prédictions et leurs conditions d'incompatibilité.
\item \textbf{ExperimentPlanner :} choisit des expériences susceptibles de séparer les candidats non équivalents encore compatibles.
\item \textbf{MechanismIR :} représentation intermédiaire typée des états, transitions, fonctions, ports, interventions et incertitudes.
\item \textbf{Synthesizer :} produit et sélectionne des mécanismes exécutables sans consulter le test réservé.
\item \textbf{InterventionRunner :} exécute l'action annoncée et ses contrôles, puis compare source et candidat.
\item \textbf{Composer/Abstractor :} compose et abstrait seulement sous les contrats sémantiques déclarés.
\item \textbf{Validator :} exécute validation indépendante et contre-exemples, sans écraser les contradictions.
\item \textbf{Auditor/Provenance :} lie chaque affirmation aux données, scripts, versions, hashes et tests correspondants.
\end{enumerate}

\begin{proposition}[La séparation des modules ne garantit pas à elle seule la correction scientifique]
Même si ces modules existent séparément, le pipeline n'est correct que si leurs contrats sont respectés et si les sorties de chaque étape sont vérifiées contre ces contrats.
\end{proposition}
\begin{proof}
Une implémentation peut contenir tous les modules mais transmettre un identifiant de modèle erroné, omettre une intervention, utiliser le test pour la sélection ou convertir un résultat manquant en zéro. L'existence nominale des modules ne prévient pas ces erreurs. Les garanties exigent la vérification des contrats et de la chaîne de provenance.\end{proof}

\section{Matrice de preuves requise}
\begin{longtable}{>{\raggedright\arraybackslash}p{0.22\textwidth}>{\raggedright\arraybackslash}p{0.35\textwidth}>{\raggedright\arraybackslash}p{0.35\textwidth}}
\toprule
\textbf{Revendication} & \textbf{Preuve minimale} & \textbf{Ce qui ne suffit pas}\\
\midrule
Empreinte reproductible & Protocole versionné, données brutes/hashes, horodatages, vérification de rejeu & Un fichier JSON de statistiques\\
Mécanisme exécutable & Équations/transition opérationnelles, exécution sans oracle caché, tests de sérialisation & Un graphe de composants seul\\
Dépendance causale candidate & Intervention ciblée, contrôle, effet mesuré, contexte et incertitude & Corrélation ou attention élevée\\
Abstraction valide & Condition de congruence démontrée ou erreur d'approximation bornée & Regroupement de neurones selon similarité seule\\
Composition valide & Contrats de ports et tests du mécanisme composé & Addition de composants dans un registre\\
Généralisation & Test tenu à l'écart, domaine décrit, contre-exemples & Bonne performance sur le jeu d'ajustement\\
Équivalence mécaniste & Identification modulo une relation déclarée et exclusion des alternatives pertinentes & Texte généré similaire au source\\
Modèle autonome & Audit des dépendances et test sans accès au source & Un checkpoint qui se charge\\
\bottomrule
\end{longtable}

\chapter{Critères d'arrêt honnêtes et continuation de recherche}
\begin{definition}[Arrêt scientifique d'une phase]
Une phase est terminée lorsque son objectif mesurable est satisfait ou lorsqu'un blocage est caractérisé avec preuves, contre-exemples, hypothèses en cause et opérations restantes. Un résultat négatif documenté peut terminer une phase d'évaluation, mais il ne doit pas être renommé réussite ni être effacé.
\end{definition}

\begin{proposition}[Un résultat négatif peut guider une expérience discriminante]
Si un candidat échoue à une prédiction sur une entrée valide, l'enregistrement de ce contre-exemple peut réduire l'ensemble compatible et guider la sélection d'une expérience qui distingue les candidats restants, sous réserve que le critère de compatibilité et la classe restent définis.
\end{proposition}
\begin{proof}
Le contre-exemple élimine chaque candidat dont la prédiction est incompatible selon le critère. Les candidats qui restent forment un sous-ensemble. Si plusieurs restent et ne sont pas équivalents, leurs prédictions permettent de rechercher une expérience où elles divergent; l'existence d'une telle expérience doit être vérifiée et non supposée.\end{proof}

\begin{critere}[Publication d'un mécanisme reconstruit]
Une publication NGeGenBrain ne doit employer le terme « mécanisme reconstruit » que si elle fournit au minimum : (i) l'empreinte et son protocole; (ii) la classe de mécanismes et la relation d'équivalence; (iii) un exécuteur autonome; (iv) des prédictions préenregistrées; (v) des interventions réellement exécutées et des contrôles; (vi) une évaluation indépendante; (vii) la provenance et les hashes; (viii) les candidats alternatifs encore compatibles; (ix) un domaine de validité et les limitations observées. Sinon, employer le statut approprié : résumé, règle candidate, simulateur comportemental, mécanisme partiel ou hypothèse non identifiée.
\end{critere}

\part{VII. Audit logiciel, limites présentes et obligations de réalisation}

\chapter{Audit du pipeline d'observation de NGeGenBrain v3.1}

\begin{remarque}[Portée historique]Ce chapitre reprend un audit consigné dans la source initiale. Il ne vaut pas nouvelle exécution du code v3.1 dans le cadre de la présente refonte ; ses constats restent attachés au périmètre et aux preuves cités dans cette source.\end{remarque}

\section{Périmètre de l'inspection}
L'inspection porte sur les fichiers livrés dans l'archive v3.1 extraite pour l'audit. Les conclusions de cette section sont des constats de code, non des jugements fondés uniquement sur l'intention documentaire.

\begin{longtable}{p{0.23\textwidth}p{0.32\textwidth}p{0.32\textwidth}}
\toprule
\textbf{Composant examiné} & \textbf{Comportement constaté} & \textbf{Portée / lacune observée}\\
\midrule
\endhead
\path{adapters/base.py} & \texttt{capture(inputs)} par défaut retourne le résultat d'un appel \texttt{forward} et un dictionnaire vide. & Aucun protocole temporel commun, aucune abstraction d'abonnement ou d'événement durable dans l'interface de base.\\
\path{adapters/torch_adapter.py} & Installe des hooks de sortie de modules pour une exécution, collecte les valeurs, puis retire les hooks. Les valeurs sont converties en structures Python sérialisables. & Capture ponctuelle, centrée sur les sorties des modules. Elle ne définit pas une histoire continue ni un ordonnanceur de sondes. La conversion en listes peut être coûteuse en mémoire et en débit.\\
\path{adapters/huggingface.py} & La méthode \texttt{capture} obtient hidden states et attentions, puis appelle aussi l'adaptateur Torch pour les activations de modules. & Deux passages de modèle sont exécutés pour un même appel de capture. L'identité de l'exécution n'est donc pas garantie pour des systèmes stochastiques ou avec état mutable, sauf hypothèses supplémentaires de déterminisme et d'absence d'effets de bord.\\
\path{analysis/causal.py} & \texttt{observe} réalise un passage et émet un dictionnaire horodaté; \texttt{intervention} compare une baseline et une sortie altérée sur une entrée. & Pas d'historique temporel continu dans ces méthodes. Un delta ponctuel est une mesure d'effet local, non une identification complète de mécanisme.\\
\path{analysis/activations.py} & Stockage compressé par couche, hash des contenus sérialisés et registre à chaîne de hachage. & Intégrité des fichiers enregistrés vérifiable; cela ne prouve ni que toutes les activations utiles furent capturées, ni qu'aucun événement ne fut omis avant l'écriture. Les valeurs passent par JSON.\\
\path{analysis/features.py} & Calcule moyennes, variances, fréquence d'activité, valeurs distinctes et corrélations sur des captures fournies. Le texte du module qualifie explicitement ces statistiques de descriptives. & Pas de sémantique garantie des dimensions; corrélation ne vaut pas causalité. Aucun mécanisme de recherche universelle des sondes n'est établi par ces fonctions seules.\\
\path{analysis/hierarchical.py} & Compose des mécanismes enregistrés en proposition d'abstraction et indique qu'une validation indépendante est requise. & La composition hiérarchique est candidate; elle n'établit pas automatiquement la validité causale ou sémantique des enfants ou du parent.\\
\texttt{web.py}, route \path{/api/model-observe} & Pour le modèle chargé, mesure un prompt ponctuel, des statistiques des états cachés et des attentions lorsque l'objet le permet. & Une requête ponctuelle n'est pas une observation live et continue. Les résumés agrégés ne constituent pas à eux seuls une empreinte multi-échelle complète.\\
\texttt{web.py}, route \path{/api/observe} & Crée un adaptateur synthétique de parité et l'observe. & Cette route n'observe pas le modèle LLM chargé; elle doit être comprise comme une démo synthétique, pas comme une inspection du modèle actif.\\
\path{web/index.html} & L'interface appelle des routes d'observation, de sweep, de résumé et de stockage à la demande. & Aucun contrat général d'abonnement durable à un flux du modèle, de gestion des événements manquants ou de reconstruction de trajectoire temporelle n'a été trouvé dans les modules inspectés.\\
\bottomrule
\end{longtable}

\section{Conclusion d'audit}
La v3.1 dispose de primitives utiles : inspection structurale, captures d'activations, capture cachée/attention pour certains adaptateurs, statistiques descriptives, interventions, registre et composition hiérarchique proposée. En revanche, son observation effective est principalement \emph{à la demande et par passage}. Elle n'implémente pas, au sens général défini ici, une observation live continue possédant une horloge, une gestion des retards et pertes, un planificateur de sondes multi-échelles, un modèle d'incertitude temporelle et une comparaison de trajectoires.

Il en résulte une distinction cruciale : la théorie proposée est un formalisme normatif pour concevoir et auditer une observation plus profonde; elle ne doit pas être lue comme l'affirmation que tous ses composants existent déjà dans v3.1.

\section{Pourquoi plusieurs façades ne signifient pas plusieurs étiquettes de texte}
Les facettes d'observation sont des familles de mesures qui interrogent des propriétés différentes du même objet : comportement externe, structure, représentations, trajectoires, sensibilité aux entrées, effets d'interventions, gradients quand accessibles, stochasticité, robustesse et changement de contexte. Produire huit catégories lexicales ou entraîner un routeur de documents ne remplace pas ce protocole. Une facette est définie par ce qu'elle mesure, son domaine de validité, son coût et les distinctions qu'elle permet d'établir.

\chapter{Matrice d'audit v3.1 par rapport au formalisme}

\begin{remarque}[Portée historique]Ce chapitre reprend un audit consigné dans la source initiale. Il ne vaut pas nouvelle exécution du code v3.1 dans le cadre de la présente refonte ; ses constats restent attachés au périmètre et aux preuves cités dans cette source.\end{remarque}

\begin{longtable}{p{0.25\textwidth}p{0.40\textwidth}p{0.26\textwidth}}
\toprule
\textbf{Exigence du formalisme} & \textbf{État constaté dans v3.1} & \textbf{Statut honnête}\\
\midrule
\endhead
Objet, hash et capacités & Modèles décrits et hashés par les adaptateurs; matrice de capacités présente. & Présent pour les adapters couverts; la granularité de la provenance doit être contrôlée par expérience.\\
Observation de sortie & Passage direct et capture des sorties disponibles. & Présent, ponctuel.\\
Activations internes & Hooks Torch et capture d'états cachés pour certaines architectures. & Présent, mais capture ponctuelle; parfois sérialisation lourde.\\
Attention & Capture Hugging Face si le modèle la renvoie. & Conditionnel à l'architecture/backend.\\
Gradients & Méthode dédiée pour certaines entrées/sorties. & Conditionnel; pas équivalent à une observation live.\\
Intervention & Interventions sur des sorties de modules et mesures de delta. & Présent; sémantique dépend du point de hook et du modèle.\\
Capture unique cohérente & Le chemin HF capture états/attentions puis exécute une autre capture via Torch hooks. & À consolider pour garantir un unique passage instrumenté ou une équivalence de passes démontrée.\\
Sondes et fermeture & Statistiques descriptives et corrélations des captures. & Pas de fermeture générale ni preuve d'équivalence de sondes dans le code inspecté.\\
Seuils et fibres & Formalisés ci-dessus, pas implémentés comme moteur général identifié dans les modules inspectés. & Théorie spécifiée; implémentation non revendiquée.\\
Flux live durable & Pas d'API générale de session d'observation continue trouvée dans les modules inspectés. & Absent du périmètre inspecté; à concevoir et tester.\\
Temps d'événement / arrivée & Observations ponctuelles ont un timestamp, sans protocole universel d'horloges et de watermark. & Insuffisant pour la théorie live complète.\\
Séquences manquantes / backpressure & Pas de contrat général découvert dans l'adaptateur de base et la route d'observation. & Absent du périmètre inspecté.\\
Provenance des activations & Enregistrements hashés et stockés par couche. & Utile pour l'intégrité après enregistrement, pas une preuve d'exhaustivité de la capture.\\
Interprétation sémantique & Les modules descriptifs ne prétendent pas assigner une sémantique par variance ou corrélation. & Conforme à la prudence requise; le protocole d'identification reste à enrichir.\\
Abstraction hiérarchique & Composition de mécanismes enregistrés comme candidat et checklist de validation. & Présent comme composition; pas équivalence causale automatique.\\
\bottomrule
\end{longtable}

\chapter{Questions ouvertes et obligations d'implémentation}
Le formalisme laisse plusieurs choix à instancier pour un modèle donné : codomaines de mesure; métriques; échelles; tolérances; fréquences; politique de ressources; modèle de bruit; modèle d'intervention; classe de mécanismes candidates; procédure d'inférence; tests de transport hors distribution. Ces choix ne sont pas des faiblesses cachées : ils déterminent la portée exacte de chaque assertion. Ils doivent être enregistrés, motivés et soumis à une analyse de sensibilité.

Pour atteindre une observation profonde et vivante, la prochaine implémentation doit ajouter un contrat commun de flux à l'adaptateur, un ordonnanceur multi-échelle, des événements avec temps d'événement et temps d'arrivée, une gestion de pertes et de backpressure, une capture cohérente d'une exécution, des sessions réversibles, un stockage adapté aux tenseurs volumineux et un pipeline de comparaison de trajectoires. Ces composants doivent être suivis par des tests de contrat, des tests d'intrusion, des tests de charge, des tests de perte/réordonnancement, de reprise après panne, ainsi que des tests de falsification sur des modèles synthétiques dont la vérité n'est pas fournie au moteur de découverte.

\appendix
\part{Registres, notations et traçabilité documentaire}

\chapter{Notation unifiée}
\begin{longtable}{>{\raggedright\arraybackslash}p{0.25\textwidth} >{\raggedright\arraybackslash}p{0.65\textwidth}}
\toprule
Symbole & Sens retenu dans le document unifié \\
\midrule\endhead
$\Omega$ & Espace des exécutions ou trajectoires source. \\
$\Ycal$ & Espace des observations ou enregistrements admissibles. \\
$\pi$ & Protocole : contexte, entrées, sondes, interventions, ordonnanceur, budgets, versions et règles de mesure. \\
$O_\pi$ & Application d'observation déterministe ; pour le cas stochastique, noyau d'observation déclaré. \\
$I_\pi$ & Empreinte ou application de signature relative au protocole. \\
$G$ & Famille déclarée de transformations relativement à laquelle une invariance peut être affirmée. \\
$S$ & Composante structurelle/invariante, seulement si son invariance est démontrée sous $G$. \\
$V$ & Composante variante ou résiduelle visant à distinguer des observations appartenant à une même fibre de $S$. \\
$\tau$ & Trace ordonnée d'événements/états et de métadonnées temporelles. \\
$\mathfrak D$ & Dossier d'évidence avec provenance, incertitudes, données et résultats négatifs. \\
$\Mcal$ & Classe déclarée de mécanismes candidats. \\
$\Inv_{\pi,C}(\mathfrak D)$ & Ensemble de mécanismes compatibles avec le dossier selon le protocole et le critère $C$. \\
$\sim_I$ & Équivalence observationnelle induite par l'égalité des empreintes. \\
$\MechEq$ & Équivalence mécaniste, toujours relative à une définition de mécanisme et de morphisme déclarée. \\
$\Lambda,\mathfrak F,\Delta$ & Ordres d'échelle de résolution, de famille de sondes et de résolution temporelle. \\
$\Sigma$ & Produit d'indexation multi-échelle ; l'ordre produit ne vaut que si chaque axe possède un sens de raffinement explicite. \\
\bottomrule
\end{longtable}
\paragraph{Conventions.} Les sources emploient parfois $I$ pour des objets de nature différente. Le présent document réserve $I_\pi$ à une empreinte (application de signature) et $\mathfrak D$ à un dossier d'évidence. Une classe d'équivalence d'états n'est jamais identifiée à l'état brut. Une quantité « stable » au sens empirique ne devient pas un invariant exact sans quantificateurs et hypothèses supplémentaires.

\chapter{Registres de résultats, checklists et conclusions des documents sources}
Les sous-sections suivantes conservent les registres synthétiques et les checklists portés par les sources. Elles sont documentaires : en cas de différence de vocabulaire, les définitions et règles d'harmonisation de ce volume prévalent. Ces registres ne remplacent pas les démonstrations dans le corps du traité.

\section*{Résumé des résultats établis}
Les propositions suivantes sont démontrées dans le domaine et sous les hypothèses annoncées :
\begin{enumerate}
\item l'indiscernabilité exacte induite par une empreinte est une relation d'équivalence;
\item l'extension d'une famille de sondes affine la partition et ne fusionne pas les classes;
\item une propriété est fonction de l'empreinte si et seulement si elle est constante sur chacune de ses fibres;
\item la fermeture observationnelle définie globalement est extensive, monotone, idempotente et maximale pour la conservation des fibres;
\item la hiérarchie des échelles engendre une hiérarchie de partitions; sur un domaine fini, les raffinements stricts sont bornés;
\item la fonction de seuil satisfait la loi du minimum double sous l'hypothèse explicite d'une chaîne d'échelles emboîtées; elle induit une pseudoultramétrique;
\item des échantillons ponctuels ne suffisent pas à reconstruire un signal arbitraire continu; une borne de reconstruction exige une hypothèse de régularité;
\item une différence après intervention établit un effet local mesuré, mais non à elle seule un mécanisme sémantique complet;
\item une famille finie de tests ne permet pas d'identifier sans hypothèse supplémentaire le comportement sur un domaine hors test;
\item la composition fonctionnelle est mathématiquement distincte de la validation causale de l'abstraction composée.
\end{enumerate}

\section{Ce qui n'est pas un théorème du présent document}
Les affirmations selon lesquelles une famille de sondes donnée découvrirait toutes les fonctions pertinentes, l'observation finirait toujours par séparer tous les objets, les mécanismes internes seraient toujours identifiables, ou une fréquence arbitrairement élevée fournirait une observation sans perturbation sont fausses en général ou exigent des hypothèses non établies. Aucune n'est revendiquée ici.

\section*{Checklist de conformité d'une étude d'observation}
\begin{enumerate}[label=\textbf{C\arabic*.}]
\item Le modèle/objet est identifié par artefact, version et hash.
\item Toutes les sondes utilisées ont un contrat et une version.
\item Les entrées, contexte, graines et paramètres sont enregistrés.
\item Les capacités disponibles et absentes sont explicites.
\item Les observations brutes sont distinguées des résumés et inférences.
\item Les facettes mesurent des propriétés différentes ou leur redondance est établie.
\item Les échelles de sondes sont emboîtées lorsqu'une loi de raffinement est invoquée.
\item Les données manquantes ne sont pas codées comme des zéros observés.
\item Pour le live, temps, séquences, retards, pertes, pause/reprise et fin de flux sont définis.
\item La charge de l'observateur et sa perturbation sont mesurées ou déclarées inconnues.
\item Les hypothèses concurrentes et les résultats discriminants sont préenregistrés.
\item Les interventions modifient réellement la cible selon une opération documentée.
\item Les baselines, contrôles et restaurations sont valides.
\item Les effets sont présentés avec le domaine et l'incertitude appropriés.
\item Les contre-exemples et résultats défavorables sont préservés.
\item Les propriétés de fermeture ou seuils ne sont pas prétendues prouvées par simple échantillon, sauf qualification empirique.
\item Les abstractions composées sont testées comme de nouveaux objets.
\item Les conclusions importantes ont une chaîne de provenance vers les données brutes.
\item Les données de validation/test n'ont pas été utilisées pour sélectionner les hypothèses ou les seuils, sauf si le protocole le déclare explicitement comme exploration.
\item L'étude indique ce qui reste indéterminé avec les sondes actuelles.
\end{enumerate}

\section*{Notation}
\begin{longtable}{p{0.23\textwidth}p{0.68\textwidth}}
\toprule
\textbf{Notation} & \textbf{Sens}\\
\midrule
\endhead
$\X$ & Domaine des objets étudiés : modèles, états, représentations ou trajectoires selon le type déclaré.\\
$\Pcal$ & Ensemble des contextes d'observation.\\
$\Fcal$ & Famille de sondes admissibles.\\
$\Gcal$ & Sous-famille de sondes utilisées pour une empreinte.\\
$\Phi(p,x,f)$ & Réponse du protocole/sonde $f$ appliqué à l'objet $x$ sous contexte $p$.\\
$I_{p,\Gcal}(x)$ & Empreinte de $x$ selon les sondes de $\Gcal$.\\
$\sim_{p,\Gcal}$ & Indiscernabilité observationnelle exacte.\\
$[x]$ & Fibre/classe d'équivalence de $x$.\\
$\Lcal_{p,\Gcal}$ & Ensemble des empreintes réalisables.\\
$\cl_p(\Gcal)$ & Fermeture observationnelle globale de la famille $\Gcal$.\\
$\Lambda,\preceq$ & Ordre d'échelle des familles de sondes.\\
$\iota(\lambda)$ & Intensité déclarée d'une échelle.\\
$\Theta_p(\{x,y\})$ & Seuil de séparation d'une paire.\\
$d_\Theta$ & Pseudoultramétrique dérivée du seuil.\\
$z(t)$ & Trajectoire d'état d'exécution.\\
$t^{\mathrm{event}}$ & Temps source déclaré de l'événement.\\
$t^{\mathrm{capture}}$ & Temps auquel le capteur a acquis la valeur.\\
$t^{\mathrm{arrival}}$ & Temps auquel le collecteur a reçu l'événement.\\
$\Fcal_t^{\mathrm{arr}}$ & Information disponible au collecteur à la date $t$.\\
$\Delta_a$ & Effet local mesuré d'une intervention $a$.\\
$\tau_a$ & Effet moyen sous la distribution de contextes déclarée.\\
\bottomrule
\end{longtable}

\section*{Contrôles de cohérence mathématique}
\begin{enumerate}
\item La relation d'indiscernabilité est définie par égalité de signatures exactes; la version tolérante est explicitement notée comme n'étant pas automatiquement transitive.
\item Le théorème de factorisation ne requiert pas que l'empreinte soit injective; il requiert exactement que $A$ soit constante sur ses fibres.
\item La fermeture est définie par factorisation globale sur le domaine; l'approximation empirique est distinguée.
\item La loi du minimum double utilise une chaîne d'échelles emboîtées et une intensité monotone; aucune généralité plus forte n'est implicitement revendiquée.
\item La fonction $e^{-\Theta}$ fournit une pseudoultramétrique; elle n'est une métrique que si les seuils entre objets distincts sont finis.
\item La reconstruction d'un signal continu à partir d'échantillons est déclarée impossible sans hypothèse de régularité; la borne Lipschitz indique les hypothèses et la constante nécessaires.
\item Un delta d'intervention est distingué de l'identification complète d'un mécanisme; la non-identifiabilité est illustrée par deux systèmes qui coïncident sur le domaine testé.
\item La chaîne de hash protège contre certaines modifications postérieures sous résistance aux collisions et référence externe; elle n'est pas une preuve de vérité ou d'exhaustivité des acquisitions.
\end{enumerate}

\backmatter

\section*{Conclusion -- source base}
Une observation profonde n'est pas une capture isolée, une liste d'activations ou un score résumé. C'est une famille versionnée de sondes, appliquées à des objets et trajectoires sous des contextes déclarés, à plusieurs échelles, avec une distinction explicite entre valeur, absence, incertitude, mesure et inférence. L'observation continue ajoute l'horloge, le flux d'événements, le traitement des retards et pertes, les ressources et l'effet de l'instrumentation. La causalité ajoute des interventions bien définies, des contrôles et une logique d'identification. La hiérarchie ajoute des abstractions et compositions dont la validation reste séparée de leur construction.

L'audit v3.1 montre que plusieurs briques utiles existent, mais que le sens fort de l'observation live, continue, multi-échelle et causalement discriminante n'est pas encore réalisé par le pipeline ponctuel inspecté. Le présent formalisme fixe les distinctions mathématiques nécessaires pour corriger ce manque sans prétendre que la théorie, à elle seule, aurait déjà produit l'implémentation ou les preuves expérimentales absentes.

\section*{Théorèmes centraux de la théorie inverse}
Le présent document établit les points suivants, toujours relativement aux hypothèses annoncées :
\begin{enumerate}
\item l'inverse d'une empreinte est naturellement un ensemble de candidats, non une fonction unique;
\item la sélection par coût ne prouve pas l'identifiabilité;
\item des mécanismes non isomorphes peuvent avoir les mêmes observations sur un protocole limité;
\item les renommages d'état préservent les sorties et empêchent l'identification des étiquettes internes sans conventions ou sondes supplémentaires;
\item des données finies ne déterminent pas un comportement universel sur un domaine infini sans structure de classe;
\item une application expérimentale injective modulo l'équivalence mécaniste est un critère suffisant d'identifiabilité relative;
\item un delta sous intervention n'est pas à lui seul une preuve de nécessité, de suffisance ou d'unicité causale;
\item une abstraction exacte existe si les transitions et sorties sont congruentes sur les classes d'abstraction;
\item une composition synchrone est exécutable si ses interfaces et règles de mise à jour sont non ambiguës;
\item la composition et la récursion ne garantissent ni fidélité au source ni convergence sans hypothèses supplémentaires;
\item les mises à jour incrémentales de l'ensemble compatible sont monotones quand elles ajoutent uniquement des contraintes valides;
\item validation comportementale, fidélité interventionnelle et identification mécaniste sont des objectifs distincts.
\end{enumerate}

\section*{Checklist finale de rigueur}
Avant d'affirmer qu'une empreinte a été transformée en mécanisme reconstructible, vérifier toutes les questions suivantes :
\begin{enumerate}[label=\textbf{R\arabic*.}]
\item La classe de mécanismes candidats et ses restrictions sont-elles écrites avant la sélection finale ?
\item Le dossier inclut-il protocole, entrées, contexte, interventions, ordre, erreurs, incertitudes et provenance ?
\item Les statistiques dérivées sont-elles distinguées des observations brutes ?
\item Les événements manquants et horloges sont-ils modélisés au lieu d'être imputés silencieusement ?
\item L'ensemble des candidats compatibles est-il conservé, ou une règle de sélection est-elle présentée comme telle ?
\item Les ambiguïtés et mécanismes non identifiables sont-ils explicités ?
\item Des expériences discriminantes ont-elles été recherchées parmi les candidats encore compatibles ?
\item Les interventions sont-elles réellement exécutées, avec contrôle de référence et contrôle négatif/positif pertinent ?
\item Chaque composant possède-t-il des équations ou une sémantique d'exécution, au-delà du graphe ?
\item L'exécuteur est-il autonome relativement au paquet déclaré ?
\item Les abstractions satisfont-elles un théorème de congruence ou disposent-elles de bornes d'erreur mesurées ?
\item Les ports et ordres d'exécution rendent-ils la composition non ambiguë ?
\item La validation est-elle indépendante des données utilisées pour choisir le mécanisme ?
\item Les contre-exemples, contradictions et résultats défavorables sont-ils gardés ?
\item Les métriques de comportement, causalité, représentation et robustesse sont-elles rapportées séparément ?
\item Les états de réussite et le domaine de validité sont-ils nommés sans généralisation excédentaire ?
\item Les poids, codes, données, configurations et rapports ont-ils des hashes et versions reproductibles ?
\item Le rapport distingue-t-il ce qui est prouvé, observé, supposé, inféré et inconnu ?
\end{enumerate}

\section*{Notation}
\begin{longtable}{p{0.22\textwidth}p{0.70\textwidth}}
\toprule
\textbf{Symbole} & \textbf{Signification}\\
\midrule
$S$ & Système source observé\\
$\pi$ & Protocole d'observation et d'expérimentation\\
$D$ & Dossier d'évidence enregistré\\
$\mathfrak D_D$ & Dossier empirique structuré dérivé du dossier\\
$\Mcal$ & Classe de mécanismes candidats déclarée\\
$M,N$ & Mécanismes candidats\\
$\Inv_{\pi,C}(D)$ & Ensemble de mécanismes compatibles avec le dossier\\
$\Psi_{\pi,\mathcal E}$ & Signature expérimentale d'un mécanisme sur le domaine $\mathcal E$\\
$\ObsEq$ & Équivalence observationnelle relative à un protocole et un domaine\\
$\MechEq$ & Équivalence mécaniste déclarée\\
$Z,z_0$ & Espace d'état et état initial\\
$\delta,\omega$ & Transition et émission déterministes\\
$K,G$ & Noyaux de transition et d'émission stochastiques\\
$\Gamma$ & Graphe de composants et dépendances candidates\\
$\alpha$ & Application d'abstraction entre états\\
$\mathcal V_t$ & Ensemble de candidats compatibles après $t$ mises à jour\\
$\mathcal D_{\mathrm{val}}$ & Domaine de validité effectivement évalué\\
\bottomrule
\end{longtable}

\section*{Conclusion -- source inverse}
Le mécanisme inverse de l'empreinte n'est pas une opération magique transformant toute collection de mesures en « compréhension » certaine. C'est un problème de synthèse et d'identification dont la solution dépend de trois objets explicites : une classe de mécanismes, un protocole expérimental qui différencie ces mécanismes, et un critère de validation indépendant. L'empreinte fournit des contraintes sur le mécanisme; les interventions ajoutent des contraintes contrefactuelles; la composition produit des candidats exécutables; la validation tente de les réfuter et précise leur domaine de validité.

La conclusion opérationnelle pour NGeGenBrain est donc précise : produire une empreinte, même profonde et multi-facettes, ne suffit pas. Il faut convertir ses éléments en variables, transitions, interfaces, mécanismes candidats et prédictions; identifier les ambiguïtés restantes; choisir des expériences discriminantes; construire un exécuteur autonome; puis valider séparément comportements, interventions, trajectoires et abstractions. Lorsqu'une seule étape ne peut être justifiée, le système doit garder le statut correspondant et poursuivre l'étude, plutôt que de substituer une explication plausible à une preuve.

\bigskip
\begin{center}
\textbf{Principe final}\par\medskip
\fbox{\parbox{0.88\textwidth}{\centering
Empreinte $\longrightarrow$ ensemble de mécanismes compatibles $\longrightarrow$ expériences discriminantes $\longrightarrow$ mécanisme exécutable candidat $\longrightarrow$ contre-tests indépendants $\longrightarrow$ abstraction et composition validées.\\[0.4em]
L'identification n'est acquise que relativement à une classe, un protocole, une équivalence et un domaine déclarés.}}
\end{center}

\section*{Registre des résultats démontrés}
Les résultats formels ci-dessus sont démontrés au niveau mathématique de leurs définitions : (1) non-reconstructibilité à partir d'une application non injective; (2) critère de reconstruction exacte et de factorisation; (3) correspondance quotient--image; (4) critère invariant plus variante; (5) non-unicité additive sans contraintes; (6) unicité d'une somme directe; (7) distinction logique entre invariance et stationnarité; (8) raffinement par ajout de coordonnées; (9) non-identifiabilité d'une trace continue depuis des temps finis sans régularité; (10) borne d'erreur sous Lipschitz; (11) confusion entre facteurs non appariés; (12) identifiabilité du modèle linéaire au rang plein. Ces théorèmes ne prouvent pas que leurs hypothèses valent pour un LLM donné : cela relève de l'audit et de l'expérimentation.

\section*{Checklist d'application}
\begin{enumerate}[label=\textbf{I\arabic*.}]
\item Quel groupe de transformations définit l'invariance ?
\item L'invariance est-elle exacte ou approximative, et avec quelle métrique et tolérance ?
\item Quelles coordonnées de l'observation sont effacées par l'empreinte ?
\item Quelle variable/trace variante est conservée pour permettre une reconstruction ?
\item L'application combinée est-elle injective sur le domaine ciblé ?
\item Si la décomposition est additive, quelle contrainte assure son unicité ?
\item Le temps est-il mesuré comme état, trajectoire, événement, fenêtre ou flux ?
\item Les pertes, délais, horloges et effets intrusifs des sondes sont-ils déclarés ?
\item Les expériences séparent-elles la sonde du contexte ?
\item Les résultats de l'étudiant ont-ils été ré-observés sous le même protocole que la source ?
\item Les conclusions négatives, modèles concurrents et limites sont-ils conservés ?
\end{enumerate}

\section*{Conclusion -- source invariant}
L'idée décisive est de séparer l'empreinte de l'observation complète. Une empreinte peut être conçue pour retenir des invariants, mais cette propriété doit être définie et démontrée relativement à un ensemble de transformations. Les variations ne forment pas automatiquement un résidu identifiable : leur séparation d'avec l'invariant requiert des contraintes de modèle ou un plan d'expériences suffisamment informatif. Le temps ajoute l'histoire et l'échantillonnage; il ne transforme pas à lui seul une information partielle en trace complète.

Pour NGeGenBrain, le programme théorique devient : \emph{trace brute et provenance} \(\to\) \emph{famille d'empreintes multi-échelles} \(\to\) \emph{décomposition invariant/variant/histoire} \(\to\) \emph{ensemble de mécanismes compatibles} \(\to\) \emph{expériences discriminantes} \(\to\) \emph{candidat exécutable} \(\to\) \emph{ré-observation et validation}. Chaque flèche possède son contrat et ses conditions de validité; aucune ne doit être remplacée par une explication narrative.

\chapter{Matrice de traçabilité des sources}
Le tableau indique l'emplacement de chaque chapitre des trois sources exploitées. « Intégration » signifie que le chapitre complet est repris dans le corps théorique, sous réserve des corrections terminologiques explicites décrites dans le présent volume. « Annexe » signifie que le contenu sert surtout de registre, checklist, notation ou conclusion récapitulative et a été conservé en section documentaire.
\begin{longtable}{>{\raggedright\arraybackslash}p{0.14\textwidth} >{\raggedright\arraybackslash}p{0.39\textwidth} >{\raggedright\arraybackslash}p{0.14\textwidth} >{\raggedright\arraybackslash}p{0.25\textwidth}}
\toprule Source & Chapitre source & Destination & Décision \\
\midrule\endhead

Formalisme initial & Objet, portée et statuts épistémiques & Corps & Contenu et preuves repris \\

Formalisme initial & Objets, contextes, sondes et mesures & Corps & Contenu et preuves repris \\

Formalisme initial & Empreinte informationnelle et identité observationnelle & Corps & Contenu et preuves repris \\

Next 1 - invariant/variant & Pourquoi l'empreinte ne doit pas être confondue avec l'observation entière & Corps & Contenu et preuves repris \\

Next 1 - invariant/variant & Quotient invariant et information résiduelle & Corps & Contenu et preuves repris \\

Next 1 - invariant/variant & Quand une décomposition additive est-elle définie de façon unique ? & Corps & Contenu et preuves repris \\

Formalisme initial & Redondance et fermeture observationnelle & Corps & Contenu et preuves repris \\

Next 0 - théorie inverse & De l'empreinte à la reconstruction : objets et statuts & Corps & Contenu et preuves repris \\

Formalisme initial & Échelles de sondes et capacité observationnelle & Corps & Contenu et preuves repris \\

Formalisme initial & Fonction de seuil et géométrie des fibres & Corps & Contenu et preuves repris \\

Next 1 - invariant/variant & Un ordre produit des échelles observationnelles & Corps & Contenu et preuves repris \\

Formalisme initial & Trajectoires et signatures temporelles & Corps & Contenu et preuves repris \\

Next 1 - invariant/variant & Traces live, histoire et reconstruction temporelle & Corps & Contenu et preuves repris \\

Formalisme initial & Stabilité, non-stationnarité et estimation en ligne & Corps & Contenu et preuves repris \\

Next 0 - théorie inverse & Empreinte de trajectoire et prédiction continue & Corps & Contenu et preuves repris \\

Next 0 - théorie inverse & Mise à jour en ligne de l'ensemble inverse & Corps & Contenu et preuves repris \\

Formalisme initial & Facettes observationnelles complémentaires & Corps & Contenu et preuves repris \\

Formalisme initial & Sélection active de sondes et recherche discriminante & Corps & Contenu et preuves repris \\

Formalisme initial & Interventions, contrôles et effets mesurés & Corps & Contenu et preuves repris \\

Next 1 - invariant/variant & Conditions d'identifiabilité des composants invariant et variant & Corps & Contenu et preuves repris \\

Next 0 - théorie inverse & Interventions et reconstruction causale & Corps & Contenu et preuves repris \\

Formalisme initial & Définition d'un mécanisme exécutable & Corps & Contenu et preuves repris \\

Next 0 - théorie inverse & Sémantique exécutable d'un mécanisme & Corps & Contenu et preuves repris \\

Next 0 - théorie inverse & Équivalence observationnelle, équivalence mécaniste et non-identifiabilité & Corps & Contenu et preuves repris \\

Next 0 - théorie inverse & Quand l'empreinte identifie-t-elle le mécanisme ? & Corps & Contenu et preuves repris \\

Next 1 - invariant/variant & Reconstruction comme génération conditionnelle d'une observation & Corps & Contenu et preuves repris \\

Next 0 - théorie inverse & Synthèse de mécanismes et représentation des dépendances & Corps & Contenu et preuves repris \\

Formalisme initial & Doubles parcours top-down et bottom-up & Corps & Contenu et preuves repris \\

Next 0 - théorie inverse & Abstraction mécaniste et quotient des états & Corps & Contenu et preuves repris \\

Next 0 - théorie inverse & Composition exécutable des mécanismes & Corps & Contenu et preuves repris \\

Next 0 - théorie inverse & Hiérarchie récursive et raffinement des mécanismes & Corps & Contenu et preuves repris \\

Formalisme initial & Incertitude et données manquantes & Corps & Contenu et preuves repris \\

Formalisme initial & Provenance et intégrité de la preuve & Corps & Contenu et preuves repris \\

Next 0 - théorie inverse & Validation indépendante du modèle reconstruit & Corps & Contenu et preuves repris \\

Next 0 - théorie inverse & Niveaux de réussite pour un modèle de langage reconstruit & Corps & Contenu et preuves repris \\

Next 0 - théorie inverse & Garanties possibles et hypothèses qu'il est interdit de cacher & Corps & Contenu et preuves repris \\

Next 1 - invariant/variant & Ce que cette théorie autorise et n'autorise pas à conclure & Corps & Contenu et preuves repris \\

Formalisme initial & Protocole d'observation profonde & Corps & Contenu et preuves repris \\

Formalisme initial & Spécification d'architecture d'un observateur continu & Corps & Contenu et preuves repris \\

Next 1 - invariant/variant & Procédure minimale invariant--variant--dynamique & Corps & Contenu et preuves repris \\

Next 0 - théorie inverse & Procédure inverse complète & Corps & Contenu et preuves repris \\

Next 0 - théorie inverse & Comment NGeGenBrain doit convertir une empreinte en mécanisme & Corps & Contenu et preuves repris \\

Next 0 - théorie inverse & Critères d'arrêt honnêtes et continuation de recherche & Corps & Contenu et preuves repris \\

Formalisme initial & Audit du pipeline d'observation de NGeGenBrain v3.1 & Corps & Contenu et preuves repris \\

Formalisme initial & Matrice d'audit v3.1 par rapport au formalisme & Corps & Contenu et preuves repris \\

Formalisme initial & Questions ouvertes et obligations d'implémentation & Corps & Contenu et preuves repris \\

Formalisme initial & Résumé des résultats établis & Annexe & Registre source conservé \\

Formalisme initial & Checklist de conformité d'une étude d'observation & Annexe & Registre source conservé \\

Formalisme initial & Notation & Annexe & Registre source conservé \\

Formalisme initial & Contrôles de cohérence mathématique & Annexe & Registre source conservé \\

Formalisme initial & Conclusion & Annexe & Registre source conservé \\

Next 0 - théorie inverse & Théorèmes centraux de la théorie inverse & Annexe & Registre source conservé \\

Next 0 - théorie inverse & Checklist finale de rigueur & Annexe & Registre source conservé \\

Next 0 - théorie inverse & Notation & Annexe & Registre source conservé \\

Next 0 - théorie inverse & Conclusion formelle & Annexe & Registre source conservé \\

Next 1 - invariant/variant & Registre des résultats démontrés & Annexe & Registre source conservé \\

Next 1 - invariant/variant & Checklist d'application & Annexe & Registre source conservé \\

Next 1 - invariant/variant & Conclusion & Annexe & Registre source conservé \\

\bottomrule
\end{longtable}
\section{Fichiers sources et limite de provenance}
Les sources utilisées pour cette refonte sont les trois fichiers locaux suivants : \url{NGeGenBrain_observation_formalisme.tex}, \url{NGeGenBrain_observation_formalisme_next_0.tex} et \url{NGeGenBrain_observation_formalisme_next_1.tex}. Le nom \url{NGeGenBrain_observation_formalisme_0.tex} ne correspondait pas à un fichier distinct dans l'environnement de travail ; il n'a donc pas été inventé ni traité comme une source indépendante. La suite \texttt{next\_1} est incluse afin de préserver la formalisation de l'idée invariant--variant--dynamique déjà produite.

\section{Correction de portée d'audit}
Les chapitres relatifs à l'audit de NGeGenBrain v3.1 décrivent le constat logiciel rapporté dans la source initiale et son périmètre d'observation à cette date. Ils ne prétendent pas constituer une ré-exécution d'un audit du logiciel actuel. Une affirmation logicielle doit rester liée à la version, au code et aux preuves effectivement examinés.

\chapter{Synthèse générale et chaîne logique complète}
La théorie entière forme la chaîne logique suivante :
\[
\begin{aligned}
&\text{exécution source et environnement}\\
&\quad\downarrow\;\text{protocole d'acquisition déclaré}\\
&\text{événements et enregistrements observationnels}\\
&\quad\downarrow\;\text{application de signature}\\
&\text{empreintes multi-facettes, multi-échelles et temporelles}\\
&\quad\downarrow\;\text{quotients, invariants déclarés et résidus}\\
&\text{contraintes empiriques et ensemble de mécanismes compatibles}\\
&\quad\downarrow\;\text{expériences discriminantes/interventions}\\
&\text{candidats identifiables relativement à une classe et un protocole}\\
&\quad\downarrow\;\text{synthèse, abstraction et composition exécutables}\\
&\text{mécanisme composé autonome}\\
&\quad\downarrow\;\text{tests indépendants et ré-observation}\\
&\text{validation, réfutation, révision ou incertitude résiduelle}.
\end{aligned}
\]
Chaque flèche constitue une obligation de preuve ou de mesure, non une transformation garantie automatiquement. Le passage à l'étape suivante n'autorise pas à déclarer l'étape précédente réussie sans ses critères. En particulier, une empreinte stable ne reconstitue pas ses variations si celles-ci ne sont pas déterminées par l'empreinte ; un mécanisme compatible n'est pas nécessairement identifié ; un mécanisme comportementalement fidèle n'est pas nécessairement mécanistiquement équivalent.

La contribution de la décomposition invariant--variant--dynamique est de rendre explicite l'information que la signature quotientée conserve et celle que le dossier doit encore porter. La paire $(I,V)$ reconstruit un enregistrement exactement si elle est injective. Une décomposition additive n'est unique qu'avec des conditions de séparation appropriées. En temps continu, les échantillons finis ne déterminent pas une trajectoire arbitraire ; toute garantie d'approximation dépend de la régularité, du maillage et du modèle d'erreur déclarés.

Pour NGeGenBrain, le résultat pratique n'est donc pas « calculer l'empreinte puis connaître le LLM », mais construire, maintenir et tester un ensemble traçable de contraintes qui sépare progressivement les mécanismes candidats sans masquer l'ambiguïté. La reconstruction est un programme expérimental récursif : observer, formaliser, intervenir, composer, re-observer, chercher les contre-exemples et réviser.

\backmatter
\end{document}
